\chapter{Graded unitarity and Lee--Yang}
\label{ch:L19}
\begingroup\setlength{\emergencystretch}{3em}

\providecommand{\lnineteenvac}{\mathbf{1}}
\providecommand{\lnineteenV}{\mathcal{V}}
\providecommand{\lnineteenL}{L}
\providecommand{\lnineteenG}{G}
\providecommand{\lnineteenLay}{\mathcal{L}}
\providecommand{\lnineteenh}{\langle\!\langle}
\providecommand{\lnineteenhr}{\rangle\!\rangle}
\providecommand{\lnineteencb}{c_{2,2n+1}}
\providecommand{\lnineteenPi}{\Pi}
\providecommand{\lnineteensl}{\mathfrak{sl}_2}
\providecommand{\lnineteenR}{\mathbb{R}}

This chapter documents the release of the LeanVOA library named in the front matter. Every Lean quotation below is checked by \texttt{scripts/textbook/sync.py}
to occur verbatim (up to whitespace and comments) in the source files of that release. Checks that were
run at an earlier state of the library and not re-run for this edition are marked as such. By topic, the chapter covers: (1) sesquilinear forms, real structures and the
graded-unitarity predicate; (2) the linear algebra of flags, leading minors and Schur complements; (3) the
Virasoro VOAs, their real structures and the length-two positivity theorem; (4) the Lee--Yang model;
(5) finite ``Program~B'' certificates; (6) Drinfeld--Sokolov reduction;
(7) weight-one positivity of affine $\lnineteensl$ at level $-8/5$ (Section~\ref{l19:sec:aff85}) and the abstract
transfer of graded unitarity along graded isometric embeddings (Section~\ref{l19:sec:transfer}); (8) the unconditional graded unitarity
of the universal affine VOA $V^k(\lnineteensl)$ for real $k<-2$ (Section~\ref{l19:sec:belowcritical}); (9) two
empty geometric files. An inverse Drinfeld--Sokolov programme (Adamovi\'c's free-field embedding of the boundary
admissible $L_k(\lnineteensl)$ into $L(c,0)\otimes\lnineteenPi(0)$) is under development and is not part of this
release.
By directory, the files covered are those
assigned to this chapter in
\texttt{Notes/Textbook/assignment.txt}: seven linear-algebra and structure files in
\leanpath{VOA/\allowbreak Core/} (\leanpath{GradedUnitarity}, \leanpath{FlagMinors},
\leanpath{MinorPivots}, \leanpath{HermitianSchur}, \leanpath{RealStructure},
\leanpath{RealSymmetricTwo}, \leanpath{SesqFormTransport}),
the example files in \leanpath{VOA/\allowbreak Examples/} (the \leanpath{LeeYang*} files, the
\leanpath{VirasoroLengthTwo*} files, the Virasoro real-structure and graded-unitarity files, the
\leanpath{InverseDS*} infrastructure files shipped with this release (free-field carriers and the general
Virasoro Verma module used by Chapter~\ref{ch:L06}; not otherwise documented here), the Drinfeld--Sokolov files
\leanpath{DSReductionSl2} and \leanpath{DSMorphismSl2}, a handful of ``Program~B'' finite certificates,
\leanpath{AffineSl2GradedUnitarityLow}, the twelve \leanpath{AffineSl2BelowCritical*} files, and the stub
\leanpath{MoonshineModule}), and the two files of \leanpath{VOA/\allowbreak Geometric/}.

\paragraph{What the chapter is about, in one paragraph.} A vertex operator algebra (VOA) that
arises from a four-dimensional $\mathcal N=2$ superconformal field theory carries, besides its
usual structure, a filtration (the ``$R$-filtration'') and a Hermitian form whose sign alternates
with the filtration degree. Arabi Ardehali, Beem, Lemos and Rastelli (henceforth \textbf{ABLR})
call a VOA with such data \emph{graded unitary}. The library formalises the ABLR datum
(\lean{GradedUnitaryAt}, \lean{GradedUnitary}), repairs a defect that once made it refutable,
reduces its nondegeneracy clause to leading principal minors of Gram matrices and proves the
converse Sylvester criterion needed for the positivity clause, constructs real structures on the
Virasoro VOAs, and proves the positivity clauses \emph{through filtration degree two} for every
boundary minimal model $M(2,2n+1)$, unconditionally
(\lean{boundary_}\allowbreak\lean{gradedUnitaryThroughTwo}). The full statement --- Kaufmann's 1994 conjecture that
$L(\lnineteencb,0)$ is graded unitary --- is \textbf{not} proved and is not claimed. Around this core
the chapter documents the Lee--Yang model $M(2,5)$ (its singular vector, its strong finiteness
\lean{leeYang_stronglyFinite}, its Zhu-algebra relation), the Drinfeld--Sokolov reduction of
affine $\lnineteensl$ (the arithmetic and the BRST differential; the cohomology is not built), and
the weight-one Hermitian positivity of the currents of $V^{-8/5}(\lnineteensl)$ (Section~\ref{l19:sec:aff85}).
The geometric files (geometric VOAs,
chiral de Rham complex) contain no declarations at all. Finally, Section~\ref{l19:sec:belowcritical} documents the
first theorem of the library that asserts the full predicate \lean{GradedUnitary} of a concrete VOA:
\lean{universal_}\allowbreak\lean{gradedUnitary_}\allowbreak\lean{belowCritical_}\allowbreak\lean{unconditional} --- for every real $k<-2$,
$V^k(\lnineteensl)$ with its normalised invariant form, the compact real structure and the current-length flag is
graded unitary --- proved with no hypothesis beyond $k<-2$, by a large-negative-level anchor, rigidity of layer
signs along $(-\infty,-2)$, and a ``no-wall'' lemma.

\paragraph{Prerequisites and pointers.} Generic Lean idioms (structures, \lean{Submodule},
\lean{finrank}, coercions $\uparrow$, \lean{∃}/\lean{∀}) are explained in Chapter~\ref{ch:01}.
The universal and simple Virasoro VOAs $V(c,0)$, $L(c,0)$, the canonical invariant form
\lean{canB}/\lean{Bsimp}, creation words \lean{negWordVir}, and the Kac determinant are in
Chapter~\ref{ch:L06}; the Lee--Yang rigidity theorem is in Chapter~\ref{ch:L07}; affine $\lnineteensl$
at level $-8/5$ is in Chapters~\ref{ch:L08} and~\ref{ch:L09}; graded flags and good filtrations are
in Chapter~\ref{ch:L11} (Definition~\ref{l11:def:flag}); invariant forms are in
Chapter~\ref{ch:L04} (Definition~\ref{l04:def:invform}); $C_2$-cofiniteness and
\lean{StronglyFinite} are in Chapter~\ref{ch:L12} (Definitions~\ref{l12:def:c2}
and~\ref{l12:def:strong}); recall that $V$ is $C_2$-cofinite when $V/C_2(V)$ is finite-dimensional, where
$C_2(V)$ is the span of all $a_{(-2)}b$ with $a,b\in V$. Whenever a statement uses one of these notions we recall it in a sentence
and point to the definition.

\paragraph{Trust model.} Every headline of this chapter was checked with \texttt{\#print axioms}
(Section~\ref{l19:sec:axioms}); all of them depend only on Lean's three standard axioms
\lean{propext}, \lean{Classical.choice}, \lean{Quot.sound}. None uses \lean{VOA.pbw} (the PBW theorem of Chapter~\ref{ch:L02}; formerly a project axiom, now a theorem). ``Axiom-clean'' does not mean ``unconditional'': many theorems below carry
explicit hypotheses, and we name every one.

%======================================================================
\section{Unitarity: the standard notions and the graded one}\label{l19:sec:notions}
%======================================================================

\subsection{Why physicists care about positivity}

In quantum mechanics the space of states is a Hilbert space: the inner product
$\langle\psi|\psi\rangle$ is a probability amplitude squared and must be positive. In a
two-dimensional conformal field theory (CFT) on the plane, radial quantisation identifies states
with local fields (the state--field correspondence $a\mapsto Y(a,z)$ of a VOA), and the inner
product of the Hilbert space becomes a sesquilinear form on the space of states $V$ for which the
modes of the stress tensor satisfy $L_n^\dagger = L_{-n}$. Equivalently, in Euclidean signature,
positivity is \emph{reflection positivity} of correlation functions (Osterwalder--Schrader). A CFT
whose state space has a positive-definite such form is called \emph{unitary}.

Unitarity has dramatic consequences. For the Virasoro algebra alone, positivity of the vacuum
module requires $c\ge0$ (already at level two, $\|L_{-2}\lnineteenvac\|^2 = c/2$), and the work of
Friedan, Qiu and Shenker (1984) and Goddard, Kent and Olive (1985--86) shows that the irreducible
highest-weight representation $L(c,h)$ is unitary only for $c\ge1$, $h\ge0$, or for the
\emph{discrete series}
\begin{equation*}
 c = 1-\frac{6}{m(m+1)},\qquad m = 2,3,4,\dots,
\end{equation*}
which are the minimal models $M(m,m+1)$ ($m=3$: Ising, $c=1/2$; $m=4$: tricritical Ising,
$c=7/10$). Every other minimal model $M(p,q)$, $|p-q|\ge2$, is non-unitary. The simplest such is the
\textbf{Lee--Yang model} $M(2,5)$ with
\begin{equation*}
 c_{2,5} = 1 - \frac{6(2-5)^2}{2\cdot5} = -\frac{22}{5},
\end{equation*}
whose vacuum module already has $\|L_{-2}\lnineteenvac\|^2 = c/2 = -11/5 < 0$.

\subsection{Unitary VOAs in the sense of Dong and Lin}

Dong and Lin (2014) formalised the physicists' notion as follows. Let $V$ be a VOA of CFT type with
conformal vector $\omega$.

\begin{definition}[Unitary VOA; Dong--Lin 2014]\label{l19:def:donglin}
A \emph{unitary VOA} is a pair $(V,\phi)$ where $\phi:V\to V$ is an \emph{anti-linear} involutive
VOA automorphism (so $\phi(za) = \bar z\phi(a)$, $\phi(a_{(n)}b) = \phi(a)_{(n)}\phi(b)$,
$\phi(\lnineteenvac)=\lnineteenvac$, $\phi(\omega)=\omega$, $\phi^2=1$), together with a
\emph{positive-definite} Hermitian form $(\cdot,\cdot)$ on $V$, linear in the first slot and
anti-linear in the second, normalised by $(\lnineteenvac,\lnineteenvac) = 1$, such that for all $a,u,v\in V$
\begin{equation*}
 \big(Y(e^{zL_1}(-z^{-2})^{L_0}a,z^{-1})u,\;v\big) = \big(u,\;Y(\phi(a),z)v\big).
\end{equation*}
\end{definition}

\begin{physics}
$\phi$ is the CPT-type anti-unitary symmetry that maps a field to its adjoint, and the displayed
identity says $Y(a,z)^\dagger = Y(e^{\bar zL_1}(-\bar z^{-2})^{L_0}\phi(a),\bar z^{-1})$: the
adjoint of a local field is again a local field, inserted at the image of $z$ under inversion
$z\mapsto 1/\bar z$. For $a=\omega$ it gives $L_n^\dagger = L_{-n}$. Positivity of $(\cdot,\cdot)$ is
the Hilbert-space positivity of the CFT.
\end{physics}

Ai and Lin (2017) extended this to vertex operator superalgebras (with a parity twist on the odd
part), and Carpi, Kawahigashi, Longo and Weiner (CKLW, 2018) proved that, under an additional
``strong locality'' hypothesis, unitary VOAs are equivalent to conformal nets of von Neumann
algebras on the circle. For Virasoro VOAs the classification is the one quoted above:
$L(c,0)$ is a unitary VOA iff $c\ge1$ or $c = 1-6/(m(m+1))$.

\begin{warning}[title={The library has no Dong--Lin unitarity}]\label{l19:warn:nodonglin}
The library contains a structure \lean{VOA.Unitary} in \leanpath{VOA/\allowbreak Modules/\allowbreak Contragredient.lean},
but Chapter~\ref{ch:L12} (Proposition~\ref{l12:prop:unitary}) shows it is \textbf{unsatisfiable} on every
module of dimension $\ge2$ (it asks a $\C$-bilinear form to be ``positive''; see the next
subsection) and that its invariance clause is not the Dong--Lin one. It has no consumers. The
positivity notion actually used in the library is the \emph{graded} one of this chapter, which is a
different notion (Remark~\ref{l19:rem:incompatible} explains that, for Virasoro VOAs with
$\omega$ in filtration degree one, the two are mutually exclusive).
\end{warning}

\subsection{Bilinear versus sesquilinear forms}

A recurring point of this chapter is elementary but crucial. A VOA carries a canonical
\emph{$\C$-bilinear} invariant form $(\cdot|\cdot)$ (Chapter~\ref{ch:L04}): it is
$\C$-linear in \emph{both} slots, and for the Virasoro VOA its Gram matrices are the Kac matrices of
Chapter~\ref{ch:L06}. A bilinear form cannot be positive on a complex vector space: if
$(x|x) = r>0$ then $(ix|ix) = i^2r = -r<0$. Positivity is a property of a \emph{sesquilinear}
form $h$, anti-linear in one slot, for which $h(ix,ix) = \bar i\,i\,h(x,x) = h(x,x)$. To pass from
the bilinear form to a sesquilinear one, one needs an anti-linear map $\rho$ and sets
$h(a,b) = (\rho a|b)$. This is exactly what the Dong--Lin $\phi$ does, and exactly what the
library's \lean{RealStructure} (Section~\ref{l19:sec:real}) does.

\subsection{Graded unitarity: where it comes from}\label{l19:sec:ablrorigin}

Beem, Lemos, Liendo, Peelaers, Rastelli and van Rees (2015) discovered that every four-dimensional
$\mathcal N=2$ superconformal field theory (SCFT) contains a protected subsector whose operator
algebra is a VOA $V$, with $c_{2d} = -12\,c_{4d}$. Since $c_{4d}>0$ in a unitary four-dimensional
theory, $c_{2d}<0$: these VOAs are \emph{never} unitary in the sense of Dong--Lin. Nevertheless they
inherit a shadow of four-dimensional unitarity. Each state of $V$ comes from a four-dimensional
operator with an $SU(2)_R$ charge $R$; the charge is not preserved by the VOA operations, but it
defines an increasing filtration $\lnineteenG_0V\subseteq \lnineteenG_1V\subseteq\cdots$ (the
\emph{$R$-filtration}). The four-dimensional Hilbert-space inner product, transported to $V$, is a
Hermitian form which on the associated graded pieces has sign $(-1)^R$. For the Virasoro VOA the
conjectured $R$-filtration is the filtration by the number of Virasoro creation modes:
$L_{-n_1}\cdots L_{-n_k}\lnineteenvac$ has filtration degree $k$.

\begin{definition}[Graded unitarity, informal; ABLR, arXiv:2507.23781, Definition~2.2]\label{l19:def:ablrinformal}
A VOA $V$ with an invariant bilinear form $(\cdot|\cdot)$ is \emph{graded unitary} if it carries
\begin{enumerate}
\item[(i)] a grading compatible with an ``$R$-parity'' $s$;
\item[(ii)] an anti-linear VOA automorphism $\rho$ with $\rho^2 = s$;
\item[(iii)] the sesquilinear form $(a,b)\mapsto(\rho a|b)$;
\item[(iv)] an increasing exhaustive filtration $\lnineteenG_pV_n$ of each weight space $V_n$ (a
  ``good'' filtration) on each step of which the form is nondegenerate;
\item[(v)] the orthogonal decomposition $\lnineteenG_pV_n = \lnineteenG_{p-1}V_n\oplus
  \lnineteenLay_p$, $\lnineteenLay_p = \lnineteenG_pV_n\cap(\lnineteenG_{p-1}V_n)^\perp$;
\item[(vi)] positivity of the twisted Hermitian form $(-1)^p(\rho a|b)$ on every layer
  $\lnineteenLay_p$.
\end{enumerate}
\end{definition}

\begin{physics}
Clause (vi) says: a state of $R$-charge $p$ (more precisely, the component of a state that is
orthogonal to everything of lower charge) has four-dimensional norm of sign $(-1)^p$ when measured
with the two-dimensional form. Graded unitarity is therefore a necessary condition for a VOA to be
the chiral algebra of a unitary 4d $\mathcal N=2$ SCFT, and a sharp tool to rule out candidate VOAs.
The boundary minimal models $M(2,2n+1)$ are the chiral algebras of the Argyres--Douglas theories
$(A_1,A_{2n})$ (Lee--Yang: $(A_1,A_2)$), so the conjecture that they are graded unitary --- first
formulated, in the language of signatures, by Kaufmann (1994) --- is a consistency check of the
SCFT/VOA dictionary.
\end{physics}

\begin{remark}[Graded unitarity and Dong--Lin unitarity are incompatible for Virasoro VOAs]\label{l19:rem:incompatible}
For $L(c,0)$ with the length filtration, $\omega = L_{-2}\lnineteenvac$ spans the degree-one layer in
weight two, so clause (vi) requires $-(\rho\omega|\omega) = -c/2 >0$, i.e.\ $c<0$; Dong--Lin
unitarity requires $(\omega,\omega) = c/2>0$. Thus no Virasoro VOA with $c\ne0$ is both. Graded
unitarity is not a weakening of unitarity: it is a different positivity, adapted to $c<0$.
\end{remark}

The library's encoding of clauses (i)--(vi) is the subject of
Sections~\ref{l19:sec:sesq}--\ref{l19:sec:reduction}.

%======================================================================
\section{Sesquilinear forms and positivity on a subspace}\label{l19:sec:sesq}
%======================================================================

\subsection{The defect that once made graded unitarity refutable}

\begin{theorem}[A bilinear form is never positive on a nonzero complex subspace]\label{l19:thm:bilpos}
Let $\beta$ be a $\C$-bilinear form on a complex vector space $E$ and $S\le E$ a complex subspace.
If $\beta(x,x)$ is a positive real number for every nonzero $x\in S$, then $S = 0$.
\end{theorem}

\noindent File \leanpath{VOA/\allowbreak Core/\allowbreak GradedUnitarity.lean}.
\begin{leancode}
theorem eq_bot_of_bilinear_pos (β : BilinForm ℂ E) (S : Submodule ℂ E)
    (hpos : ∀ x ∈ S, x ≠ 0 → ∃ r : ℝ, 0 < r ∧ β x x = (r : ℂ)) : S = ⊥ := by
\end{leancode}

\begin{dictionary}
\lean{BilinForm ℂ E} is Mathlib's type of $\C$-bilinear forms $E\times E\to\C$ (a bilinear map,
curried: \lean{β x y}). \lean{Submodule ℂ E} is a complex subspace; \lean{S = ⊥} says it is the zero
subspace. ``$\beta(x,x)$ is a positive real'' is spelled \lean{∃ r : ℝ, 0 < r ∧ β x x = (r : ℂ)},
where \lean{(r : ℂ)} is the real number $r$ viewed as a complex number. The statement is exactly
Theorem~\ref{l19:thm:bilpos}. It is not a VOA statement: it is kept in the library as a
\emph{recorded defect}, see the warning below.
\end{dictionary}

\paragraph{Proof on paper.} If $x\in S$ is nonzero with $\beta(x,x) = r>0$, then $ix\in S$ is
nonzero and $\beta(ix,ix) = i^2\beta(x,x) = -r<0$, which is not a positive real. \qed

\begin{physics}
A ``probability'' computed with a bilinear form changes sign when the state is multiplied by the
phase $i$; a physical norm must be phase-invariant. This is why every positivity statement in
quantum theory is about a Hermitian form.
\end{physics}

\begin{status}
Proved; \texttt{\#print axioms}: \lean{[propext, Classical.choice, Quot.sound]}. The same theorem is
discussed in Chapter~\ref{ch:15} (Theorem~\ref{ch15:thm:bilpos}).
\end{status}

\begin{warning}[title={Recorded defect, now repaired: the old clause (vi) was false for every VOA}]\label{l19:warn:olddefect}
An earlier version of \leanpath{VOA/\allowbreak Core/\allowbreak GradedUnitarity.lean} stated clause (vi) as positivity
of the $\C$-\emph{bilinear} invariant form on the layers $\lnineteenLay_p$, via a predicate
\lean{IsPosOn (β : BilinForm ℂ E) (S : Submodule ℂ E)} of exactly the shape of the hypothesis of
Theorem~\ref{l19:thm:bilpos}. By that theorem every layer would have to be zero, contradicting
exhaustiveness on any nonzero weight space; so the old \lean{GradedUnitary} and
\lean{VirGradedUnitary} were \textbf{provably false} for every VOA with $V_0\ne0$, while their
docstrings called the conjecture ``open''. The repair, described in the file's §0 and
checked below, is to state positivity for a \emph{sesquilinear} form built from an anti-linear
$\rho$ (ABLR clause (iii)). Theorem~\ref{l19:thm:bilpos} is kept as a theorem precisely so that the
mistake cannot be reintroduced. (``ABLR'' is Arabi Ardehali--Beem--Lemos--Rastelli, as the source files cite.)
\end{warning}

\subsection{Sesquilinear forms}

\begin{definition}[Sesquilinear form]\label{l19:def:sesq}
A \emph{sesquilinear form} on a complex vector space $E$ is a map $h:E\times E\to\C$ that is
additive in each slot, conjugate-homogeneous in the first slot and homogeneous in the second:
$h(za,b) = \bar z\,h(a,b)$, $h(a,zb) = z\,h(a,b)$.
\end{definition}

\noindent File \leanpath{VOA/\allowbreak Core/\allowbreak GradedUnitarity.lean}.
\begin{leancode}
structure SesqForm (E : Type*) [AddCommGroup E] [Module ℂ E] where
  /-- The underlying function. -/
  toFun : E → E → ℂ
  /-- Additivity in the first slot. -/
  add_left : ∀ a a' b : E, toFun (a + a') b = toFun a b + toFun a' b
  /-- **Conjugate**-homogeneity in the first slot: `⟨z a, b⟩ = z̄ ⟨a,b⟩`. -/
  smul_left : ∀ (z : ℂ) (a b : E), toFun (z • a) b = (starRingEnd ℂ) z * toFun a b
  /-- Additivity in the second slot. -/
  add_right : ∀ a b b' : E, toFun a (b + b') = toFun a b + toFun a b'
  /-- ℂ-homogeneity in the second slot: `⟨a, z b⟩ = z ⟨a,b⟩`. -/
  smul_right : ∀ (z : ℂ) (a b : E), toFun a (z • b) = z * toFun a b
\end{leancode}

\begin{dictionary}
A Lean \lean{structure} bundles a function \lean{toFun} with proofs of its four laws.
\lean{starRingEnd ℂ} is complex conjugation $z\mapsto\bar z$ (Mathlib's name for the ``star''
ring endomorphism). \lean{z • a} is scalar multiplication. A coercion (declared in the file) lets
one write \lean{h a b} for \lean{h.toFun a b}. Convention: anti-linear in the \emph{first} slot,
as in ABLR's (2.30) and in physics ($\langle a|b\rangle$), the opposite of Dong--Lin's
convention; this is a harmless relabelling.

\textbf{What is not required.} Hermitian symmetry $h(b,a) = \overline{h(a,b)}$ is \emph{not} a field.
The docstring says this is deliberate: positivity is stated as ``$h(x,x)$ equals a positive
real'', which already asserts reality of the diagonal where it is needed. Since a sesquilinear form
whose diagonal is real on a subspace is automatically Hermitian there (polarisation), nothing is lost
for positivity statements; but the reader should know that \lean{SesqForm} alone is the weaker
notion.
\end{dictionary}

\begin{physics}
$h(a,b)$ is the transition amplitude $\langle a|b\rangle$ between two states. Anti-linearity in the
bra is what makes $\langle a|a\rangle$ insensitive to the phase of $a$.
\end{physics}

\begin{status}
A definition. The coercion and the lemmas \lean{SesqForm.zero_left}, \lean{zero_right},
\lean{smul_self} are proved in the same file.
\end{status}

\begin{theorem}[The obstruction disappears]\label{l19:thm:ismul}
For every sesquilinear form $h$ and vector $x$, $h(ix,ix) = h(x,x)$.
\end{theorem}

\noindent File \leanpath{VOA/\allowbreak Core/\allowbreak GradedUnitarity.lean}.
\begin{leancode}
theorem I_smul_self (h : SesqForm E) (x : E) :
    h (Complex.I • x) (Complex.I • x) = h x x := by
\end{leancode}

\begin{dictionary}
\lean{Complex.I} is $i$. The theorem is \lean{SesqForm.I_smul_self}. Exactly the statement.
\end{dictionary}

\paragraph{Proof on paper.} $h(ix,ix) = \bar i\,i\,h(x,x) = |i|^2h(x,x) = h(x,x)$. \qed

\begin{physics}
Multiplying a state by a phase does not change its norm.
\end{physics}

\begin{status}
Proved, axiom-clean (Section~\ref{l19:sec:axioms}).
\end{status}

\begin{definition}[Positivity on a subspace; the sign twist]\label{l19:def:isposon}
A sesquilinear form $h$ is \emph{positive on} a subspace $S$ if $h(x,x)$ is a positive real number
for every nonzero $x\in S$. For $p\in\N$, the \emph{signed form} of degree $p$ is $(-1)^ph$.
\end{definition}

\noindent File \leanpath{VOA/\allowbreak Core/\allowbreak GradedUnitarity.lean}.
\begin{leancode}
def IsPosOn (h : SesqForm E) (S : Submodule ℂ E) : Prop :=
  ∀ x ∈ S, x ≠ 0 → ∃ r : ℝ, 0 < r ∧ h x x = (r : ℂ)

def signedForm (h : SesqForm E) (p : ℕ) : SesqForm E where
  toFun a b := (-1 : ℂ) ^ p * h a b
  add_left a a' b := by rw [h.add_left]; ring
  smul_left z a b := by rw [h.smul_left]; ring
  add_right a b b' := by rw [h.add_right]; ring
  smul_right z a b := by rw [h.smul_right]; ring
\end{leancode}

\begin{dictionary}
\lean{IsPosOn h S} is literally Definition~\ref{l19:def:isposon}; note it now takes a
\lean{SesqForm}, not a \lean{BilinForm}, which is the repair of
Warning~\ref{l19:warn:olddefect}. \lean{signedForm h p} is the form $(a,b)\mapsto(-1)^ph(a,b)$; the
four \lean{by ...} lines are the proofs that it is again sesquilinear. The exponent is $p$, the
filtration degree, matching ABLR's $\mathrm{sign}\langle O|O\rangle = (-1)^{R_O}$ (their eq.~(3.16),
Kaufmann's (4.4)) under the library's normalisation $(\lnineteenvac|\lnineteenvac) = +1$.

\textbf{Could this be vacuous or refutable?} The file contains three witnesses, all proved:
\lean{isPosOn_stdHerm_top} (the standard form $\bar ab$ on $\C$ is positive on all of $\C$, a
nonzero subspace --- so \lean{IsPosOn} is satisfiable where the old predicate was not);
\lean{isPosOn_}\allowbreak\lean{signedForm_}\allowbreak\lean{stdHerm_}\allowbreak\lean{zero} (the twist at $p=0$ keeps it positive); and
\lean{not_}\allowbreak\lean{isPosOn_}\allowbreak\lean{signedForm_}\allowbreak\lean{stdHerm_}\allowbreak\lean{one} (at $p=1$ the same form on the same subspace \emph{fails}),
so the predicate has teeth and the two parities impose opposite conditions. On the zero subspace it
holds vacuously (\lean{isPosOn_bot}); this is why later we check that the layers it is applied to are
nonzero.
\end{dictionary}

\begin{physics}
\lean{IsPosOn} is ``all states in $S$ have positive norm''; \lean{signedForm} implements the
$(-1)^R$ of the SCFT/VOA dictionary: in odd $R$-charge the two-dimensional form is the
\emph{negative} of the four-dimensional norm.
\end{physics}

\begin{status}
Definitions; the witnesses \lean{isPosOn_stdHerm_top}, \lean{isPosOn_}\allowbreak\lean{signedForm_}\allowbreak\lean{stdHerm_}\allowbreak\lean{zero},
\lean{not_}\allowbreak\lean{isPosOn_}\allowbreak\lean{signedForm_}\allowbreak\lean{stdHerm_}\allowbreak\lean{one}, \lean{isPosOn_I_smul} are proved and axiom-clean.
\end{status}

\begin{theorem}[Positivity pulls back along injective maps]\label{l19:thm:pullback}
Let $\iota:E\to F$ be an injective $\C$-linear map, $h_F$ a sesquilinear form on $F$, and
$h_E = \iota^*h_F$, i.e.\ $h_E(x,y) = h_F(\iota x,\iota y)$. If $(-1)^ph_F$ is positive on
$\iota(S)$, then $(-1)^ph_E$ is positive on $S$.
\end{theorem}

\noindent File \leanpath{VOA/\allowbreak Core/\allowbreak SesqFormTransport.lean}.
\begin{leancode}
def pullback (h : SesqForm F) (iota : E →ₗ[ℂ] F) : SesqForm E where
  toFun x y := h (iota x) (iota y)
  add_left x x' y := by rw [map_add, h.add_left]
  smul_left a x y := by rw [map_smul, h.smul_left]
  add_right x y y' := by rw [map_add, h.add_right]
  smul_right a x y := by rw [map_smul, h.smul_right]

theorem isPosOn_signedForm_of_eq_pullback {hE : SesqForm E} {hF : SesqForm F}
    {iota : E →ₗ[ℂ] F} (hiota : Function.Injective iota)
    (hcompat : hE = hF.pullback iota) {S : Submodule ℂ E} (p : ℕ)
    (hpos : IsPosOn (signedForm hF p) (S.map iota)) :
    IsPosOn (signedForm hE p) S :=
\end{leancode}

\begin{dictionary}
\lean{E →ₗ[ℂ] F} is the type of $\C$-linear maps. \lean{S.map iota} is the image subspace
$\iota(S)$. \lean{hcompat : hE = hF.pullback iota} says $h_E$ \emph{is} the pullback (an exact
isometry). The theorem is \lean{SesqForm.}\allowbreak\lean{isPosOn_}\allowbreak\lean{signedForm_}\allowbreak\lean{of_}\allowbreak\lean{eq_}\allowbreak\lean{pullback}; it is exactly
Theorem~\ref{l19:thm:pullback}. Injectivity is needed: a nonzero $x$ with $\iota x = 0$ would have
$h_E(x,x) = 0$.
\end{dictionary}

\paragraph{Proof on paper.} For $0\ne x\in S$, $\iota x\ne0$ lies in $\iota(S)$, so
$(-1)^ph_E(x,x) = (-1)^ph_F(\iota x,\iota x)$ is a positive real. \qed

\begin{physics}
If a theory is realised inside a bigger one by an isometric embedding of state spaces (a free-field
realisation, say), then positivity of norms in the bigger theory, restricted to the image,
gives positivity in the smaller one.
This is the engine of the transfer theorem of Section~\ref{l19:sec:transfer}.
\end{physics}

\begin{status}
Proved, axiom-clean. Also proved: \lean{pullback_comp}, \lean{pullback_signedForm},
\lean{isPosOn_pullback}, \lean{isPosOn_of_eq_pullback}.
\end{status}

%======================================================================
\section{Real structures}\label{l19:sec:real}
%======================================================================

\begin{definition}[Anti-linear automorphism with parity; ABLR clauses (i)--(ii)]\label{l19:def:real}
A \emph{real structure} on a VOA $W$ (in the library's sense) consists of an additive map
$\rho:V\to V$ such that $\rho(zv) = \bar z\rho(v)$, $\rho(u_{(n)}v) = \rho(u)_{(n)}\rho(v)$ for all
$n\in\Z$, $\rho(\lnineteenvac) = \lnineteenvac$, $\rho(\omega) = \omega$, together with a function
$s:\Z\to\C$ (the ``$R$-parity'') such that $\rho^2 = s(h)$ on the weight-$h$ space.
\end{definition}

\noindent File \leanpath{VOA/\allowbreak Core/\allowbreak RealStructure.lean}.
\begin{leancode}
structure RealStructure (W : VertexOperatorAlgebra 𝒱) where
  /-- The underlying additive map. -/
  rho : 𝒱 →+ 𝒱
  /-- **Anti-linearity**: `ρ(z·v) = z̄·ρ(v)`. -/
  antilinear : ∀ (z : ℂ) (v : 𝒱), rho (z • v) = (starRingEnd ℂ) z • rho v
  /-- **Vertex-algebra compatibility**: `ρ(u₍ₙ₎v) = ρ(u)₍ₙ₎ρ(v)` for every mode. -/
  map_smul : ∀ (u : 𝒱) (n : ℤ) (v : 𝒱), rho (W.smul u n v) = W.smul (rho u) n (rho v)
  /-- `ρ` fixes the vacuum. -/
  map_vacuum : rho W.vacuum = W.vacuum
  /-- `ρ` fixes the conformal vector. -/
  map_omega : rho W.conformalVector = W.conformalVector
  /-- ABLR's `R`-parity `s`, as a function of the conformal weight.  Carried as data. -/
  parity : ℤ → ℂ
  /-- **`ρ² = s`** on each weight space. -/
  sq : ∀ {h : ℤ} (v : 𝒱), v ∈ W.grading h → rho (rho v) = parity h • v
\end{leancode}

\begin{dictionary}
\lean{𝒱 →+ 𝒱} is an additive group homomorphism (not a linear map; anti-linearity is a separate
field, a design choice explained in the file docstring). \lean{W.smul u n v} is the $n$-th product
$u_{(n)}v$ (Chapter~\ref{ch:01}). \lean{W.grading h} is the weight-$h$ space $V_h$. The structure is
exactly Definition~\ref{l19:def:real}.

\textbf{Comparison with Dong--Lin.} With $s\equiv1$ a \lean{RealStructure} is precisely Dong--Lin's
anti-linear involution $\phi$ (Definition~\ref{l19:def:donglin}). With general $s$ it is ABLR's
clause (ii) $\rho^2 = s$. The file's docstring records that ABLR's $s = (-1)^{h+d/2}$ was not checked
against the source (the paper was not in the repository) and that therefore $s$ is a free field; no
theorem of the library depends on its value except through $s(n)\ne0$
(Theorem~\ref{l19:thm:orthsame}). The Virasoro instances below all have $s\equiv1$.

\textbf{Not required:} that $\rho$ preserve weight spaces. It follows from $\rho(\omega)=\omega$
and \lean{map_smul} (so $\rho L_0 = L_0\rho$, \lean{RealStructure.rho_L}) together with
anti-linearity and the reality of the $L_0$-eigenvalues.
\end{dictionary}

\begin{physics}
$\rho$ is the anti-unitary ``conjugation'' that sends a field to its Hermitian conjugate up to the
$R$-parity sign; for the Virasoro VOA it fixes every state $L_{-n_1}\cdots L_{-n_k}\lnineteenvac$
because the stress tensor is a Hermitian field with real OPE coefficients.
\end{physics}

\begin{status}
A definition. \lean{RealStructure.rho_L} ($\rho L_n = L_n\rho$) and
\lean{RealStructure.}\allowbreak\lean{parity_}\allowbreak\lean{eq_}\allowbreak\lean{one_}\allowbreak\lean{of_}\allowbreak\lean{fixes} are proved in the same file. Instances: Virasoro, below.
\end{status}

\begin{definition}[The sesquilinear form of clause (iii)]\label{l19:def:gradedsesq}
Given an invariant bilinear form $B$ and a real structure $\rho$ on $W$, set
$h(a,b) = B(\rho a,b)$, and let $h_n$ be its restriction to $V_n\times V_n$.
\end{definition}

\noindent File \leanpath{VOA/\allowbreak Core/\allowbreak GradedUnitarity.lean}.
\begin{leancode}
def sesqForm (B : InvariantForm W) (R : RealStructure W) (a b : 𝒱) : ℂ := B.form (R.rho a) b

def gradedSesq (B : InvariantForm W) (R : RealStructure W) (n : ℤ) : SesqForm (W.grading n) where
\end{leancode}

\begin{dictionary}
\lean{InvariantForm W} is the library's invariant bilinear form (Chapter~\ref{ch:L04},
Definition~\ref{l04:def:invform}); \lean{B.form} is the underlying bilinear map. \lean{gradedSesq B R n}
bundles $(a,b)\mapsto B(\rho a,b)$ on the weight space \lean{W.grading n} as a \lean{SesqForm}; its
four laws are proved from \lean{R.antilinear} (we omit the four proof lines; the simp lemma
\lean{gradedSesq_apply} states \lean{gradedSesq B R n a b = B.form (R.rho (a : 𝒱)) (b : 𝒱)}).
On a set fixed by $\rho$ the two forms agree: \lean{sesqForm_eq_form_of_fixes} and
\lean{sesq_diag_eq_gradedForm}. This is why the library's \emph{bilinear} Gram matrices (whose
entries on Virasoro creation monomials are real) compute the \emph{Hermitian} values of clause (vi).
\end{dictionary}

\begin{physics}
$h(a,b) = (\rho a|b)$ is the BPZ two-point function with the bra conjugated: the Hilbert-space
inner product of radial quantisation.
\end{physics}

\begin{status}
Definitions; homogeneity lemmas \lean{sesqForm_smul_left}, \lean{sesqForm_smul_right} proved.
\end{status}

%======================================================================
\section{The graded-unitarity predicate}\label{l19:sec:gu}
%======================================================================

\subsection{Flags, layers, and the two clauses}

Recall from Chapter~\ref{ch:L11} (Definition~\ref{l11:def:flag}) that a \lean{GradedFlag W} is a family
of subspaces $\lnineteenG_pV_n$ ($p\in\N$, $n\in\Z$), each inside $V_n$, increasing in $p$, exhausting
$V_n$, with $\lnineteenG_0V_n = 0$ for $n>0$ and $\lnineteenG_0V_0 = \C\lnineteenvac$. In Lean,
\lean{F.piece p n} is $\lnineteenG_pV_n$ as a subspace of $V$, \lean{F.pieceIn p n} is the same subspace
viewed inside the finite-dimensional space $V_n$, and \lean{F.pieceLT p n} is
$\lnineteenG_{p-1}V_n$ (with $\lnineteenG_{-1} = 0$). The flag is \emph{not} required to be compatible
with the vertex operators (Chapter~\ref{ch:L11} records that the ``good'' clause
\lean{ModeCompatible} is a reconstruction and is never used here).

\begin{definition}[Layers; clauses (iv) and (vi); graded unitarity]\label{l19:def:gu}
Let $F$ be a graded flag, $\beta$ a bilinear form and $h$ a sesquilinear form on $V_n$.
\begin{enumerate}
\item The \emph{layer} $\lnineteenLay_p = \lnineteenG_pV_n\cap(\lnineteenG_{p-1}V_n)^{\perp_\beta}$.
\item $(F,\beta)$ is \emph{layer-nondegenerate} at $n$ if $\beta$ restricted to every
  $\lnineteenG_pV_n$ is nondegenerate.
\item $(F,\beta,h)$ is \emph{layer-positive} at $n$ if $(-1)^ph$ is positive on $\lnineteenLay_p$ for
  every $p$.
\item \emph{Graded unitarity at $n$} is the conjunction of 2 and 3. A VOA with invariant form $B$, real
  structure $\rho$ and flag $F$ is \emph{graded unitary} if it is graded unitary at every $n$ with
  $\beta = B|_{V_n}$ and $h = h_n$ of Definition~\ref{l19:def:gradedsesq}.
\end{enumerate}
\end{definition}

\noindent File \leanpath{VOA/\allowbreak Core/\allowbreak GradedUnitarity.lean}.
\begin{leancode}
def layer (F : GradedFlag W) (n : ℤ) (β : BilinForm ℂ (W.grading n)) (p : ℕ) :
    Submodule ℂ (W.grading n) :=
  F.pieceIn p n ⊓ β.orthogonal (F.pieceLT p n)

def LayerNondegenerate (F : GradedFlag W) (n : ℤ) (β : BilinForm ℂ (W.grading n)) : Prop :=
  ∀ p : ℕ, (β.restrict (F.pieceIn p n)).Nondegenerate

def LayerPositive (F : GradedFlag W) (n : ℤ) (β : BilinForm ℂ (W.grading n))
    (h : SesqForm (W.grading n)) : Prop :=
  ∀ p : ℕ, IsPosOn (signedForm h p) (F.layer n β p)

def GradedUnitaryAt (F : GradedFlag W) (n : ℤ) (β : BilinForm ℂ (W.grading n))
    (h : SesqForm (W.grading n)) : Prop :=
  LayerNondegenerate F n β ∧ LayerPositive F n β h

def GradedUnitary (B : InvariantForm W) (R : RealStructure W) (F : GradedFlag W) : Prop :=
  ∀ n : ℤ, GradedUnitaryAt F n (LiRadical.gradedForm B n) (gradedSesq B R n)
\end{leancode}

\begin{dictionary}
\textbf{Decoding.} \lean{⊓} is intersection of subspaces. \lean{β.orthogonal S} is Mathlib's
orthogonal complement $\{y : \beta(x,y) = 0\ \forall x\in S\}$ (for a reflexive form the side does not
matter). \lean{β.restrict S} is $\beta$ on $S\times S$, and \lean{.Nondegenerate} means the only
$x\in S$ with $\beta(x,S) = 0$ is $0$. \lean{LiRadical.gradedForm B n} is $B$ restricted to
$V_n\times V_n$ (the level-$n$ Gram form; \leanpath{VOA/\allowbreak Core/\allowbreak InvariantFormRadical.lean}).

\textbf{Clause by clause against ABLR.}
(i)--(ii): the \lean{RealStructure} argument \lean{R}. (iii): \lean{gradedSesq B R n}. (iv):
\lean{LayerNondegenerate}. (v): \emph{not an assumption}; it is derived from (iv) in
Theorem~\ref{l19:thm:layerdec}. (vi): \lean{LayerPositive}, for the \emph{sesquilinear} form, with
the twist $(-1)^p$.

\textbf{Two forms in one definition.} The layer is cut out by the \emph{bilinear}
$\beta$, positivity is for the \emph{sesquilinear} $h$. This is legitimate when the steps are
$\rho$-stable: Theorem~\ref{l19:thm:orthsame} shows the two orthogonal complements coincide. For
the Virasoro flag every step is $\rho$-stable for every real structure
(Theorem~\ref{l19:thm:rhostable}).

\textbf{What the definition does \emph{not} pin down.}
\begin{itemize}
\item \emph{The flag is a free parameter.} \lean{GradedUnitary B R F} is a statement about a given
  \lean{F}. ABLR's flag is a specific one (the $R$-filtration); the library's \lean{GradedFlag} does
  not even require compatibility with the vertex operators. See Warning~\ref{l19:warn:freeflag}.
\item \emph{$\rho$-stability of the steps} is not required (it holds in the Virasoro instance).
\item \emph{The $R$-parity} $s$ plays no role in the predicate.
\end{itemize}
\end{dictionary}

\begin{physics}
$\lnineteenLay_p$ is the space of states of ``exact $R$-charge $p$'': the part of
$\lnineteenG_pV_n$ orthogonal to all states of lower charge. Clause (iv) says no state of charge $\le p$
is null with respect to all of charge $\le p$ (no accidental null states in a truncated sector);
clause (vi) says the exact-charge-$p$ states have four-dimensional norm sign $(-1)^p$.
\end{physics}

\begin{status}
Definitions. The docstring stresses: ``This is a predicate. No theorem in this repository asserts
it of anything.'' That was accurate at the earlier state of the library against which this chapter was first written:
\lean{GradedUnitary} was then never proved for any VOA. At that snapshot, what was proved were reductions (below), the vacuum slice, and the truncated
version through degree two (Section~\ref{l19:sec:lengthtwo}). \emph{Later addition:} in a later revision of the library
the predicate is asserted, and proved unconditionally, for the universal affine VOA $V^k(\lnineteensl)$ at every real
$k<-2$ with a fixed (not free) flag (Theorem~\ref{l19:thm:bcheadline}); the docstring sentence quoted above, still
present at a pre-release build commit, is therefore stale.
\end{status}

\begin{warning}[title={The abstract predicate is only as meaningful as the flag supplied}]\label{l19:warn:freeflag}
Because \lean{F} is an arbitrary \lean{GradedFlag}, the bare statement ``there is a flag $F$ with
\lean{GradedUnitary B R F}'' is nearly empty. \emph{Paper argument (not machine-checked):} suppose $B$
is nondegenerate on each $V_n$, $V_0 = \C\lnineteenvac$, and $\rho$ is a real structure with $s\equiv1$
for which $h_n$ is Hermitian (as for the Virasoro VOAs below). For $n>0$ choose the flag
$\lnineteenG_1V_n = N_n$, a maximal negative-definite subspace of $h_n$ spanned by $\rho$-fixed vectors,
and $\lnineteenG_pV_n = V_n$ for $p\ge2$. Then $\lnineteenLay_1 = N_n$ is negative, $\lnineteenLay_2$ is
the $h_n$-orthogonal complement of $N_n$, which is positive, and higher layers vanish; clause (iv)
holds because $h_n$ is nondegenerate on $N_n$ and on $V_n$. So \emph{every} such VOA --- including
unitary ones after a sign change, and non-unitary ones like Lee--Yang --- is ``graded unitary for some
flag''. The content of graded unitarity lies entirely in the choice of the flag. The library's
Virasoro statement (\lean{VirGradedUnitary}, Definition~\ref{l19:def:virgu}) correctly fixes the
length flag, so it is not affected; but any future theorem of the form
``$\exists F$, \lean{GradedUnitary B R F}'', or any transfer theorem whose flags are free, must be read with
this in mind.
\end{warning}

\begin{theorem}[The two orthogonal complements agree on $\rho$-stable subspaces]\label{l19:thm:orthsame}
Let $S\le V_n$ be $\rho$-stable ($\rho(S)\subseteq S$) and suppose $s(n)\ne0$. Then
$\{y : h_n(x,y) = 0\ \forall x\in S\} = \{y : B(x,y) = 0\ \forall x\in S\}$.
\end{theorem}

\noindent File \leanpath{VOA/\allowbreak Core/\allowbreak GradedUnitarity.lean}.
\begin{leancode}
def RhoStableOn (R : RealStructure W) {n : ℤ} (S : Submodule ℂ (W.grading n)) : Prop :=
  ∀ x ∈ S, ∃ y ∈ S, (y : 𝒱) = R.rho (x : 𝒱)

theorem orthogonal_gradedSesq_eq (B : InvariantForm W) (R : RealStructure W) {n : ℤ}
    {S : Submodule ℂ (W.grading n)} (hS : RhoStableOn R S) (hp : R.parity n ≠ 0) :
    (gradedSesq B R n).orthogonal S = (LiRadical.gradedForm B n).orthogonal S := by
\end{leancode}

\begin{dictionary}
\lean{RhoStableOn R S}: for every $x\in S$ the vector $\rho x$ (computed in $V$) lies in $S$.
\lean{SesqForm.orthogonal} is defined in the file with the convention ``killed in the first slot'',
matching the bilinear convention. The statement is exactly Theorem~\ref{l19:thm:orthsame}.
\end{dictionary}

\paragraph{Proof on paper.} If $B(x,y)=0$ for all $x\in S$ then $h(x,y) = B(\rho x,y) = 0$ since
$\rho x\in S$. Conversely, if $h(z,y) = 0$ for all $z\in S$, take $z = \rho x$: $0 = B(\rho^2x,y) =
s(n)B(x,y)$, and $s(n)\ne0$. \qed

\begin{physics}
Whether one says ``orthogonal'' with the BPZ bilinear pairing or with the Hilbert-space inner
product makes no difference on a sector closed under conjugation; so the ``exact-charge'' states are
well defined independently of which form one uses.
\end{physics}

\begin{status}
Proved, axiom-clean.
\end{status}

%======================================================================
\section{The layer reduction: clause (v) follows from clause (iv)}\label{l19:sec:reduction}
%======================================================================

\begin{theorem}[Gram--Schmidt along a nondegenerate flag]\label{l19:thm:layerdec}
Let $V_n$ be finite-dimensional, $\beta$ a reflexive bilinear form on $V_n$ (i.e.\ $\beta(x,y) = 0\iff
\beta(y,x) = 0$), and assume clause (iv). Then for every $p$
\begin{equation*}
 \lnineteenG_{p-1}V_n + \lnineteenLay_p = \lnineteenG_pV_n,\qquad
 \lnineteenG_{p-1}V_n\cap\lnineteenLay_p = 0,\qquad
 \dim\lnineteenG_{p-1}V_n + \dim\lnineteenLay_p = \dim\lnineteenG_pV_n,
\end{equation*}
the layers span $V_n$, and if $V_n\ne0$ some layer is nonzero.
\end{theorem}

\noindent File \leanpath{VOA/\allowbreak Core/\allowbreak GradedUnitarity.lean}.
\begin{leancode}
theorem layerDecomposition [FiniteDimensional ℂ (W.grading n)] (href : β.IsRefl)
    (hnd : LayerNondegenerate F n β) (p : ℕ) :
    F.pieceLT p n ⊔ F.layer n β p = F.pieceIn p n ∧
      Disjoint (F.pieceLT p n) (F.layer n β p) :=

theorem finrank_pieceLT_add_layer [FiniteDimensional ℂ (W.grading n)] (href : β.IsRefl)
    (hnd : LayerNondegenerate F n β) (p : ℕ) :
    finrank ℂ (F.pieceLT p n) + finrank ℂ (F.layer n β p) = finrank ℂ (F.pieceIn p n) :=

theorem iSup_layer_eq_top [FiniteDimensional ℂ (W.grading n)] (href : β.IsRefl)
    (hnd : LayerNondegenerate F n β) : (⨆ p : ℕ, F.layer n β p) = ⊤ := by

theorem exists_layer_ne_bot [FiniteDimensional ℂ (W.grading n)] (href : β.IsRefl)
    (hnd : LayerNondegenerate F n β) (hne : (⊤ : Submodule ℂ (W.grading n)) ≠ ⊥) :
    ∃ p : ℕ, F.layer n β p ≠ ⊥ := by
\end{leancode}

\begin{dictionary}
\lean{⊔} is the sum of subspaces, \lean{Disjoint} means zero intersection, \lean{⨆ p} the span of all
layers, \lean{⊤} all of $V_n$. \lean{[FiniteDimensional ℂ (W.grading n)]} is an instance argument
(supplied automatically from the VOA axiom that weight spaces are finite-dimensional).
\lean{β.IsRefl} is reflexivity; for the canonical forms it follows from symmetry, which the library
proves (\lean{VirQuotForm.}\allowbreak\lean{gradedForm_}\allowbreak\lean{Bsimp_}\allowbreak\lean{isRefl} for Virasoro). Exactly
Theorem~\ref{l19:thm:layerdec}. The point of the theorem: ABLR's clause (v) is \emph{not} an extra
hypothesis once (iv) holds; and positivity cannot be satisfied vacuously by all layers being zero.
\end{dictionary}

\paragraph{Proof on paper.} If $\beta$ is nondegenerate on $S = \lnineteenG_{p-1}V_n$ then
$V_n = S\oplus S^\perp$ (standard linear algebra for a reflexive form on a finite-dimensional
space). Intersecting with $T = \lnineteenG_pV_n\supseteq S$ gives $T = S\oplus(T\cap S^\perp)$ by the
modular law. Dimensions add. By induction on $p$ each $\lnineteenG_pV_n$ lies in the span of the layers, and
by exhaustiveness so does $V_n$. If all layers were zero, $V_n$ would be zero. \qed

\begin{physics}
Gram--Schmidt orthogonalisation of states ordered by $R$-charge: the space of states is the
orthogonal direct sum of the exact-charge sectors, and the number of states of exact charge $p$ is
the jump of the cumulative counts. Combined with a refined character (counting states by weight and
length, Salazar's theorem cited in \lean{AdmissiblePartitions}), this is what turns graded unitarity
into the alternating-sum signature formula $\operatorname{sig} = \sum_r(-1)^ra_{N,r}$ studied
numerically by the project.
\end{physics}

\begin{status}
All four proved, axiom-clean; unconditional (the hypotheses are the clause-(iv) assumption and
reflexivity, both explicit). \lean{nondegenerate_restrict_bot} handles $p = 0$.
\end{status}

The next theorem is a \emph{calibration check}, not a step of the positivity argument. Its job is to show that
the repaired, sesquilinear encoding of clause (vi) is satisfiable on a nonzero layer of a genuine VOA --- exactly
what the old bilinear encoding could never achieve (Theorem~\ref{l19:thm:bilpos}). It is used once later, as the
trivial layer $p=0$ of the through-two theorem (Theorem~\ref{l19:thm:throughtwo}); everything with mathematical
content happens in layers $p\ge1$.

\begin{theorem}[The vacuum slice of clause (vi); calibration check]\label{l19:thm:vacuum}
For every VOA $W$ with $B(\lnineteenvac,\lnineteenvac) = 1$, every real structure, every flag and every
form $\beta$ on $V_0$, the twisted form $(-1)^0h_0$ is positive on the bottom layer
$\lnineteenLay_0\subseteq V_0$, which equals $\C\lnineteenvac$; if $\lnineteenvac\ne0$ this layer is nonzero.
\end{theorem}

\noindent File \leanpath{VOA/\allowbreak Core/\allowbreak GradedUnitarity.lean}.
\begin{leancode}
theorem isPosOn_layer_zero_vacuum (B : InvariantForm W) (R : RealStructure W) (F : GradedFlag W)
    (β : BilinForm ℂ (W.grading 0)) (hn : B.Normalized) :
    IsPosOn (signedForm (gradedSesq B R 0) 0) (F.layer 0 β 0) := by

theorem layer_zero_ne_bot_and_isPosOn (B : InvariantForm W) (R : RealStructure W)
    (F : GradedFlag W) (β : BilinForm ℂ (W.grading 0)) (hn : B.Normalized)
    (hvac : W.vacuum ≠ 0) :
    F.layer 0 β 0 ≠ ⊥ ∧ IsPosOn (signedForm (gradedSesq B R 0) 0) (F.layer 0 β 0) :=
\end{leancode}

\begin{dictionary}
\lean{B.Normalized} is $B(\lnineteenvac,\lnineteenvac) = 1$ (\leanpath{VOA/\allowbreak Core/\allowbreak BilinearForm.lean}).
The bottom layer is \lean{F.layer 0 β 0} $= \lnineteenG_0V_0 = \C\lnineteenvac$ (a flag axiom). Exactly
Theorem~\ref{l19:thm:vacuum}. Its role is calibration, not progress: it certifies that under the
repaired encoding clause (vi) holds on a nonzero layer of a genuine VOA, where the old encoding was
false, and that the sign $(-1)^0 = +1$ at $p=0$ is consistent with $(\lnineteenvac|\lnineteenvac) = +1$
(\lean{signedForm_}\allowbreak\lean{gradedSesq_}\allowbreak\lean{vacuum}, \lean{pos_gradedSesq_vacuum}).
\end{dictionary}

\paragraph{Proof on paper.} $x = z\lnineteenvac$ gives $h(x,x) = B(\bar z\lnineteenvac,z\lnineteenvac) =
|z|^2$. \qed

\begin{physics}
The vacuum has norm one.
\end{physics}

\begin{status}
Proved, axiom-clean, unconditional. The docstring warns that this ``is not progress towards graded
unitarity'', which is the correct reading: the theorem certifies that, under the repaired sesquilinear
encoding, clause (vi) holds on a nonzero layer of a genuine VOA (where the old bilinear encoding was
impossible by Theorem~\ref{l19:thm:bilpos}). It is a consistency check of the definitions, not a step towards
unconditional graded unitarity; the only later use is as the trivial layer $p = 0$ in
Theorem~\ref{l19:thm:throughtwo}.
\end{status}

%======================================================================
\section{The minors criterion}\label{l19:sec:minors}
%======================================================================

Clauses (iv) and (vi) are statements about subspaces. When the flag is the family of initial
segments of an ordered basis $e_1,\dots,e_m$ of $V_n$, both become statements about the
\emph{leading principal minors} $D_k = \det(\beta(e_i,e_j))_{i,j\le k}$ of the Gram matrix.

\subsection{Clause (iv) and nonzero minors}

\begin{definition}[Leading principal submatrix; flags refining a basis]\label{l19:def:refines}
For an $m\times m$ matrix $M$ and $k\le m$, $M_{\le k}$ is the top-left $k\times k$ block. A family of
subspaces $F_p$ \emph{refines} the ordered basis $e$ if there are integers $d_p\le m$ with
$F_p = \operatorname{span}(e_1,\dots,e_{d_p})$.
\end{definition}

\noindent File \leanpath{VOA/\allowbreak Core/\allowbreak FlagMinors.lean}.
\begin{leancode}
def leadingPrincipal (M : Matrix (Fin m) (Fin m) K) {k : ℕ} (hk : k ≤ m) :
    Matrix (Fin k) (Fin k) K :=
  M.submatrix (Fin.castLE hk) (Fin.castLE hk)

structure RefinesBasis (F : ℕ → Submodule K E) (e : Module.Basis (Fin m) K E) where
  /-- The dimension of the `p`-th step. -/
  dim : ℕ → ℕ
  /-- Each step dimension is at most the dimension of the ambient space. -/
  dim_le : ∀ p, dim p ≤ m
  /-- Step `p` is spanned by the first `dim p` basis vectors, **in the given order**. -/
  eq : ∀ p, F p = initFlag e (dim_le p)
\end{leancode}

\begin{dictionary}
\lean{Fin m} is $\{0,\dots,m-1\}$; \lean{Fin.castLE hk} embeds $\{0,\dots,k-1\}$ into it, so
\lean{M.submatrix (Fin.castLE hk) (Fin.castLE hk)} is the leading $k\times k$ block.
\lean{Module.Basis (Fin m) K E} is an ordered basis; \lean{initFlag e hk} is the span of its first $k$
vectors (defined in the same file). $K$ is any field. Exactly Definition~\ref{l19:def:refines}.
\end{dictionary}

\begin{physics}
An ordered basis of states sorted by $R$-charge (PBW length), so that each charge-$\le p$ sector is
spanned by an initial segment.
\end{physics}

\begin{status}
Definitions.
\end{status}

\begin{theorem}[Clause (iv) $\iff$ nonzero leading minors]\label{l19:thm:minors}
If $F$ refines the ordered basis $e$, then $\beta$ is nondegenerate on every $F_p$ if and only if
$\det\big(\mathrm{Gram}_e(\beta)_{\le d_p}\big)\ne0$ for every $p$.
\end{theorem}

\noindent File \leanpath{VOA/\allowbreak Core/\allowbreak FlagMinors.lean}.
\begin{leancode}
theorem nondegenerate_flag_iff_minors {F : ℕ → Submodule K E} {e : Module.Basis (Fin m) K E}
    (h : RefinesBasis F e) (β : BilinForm K E) :
    (∀ p, (β.restrict (F p)).Nondegenerate)
      ↔ ∀ p, (leadingPrincipal (LinearMap.BilinForm.toMatrix e β) (h.dim_le p)).det ≠ 0 := by
\end{leancode}

\begin{dictionary}
\lean{LinearMap.BilinForm.toMatrix e β} is the Gram matrix $(\beta(e_i,e_j))$. The theorem is an
equivalence; exactly Theorem~\ref{l19:thm:minors}. The file docstring warns that the
\emph{ordering} matters for the companion sign statement: at $c=-22/5$, level $6$, in the order
$(L_{-6}\lnineteenvac,L_{-4}L_{-2}\lnineteenvac)$ the Gram matrix is
$\begin{psmallmatrix}-77&-22\\-22&66/5\end{psmallmatrix}$ with pivots $(-77,\,682/35)$, while in the
reversed order the pivots are $(66/5,\,-341/3)$. (We recomputed these: $(L_{-6}\lnineteenvac|L_{-6}\lnineteenvac)
= \tfrac{c}{12}(6^3-6) = \tfrac{35c}{2}$, $(L_{-6}\lnineteenvac|L_{-4}L_{-2}\lnineteenvac) = 5c$,
$(L_{-4}L_{-2}\lnineteenvac|L_{-4}L_{-2}\lnineteenvac) = \tfrac c2(5c+16)$, and $\det = -7502/5$.)
\end{dictionary}

\paragraph{Proof on paper.} The Gram matrix of $\beta|_{F_p}$ in the basis $e_1,\dots,e_{d_p}$ is the
leading block (\lean{toMatrix_restrict_initFlag}), and a bilinear form on a finite-dimensional space
is nondegenerate iff its Gram determinant is nonzero. \qed

\begin{physics}
``No null state in any truncated charge sector'' is the non-vanishing of the cumulative Kac-type
determinants.
\end{physics}

\begin{status}
Proved, axiom-clean, unconditional.
\end{status}

\subsection{Clause (vi) and the signs of minor ratios}

Over $\lnineteenR$, Sylvester's criterion says a symmetric matrix is positive definite iff all its leading
minors are positive. Mathlib has only the easy direction (\lean{Matrix.PosDef.det_pos}). The library
supplies the converse and its layered refinement.

\begin{theorem}[Jacobi's quotient identity, frame invariance, converse Sylvester]\label{l19:thm:sylvester}
Let $M$ be an $m\times m$ matrix over a field.
\begin{enumerate}
\item If $D_l = \det M_{\le l}$ is invertible, then $D_{l+j} = D_l\cdot\det S_{l,j}$, where $S_{l,j}$ is
  the Schur complement of $M_{\le l}$ in $M_{\le l+j}$.
\item If $U$ is upper unitriangular, $U^TMU$ has the same leading minors as $M$.
\item Over $\lnineteenR$: if $M$ is symmetric and $D_k>0$ for all $1\le k\le m$, then $M$ is positive definite.
\item Over $\lnineteenR$, $M$ symmetric: if $D_l>0$ and $D_{l+i}/D_l>0$ for $1\le i\le j$, then $S_{l,j}$ is
  positive definite; if instead $(-1)^iD_{l+i}/D_l>0$, then $-S_{l,j}$ is positive definite.
\end{enumerate}
\end{theorem}

\noindent File \leanpath{VOA/\allowbreak Core/\allowbreak MinorPivots.lean}.
\begin{leancode}
theorem det_leadingPrincipal_eq_mul_det_schurOfLeading
    (M : Matrix (Fin m) (Fin m) K) (l j : ℕ) (hlj : l + j ≤ m)
    (hl : l ≤ m) (hdet : IsUnit (leadingPrincipal M hl).det) :
    (leadingPrincipal M hlj).det
      = (leadingPrincipal M hl).det * (schurOfLeading M l j hlj).det := by

theorem det_leadingPrincipal_unitriangular_congr
    (B : Matrix (Fin m) (Fin m) K) {U : Matrix (Fin m) (Fin m) K}
    (hU : U.BlockTriangular id) (hdiag : ∀ i, U i i = 1)
    {k : ℕ} (hk : k ≤ m) :
    (leadingPrincipal (Uᵀ * B * U) hk).det = (leadingPrincipal B hk).det := by

theorem posDef_of_det_leadingPrincipal_pos :
    ∀ {m : ℕ} (M : Matrix (Fin m) (Fin m) ℝ), M.IsHermitian →
      (∀ k (hk : k ≤ m), 0 < k → 0 < (leadingPrincipal M hk).det) →
      M.PosDef := by

theorem schurOfLeading_posDef_of_sign_ratios
    (M : Matrix (Fin m) (Fin m) ℝ) (hM : M.IsHermitian) (l j : ℕ)
    (hlj : l + j ≤ m) (hl : l ≤ m)
    (hDl : (leadingPrincipal M hl).det ≠ 0)
    (hratios : ∀ i (hi0 : 0 < i) (hli : l + i ≤ m), i ≤ j →
      0 < (leadingPrincipal M hli).det / (leadingPrincipal M hl).det) :
    (schurOfLeading M l j hlj).PosDef := by

theorem schurOfLeading_negDef_of_sign_ratios
    (M : Matrix (Fin m) (Fin m) ℝ) (hM : M.IsHermitian) (l j : ℕ)
    (hlj : l + j ≤ m) (hl : l ≤ m)
    (hDl : (leadingPrincipal M hl).det ≠ 0)
    (hratios : ∀ i (hi0 : 0 < i) (hli : l + i ≤ m), i ≤ j →
      0 < (-1 : ℝ) ^ i
        * ((leadingPrincipal M hli).det / (leadingPrincipal M hl).det)) :
    (-(schurOfLeading M l j hlj)).PosDef := by
\end{leancode}

\begin{dictionary}
\lean{schurOfLeading M l j hlj} is defined in the file as
$M_{TT} - M_{TL}(M_{\le l})^{-1}M_{LT}$ with $L = \{1..l\}$, $T = \{l+1..l+j\}$.
\lean{U.BlockTriangular id} means upper triangular; \lean{Uᵀ} transpose. \lean{M.IsHermitian} over
$\lnineteenR$ is symmetry; \lean{M.PosDef} is Mathlib's positive definiteness ($x^TMx>0$ for $x\ne0$, plus
symmetry). Items 1--4 of Theorem~\ref{l19:thm:sylvester} are the four groups of declarations in
order (the last two make item 4).

\textbf{A hypothesis stronger than needed --- since weakened.} When this chapter was written, item 4 assumed $D_l>0$;
the current statements, quoted above, assume only \lean{hDl : (leadingPrincipal M hl).det ≠ 0}, which is all that is used (the
leading minors of $S_{l,j}$ are $D_{l+i}/D_l$ whatever the sign of $D_l$). In graded unitarity
the cumulative minors alternate in sign, so $D_l<0$ happens (e.g.\ after a one-dimensional negative
layer). The library contains the general version elsewhere:
\lean{PolarizedRees.}\allowbreak\lean{schurOfLeading_}\allowbreak\lean{posDef_}\allowbreak\lean{of_}\allowbreak\lean{ratios_}\allowbreak\lean{of_}\allowbreak\lean{pivot_}\allowbreak\lean{ne_}\allowbreak\lean{zero} and its negative twin in
\leanpath{VOA/\allowbreak Examples/\allowbreak PolarizedReesStrongGenerator.lean}, hypothesis \lean{hDl : … ≠ 0}. Since the
library update the \leanpath{MinorPivots} versions have the same, optimal, hypothesis.
\end{dictionary}

\paragraph{Proof on paper.} (1) Block decomposition
$M_{\le l+j} = \begin{psmallmatrix}A&B\\C&D\end{psmallmatrix}$ with $A$ invertible gives
$\det = \det A\cdot\det(D-CA^{-1}B)$. (2) The leading block of $U^TMU$ is
$U_{\le k}^TM_{\le k}U_{\le k}$ because $U$ is triangular, and $\det U_{\le k} = 1$. (3) Induction on
$m$: $A = M_{\le m-1}$ is positive definite by induction; the $1\times1$ Schur complement has determinant
$D_m/D_{m-1}>0$; a symmetric block matrix with positive-definite $A$ is positive definite iff its Schur
complement is (\lean{fromBlocks_}\allowbreak\lean{posDef_}\allowbreak\lean{iff_}\allowbreak\lean{schur₁₁}, \leanpath{VOA/\allowbreak Core/\allowbreak HermitianSchur.lean}).
(4) The leading $i\times i$ block of $S_{l,j}$ is $S_{l,i}$ (\lean{leadingPrincipal_}\allowbreak\lean{schurOfLeading}),
whose determinant is $D_{l+i}/D_l$ by (1); apply (3) to $\pm S_{l,j}$. \qed

\begin{physics}
The Schur complement $S_{l,j}$ is the Gram matrix of the states of the next charge sector after
projecting out the lower-charge states --- i.e.\ the Gram matrix of the layer. Item 4 says: graded
unitarity at one weight is equivalent to the statement that the pivots
$D_k/D_{k-1}$ have sign $(-1)^{(\text{length of the }k\text{-th basis state})}$, a condition one can
test with exact rational arithmetic, which is how the project's numerical campaign was run.
\end{physics}

\begin{status}
All proved, axiom-clean, unconditional. Also proved: \lean{fromBlocks_}\allowbreak\lean{posDef_}\allowbreak\lean{iff_}\allowbreak\lean{schur₁₁}
(\leanpath{HermitianSchur}), \lean{realSymmetricTwo_}\allowbreak\lean{posDef_}\allowbreak\lean{iff} and
\lean{realSymmetricTwo_}\allowbreak\lean{posSemidef_}\allowbreak\lean{iff} (the $2\times2$ Sylvester criteria,
\leanpath{VOA/\allowbreak Core/\allowbreak RealSymmetricTwo.lean}). The docstrings of \leanpath{FlagMinors.lean},
\leanpath{GradedUnitarity.lean} and \leanpath{VirasoroGradedUnitarity.lean} still say that ``the
pivot/positivity half is absent'' and ``not reducible to minors with what Mathlib currently has'';
this is \textbf{stale} since \leanpath{MinorPivots.lean} landed (documentation defect, see the LEAN
DEFECTS list at the end of the chapter).
\end{status}

\begin{example}[Lee--Yang, weight $6$]\label{l19:ex:ly6}
With $M = \begin{psmallmatrix}-77&-22\\-22&66/5\end{psmallmatrix}$, $D_1 = -77<0$ and
$D_2/D_1 = 682/35>0$. The first basis vector $L_{-6}\lnineteenvac = \tfrac1{24}L_{-1}^4\omega$ has length
one and the second has length two, so the predicted signs are $(-1)^1$ and $(-1)^2$: the
pivots agree. This is the content of graded unitarity of $L(-22/5,0)$ at weight six, and it is a
special case of the length-two theorem of Section~\ref{l19:sec:lengthtwo}.
\end{example}

%======================================================================
\section{Graded unitarity of the Virasoro VOAs: statement and reduction}\label{l19:sec:vir}
%======================================================================

\subsection{The length flag}

Recall (Chapter~\ref{ch:L06}) that $V(c,0)$ is the universal Virasoro vacuum module, with basis the
sorted creation words: \lean{negWordVir c w} for a weakly decreasing list $w = [i_1,\dots,i_k]$ of
natural numbers is $L_{-(i_1+2)}\cdots L_{-(i_k+2)}\lnineteenvac$, of weight
\lean{energyVir w} $= \sum_j(i_j+2)$ and length $k$; and $L(c,0) = V(c,0)/\mathcal M_c$ is the simple
quotient by the maximal proper ideal \lean{maxIdealVir c}, with canonical normalised invariant form
\lean{Bsimp c} (Chapter~\ref{ch:L07}, \S\ref{l07:sec:kac}).

\begin{definition}[Length filtration of $L(c,0)$; ABLR (3.11), Kaufmann $\tilde V_{n,k}$]\label{l19:def:lengthflag}
$\mathcal F_pV(c,0)_n$ is the span of the sorted creation words of weight $n$ and length $\le p$, and
$\lnineteenG_pL(c,0)_n$ is its image in the simple quotient.
\end{definition}

\noindent File \leanpath{VOA/\allowbreak Examples/\allowbreak VirasoroGoodFiltration.lean} (a file of
Chapter~\ref{ch:L06}).
\begin{leancode}
noncomputable def virFil (p : ℕ) (n : ℤ) : Submodule ℂ (VirUniv c) :=
  Submodule.span ℂ
    { x | ∃ w : List ℕ, SortedWord w ∧ (energyVir w : ℤ) = n ∧ w.length ≤ p ∧
            negWordVir c w = x }

noncomputable def goodFil (p : ℕ) (n : ℤ) : Submodule ℂ (VirUniv c ⧸ maxIdealVir c) :=
  (virFil c p n).map (maxIdealVir c).mkQ

noncomputable def virGradedFlag : GradedFlag (virasoroSimple c) where
  piece := goodFil c
  le_grading := goodFil_le_grading c
  mono := fun n h => goodFil_mono c n h
  bot := fun _ hn => goodFil_zero_of_pos c hn
  vacuum_zero := goodFil_zero_zero c
  exhaustive := iSup_goodFil c
\end{leancode}

\begin{dictionary}
\lean{SortedWord w} is \lean{w.Pairwise (· ≥ ·)} (weakly decreasing). \lean{(maxIdealVir c).mkQ} is
the quotient map $V(c,0)\to L(c,0)$, and \lean{.map} takes the image. The \lean{GradedFlag} fields are
the axioms of Chapter~\ref{ch:L11}, all proved for every $c$. Note the flag is an \emph{image}: no
basis of $L(c,0)$ is assumed. Whether the length filtration \emph{is} ABLR's $R$-filtration for the
boundary models is part of the conjecture's physics, not of the library.
\end{dictionary}

\begin{physics}
Each $L_{-n}$ carries one unit of $SU(2)_R$ charge (the stress tensor comes from the 4d
$SU(2)_R$ current multiplet, of $R=1$), so a word of length $k$ has $R\le k$; the length filtration is
the natural candidate for the $R$-filtration of the Virasoro chiral algebras of the
$(A_1,A_{2n})$ Argyres--Douglas theories.
\end{physics}

\begin{status}
Definitions with all fields proved (Chapter~\ref{ch:L06} files), axiom-clean.
\end{status}

\begin{definition}[Graded unitarity of $L(c,0)$]\label{l19:def:virgu}
$L(c,0)$ is graded unitary with respect to a real structure $\rho$ if
\lean{GradedUnitary} holds for the canonical form, $\rho$, and the length flag.
\end{definition}

\noindent File \leanpath{VOA/\allowbreak Examples/\allowbreak VirasoroGradedUnitarity.lean}.
\begin{leancode}
def VirGradedUnitary (R : RealStructure (virasoroSimple c)) : Prop :=
  GradedUnitary (Bsimp c) R (virGradedFlag c)
\end{leancode}

\begin{dictionary}
\lean{virasoroSimple c} is the simple Virasoro VOA $L(c,0)$, \lean{Bsimp c} its canonical form with
$(\lnineteenvac|\lnineteenvac) = 1$ (\lean{Bsimp_normalized}). The flag is \emph{pinned} to the length flag,
so Warning~\ref{l19:warn:freeflag} does not apply here. The real structure is a parameter; canonical
choices are constructed in Section~\ref{l19:sec:virreal}. The statement that
\lean{VirGradedUnitary (boundaryCharge n) R} holds for $c = \lnineteencb$ is \textbf{Kaufmann's
conjecture} (1994), open; the library does not claim it.
\end{dictionary}

\begin{physics}
``The Virasoro chiral algebra of the $(A_1,A_{2n})$ Argyres--Douglas theory satisfies all the
positivity constraints inherited from four-dimensional unitarity.''
\end{physics}

\begin{status}
Definition. Never asserted of any $c$. \lean{virGradedUnitary_iff} unfolds it weight by weight;
\lean{isPosOn_}\allowbreak\lean{layer_}\allowbreak\lean{zero_}\allowbreak\lean{vacuum_}\allowbreak\lean{simple} proves the weight-zero slice of clause (vi) for every $c$ and
every $\rho$.
\end{status}

\begin{theorem}[The Virasoro reduction to minors]\label{l19:thm:virreduction}
Let $N\ge0$ and let $e$ be an ordered basis of $L(c,0)_N$ whose initial segments are the steps of the
length flag. For every sesquilinear form $h$ on $L(c,0)_N$, graded unitarity at weight $N$ holds iff
all leading principal minors of the Gram matrix of \lean{Bsimp c} in the basis $e$ (cut at the step
dimensions) are nonzero and the layers are positive. Moreover, if each $e_i$ is the image of a
creation word $w_i$, the Gram entries are the universal values $(L_{-w_i}\lnineteenvac|L_{-w_j}\lnineteenvac)$
of the canonical form on $V(c,0)$.
\end{theorem}

\noindent File \leanpath{VOA/\allowbreak Examples/\allowbreak VirasoroGradedUnitarity.lean}.
\begin{leancode}
theorem gradedUnitaryAt_iff_minors_and_positive {c : ℂ} {N m : ℕ}
    (e : Module.Basis (Fin m) ℂ ((virasoroSimple c).grading (N : ℤ)))
    (hrb : FlagMinors.RefinesBasis (fun p => (virGradedFlag c).pieceIn p (N : ℤ)) e)
    (h : SesqForm ((virasoroSimple c).grading (N : ℤ))) :
    GradedUnitaryAt (virGradedFlag c) (N : ℤ) (LiRadical.gradedForm (Bsimp c) (N : ℤ)) h
      ↔ (∀ p : ℕ, (FlagMinors.leadingPrincipal
            (LinearMap.BilinForm.toMatrix e (LiRadical.gradedForm (Bsimp c) (N : ℤ)))
            (hrb.dim_le p)).det ≠ 0) ∧
        LayerPositive (virGradedFlag c) (N : ℤ) (LiRadical.gradedForm (Bsimp c) (N : ℤ)) h :=

theorem toMatrix_eq_canB_of_monomials (w : Fin m → List ℕ)
    (hw : ∀ i, (e i : VirUniv c ⧸ maxIdealVir c)
      = (maxIdealVir c).mkQ (negWordVir c (w i))) (i j : Fin m) :
    LinearMap.BilinForm.toMatrix e (LiRadical.gradedForm (Bsimp c) (N : ℤ)) i j
      = (canB c).form (negWordVir c (w i)) (negWordVir c (w j)) := by
\end{leancode}

\begin{dictionary}
\lean{hrb} is the hypothesis that the flag refines $e$; a basis of the simple quotient adapted to the
flag is \emph{not} constructed in the library --- it is the content of the cited
\lean{AdmissiblePartitions.}\allowbreak\lean{FlagBasisHyp} (Salazar's admissible-partition basis; Chapter~\ref{ch:L09}
explains this hypothesis), so the theorem is a reduction \emph{given} such a basis. The right side
keeps \lean{LayerPositive} untouched: the sign half is \emph{not} rewritten into minors here (the
pivot machinery of Theorem~\ref{l19:thm:sylvester} would do it, but no Virasoro statement composes the
two). In \lean{toMatrix_}\allowbreak\lean{eq_}\allowbreak\lean{canB_}\allowbreak\lean{of_}\allowbreak\lean{monomials} the section variables are an ordered basis \lean{e} of
the weight-$N$ space and $m$ its size; \lean{canB c} is the canonical form of $V(c,0)$ and the identity
is \lean{VirQuotForm.canB_quot_mkQ}: the form on the quotient is computed by representatives.
\end{dictionary}

\paragraph{Proof on paper.} Combine Theorem~\ref{l19:thm:minors} with the definition
\lean{GradedUnitaryAt} = (iv) $\wedge$ (vi). The second statement is the definition of the quotient form.
\qed

\begin{physics}
To test graded unitarity of a minimal model at a given level one needs only the Kac-type Gram matrix
of creation monomials (computable in the universal module) and an admissible basis of the quotient.
\end{physics}

\begin{status}
Proved, axiom-clean; conditional on the explicit hypothesis \lean{hrb} (an adapted basis), never
discharged for any $N$ in the library. Also proved: \lean{virLayerDecomposition},
\lean{virFinrank_layer}, \lean{virISup_layer_eq_top} (Theorem~\ref{l19:thm:layerdec} at $L(c,0)$,
conditional on clause (iv)).
\end{status}

\subsection{Real structures on the Virasoro VOAs}\label{l19:sec:virreal}

\begin{definition}[Partition-basis conjugation]\label{l19:def:virconj}
On $V(c,0)$, \lean{virBasisConj c} conjugates the coordinates in the basis of sorted creation words:
$\sum_w a_wL_{-w}\lnineteenvac\mapsto\sum_w\bar a_wL_{-w}\lnineteenvac$.
\end{definition}

\noindent File \leanpath{VOA/\allowbreak Examples/\allowbreak VirasoroFockConjugation.lean} (Chapter~\ref{ch:L06}).
\begin{leancode}
noncomputable def virBasisConj (c : ℂ) : VirUniv c →ₛₗ[starRingEnd ℂ] VirUniv c :=
  basisConj (sortedBasisVir c)
\end{leancode}

\begin{dictionary}
\lean{→ₛₗ[starRingEnd ℂ]} is a semilinear (here: anti-linear) map. \lean{sortedBasisVir c} is the
sorted-word basis of $V(c,0)$ (Chapter~\ref{ch:L06}; its linear independence is the Feigin--Fuchs
argument, unconditional). It is an involution (\lean{virBasisConj_sq}) at every $c$; whether it
commutes with the modes is the question answered next.
\end{dictionary}

\begin{physics}
``Complex conjugation of wave-functions in the PBW basis.''
\end{physics}

\begin{status}
Definition.
\end{status}

\begin{theorem}[Real structures on $V(c,0)$ and $L(c,0)$ for $c\le1$]\label{l19:thm:virreal}
Let $\lambda\in\lnineteenR$ and $c = 1-12\lambda^2$. Then the partition-basis conjugation is a real structure on
$V(c,0)$ with parity $s\equiv1$; it preserves the maximal ideal and descends to a real structure on
$L(c,0)$.
\end{theorem}

\noindent File \leanpath{VOA/\allowbreak Examples/\allowbreak VirasoroRealStructure.lean}.
\begin{leancode}
noncomputable def virasoroUniversalRealStructure {c lam : ℂ}
    (hlam : 1 - 12 * lam ^ 2 = c) (hreal : (starRingEnd ℂ) lam = lam) :
    RealStructure (virasoroVOA c) where

noncomputable def virasoroSimpleRealStructure {c lam : ℂ}
    (hlam : 1 - 12 * lam ^ 2 = c) (hreal : (starRingEnd ℂ) lam = lam) :
    RealStructure (virasoroSimple c) where
\end{leancode}

\begin{dictionary}
A Lean \lean{def ... : RealStructure ... where} \emph{constructs} the structure: every field of
Definition~\ref{l19:def:real} is proved (we omit the field proofs). \lean{hreal} says $\bar\lambda =
\lambda$, i.e.\ $\lambda$ is real. So the hypotheses say: $c = 1-12\lambda^2$ for a real $\lambda$,
equivalently $c$ real and $c\le1$. Parity is \lean{fun _ => 1}.

\textbf{A restriction, not a defect.} Mathematically the partition-basis conjugation is a real
structure for every real $c$ (the structure constants of the Virasoro algebra in the PBW basis are real
polynomials in $c$). The library's proof goes through the Feigin--Fuchs free-field realisation, which
needs a real background charge $\lambda$; for real $c>1$ ($\lambda$ imaginary) no real structure is
constructed. All the boundary models have $c<0$, so nothing used later is affected.
\end{dictionary}

\paragraph{Proof on paper.} With $\lambda$ real, the Feigin--Fuchs embedding $V(c,0)\hookrightarrow$ Fock
(Chapter~\ref{ch:L06}, Theorem~\ref{l06:thm:ffembed}) intertwines \lean{virBasisConj} with coefficientwise
conjugation of Fock polynomials, which commutes with the free-boson Virasoro modes $L_n^\lambda$ because
their coefficients are real. Hence $\rho L_n = L_n\rho$ on $V(c,0)$ (\lean{virBasisConj_Lop}). Since
$V(c,0)$ is strongly generated by $\omega$, every field $Y(u,z)$ is built from $Y(\omega,z)$ by
derivatives and normal-ordered products with real coefficients, and conjugation-covariance propagates
(\lean{VirConjCov.nprod}, \lean{virBasisConj_virYA}); this gives \lean{map_smul}. The preimage of a VOA
ideal under $\rho$ is an ideal not containing $\lnineteenvac$, hence $\rho(\mathcal M_c)\subseteq\mathcal M_c$
by maximality, and $\rho$ descends. \qed

\begin{physics}
The Virasoro stress tensor is a Hermitian field and its OPE has real coefficients; so ``complex
conjugation in the PBW basis'' is an anti-unitary symmetry of the chiral algebra. For $c\le1$ this is
manifest in the Coulomb-gas (Dotsenko--Fateev) description with real background charge $\alpha_0$.
\end{physics}

\begin{status}
Constructed; \texttt{\#print axioms VOA.virasoroSimpleRealStructure}: standard three. Unconditional
given the stated hypotheses.
\end{status}

\begin{theorem}[Canonical real structures on the boundary models]\label{l19:thm:boundaryreal}
For every $n\in\N$, $c_{2,2n+1} = 1-12\lambda_n^2$ with the real number
$\lambda_n = \sqrt{(2n-1)^2/(4(2n+1))}$, so $L(\lnineteencb,0)$ carries a canonical real structure.
\end{theorem}

\noindent File \leanpath{VOA/\allowbreak Examples/\allowbreak VirasoroBoundaryRealStructure.lean}.
\begin{leancode}
noncomputable def boundaryBackgroundChargeSq (n : ℕ) : ℝ :=
  (2 * (n : ℝ) - 1) ^ 2 / (4 * (2 * n + 1))

theorem boundaryBackgroundCharge_spec (n : ℕ) :
    1 - 12 * boundaryBackgroundCharge n ^ 2 = boundaryCharge n := by

noncomputable def boundaryVirasoroRealStructure (n : ℕ) :
    RealStructure (virasoroSimple (boundaryCharge n)) :=
  virasoroSimpleRealStructure (boundaryBackgroundCharge_spec n)
    (boundaryBackgroundCharge_fixed n)
\end{leancode}

\begin{dictionary}
\lean{boundaryBackgroundCharge n} is \lean{(Real.sqrt (boundaryBackgroundChargeSq n) : ℂ)}, and
\lean{boundaryCharge n} is the complex number $\lnineteencb = 1-3(2n-1)^2/(2n+1)$
(Definition~\ref{l19:def:boundarycharge} below). Check: $1-12\cdot\frac{(2n-1)^2}{4(2n+1)} =
1-\frac{3(2n-1)^2}{2n+1}$. Exactly Theorem~\ref{l19:thm:boundaryreal}.
\end{dictionary}

\begin{physics}
$\lambda_n$ is the Coulomb-gas background charge $\alpha_0/\sqrt2$ of $M(2,2n+1)$.
\end{physics}

\begin{status}
Constructed, axiom-clean, unconditional.
\end{status}

\begin{theorem}[Every step of the length flag is $\rho$-stable, for every real structure]\label{l19:thm:rhostable}
For every $c$, every real structure $\rho$ on $L(c,0)$, every $p$ and $N$, $\rho(\lnineteenG_pL(c,0)_N)\subseteq
\lnineteenG_pL(c,0)_N$.
\end{theorem}

\noindent File \leanpath{VOA/\allowbreak Examples/\allowbreak VirasoroLengthTwoBoundarySemantic.lean}.
\begin{leancode}
theorem rhoStableOn_vir_pieceIn {p : ℕ} {N : ℤ} :
    RhoStableOn R ((virGradedFlag c).pieceIn p N) := by
\end{leancode}

\begin{dictionary}
The section variables are \lean{{c : ℂ} (R : RealStructure (virasoroSimple c))}: an \emph{arbitrary}
real structure. Exactly the theorem. Consequence (via Theorem~\ref{l19:thm:orthsame}, with $s(N)\ne0$
--- e.g.\ $s\equiv1$): for the Virasoro flag the bilinear layers of \lean{GradedUnitary} \emph{are} ABLR's
sesquilinear layers.
\end{dictionary}

\paragraph{Proof on paper.} $\rho$ commutes with all $L_n$ (\lean{rho_L}) and fixes $\lnineteenvac$, so it
fixes every creation monomial; it maps the span of monomials of length $\le p$ to itself because it only
conjugates coefficients. \qed

\begin{physics}
Conjugation does not change the number of stress-tensor modes in a state.
\end{physics}

\begin{status}
Proved, axiom-clean, unconditional.
\end{status}

%======================================================================
\section{Length-two positivity for the boundary models}\label{l19:sec:lengthtwo}
%======================================================================

This is the main unconditional positivity theorem of the chapter. We first set up the arithmetic
and the quasi-primaries.

\subsection{The boundary central charges and the root \texorpdfstring{$\rho_m$}{rho m}}

\begin{definition}[Boundary central charge; the rational function $\rho_m$]\label{l19:def:boundarycharge}
For $n\in\N$, $\lnineteencb = 1-\dfrac{3(2n-1)^2}{2n+1}$. For $m\in\N$,
\begin{equation*}
 \rho_m = \frac{36(4m^2-2m-1)}{m(m-1)(m+1)(2m-3)(2m-1)(2m+1)}.
\end{equation*}
\end{definition}

\noindent File \leanpath{VOA/\allowbreak Examples/\allowbreak VirasoroLengthTwoArithmetic.lean}, then
\leanpath{VOA/\allowbreak Examples/\allowbreak VirasoroLengthTwoBoundary.lean}.
\begin{leancode}
def boundaryMinimalCentralCharge (n : ℕ) : ℚ :=
  1 - 3 * (2 * (n : ℚ) - 1) ^ 2 / (2 * n + 1)

def lengthTwoRho (m : ℕ) : ℚ :=
  36 * (4 * (m : ℚ) ^ 2 - 2 * m - 1) /
    ((m : ℚ) * (m - 1) * (m + 1) * (2 * m - 3) * (2 * m - 1) * (2 * m + 1))

theorem lengthTwoRho_mem_Ioc (m : ℕ) (hm : 2 ≤ m) :
    lengthTwoRho m ∈ Set.Ioc (0 : ℚ) (22 / 5) :=

def boundaryCharge (n : ℕ) : ℂ := (boundaryMinimalCentralCharge n : ℂ)
\end{leancode}

\begin{dictionary}
Rational numbers \lean{ℚ}; \lean{Set.Ioc a b} is the half-open interval $(a,b]$; \lean{(q : ℂ)} is the
embedding $\Q\subset\C$. Check: $1-6(2-(2n+1))^2/(2(2n+1)) = 1-3(2n-1)^2/(2n+1)$, the $(2,2n+1)$
minimal-model value ($n=2$: $-22/5$, \lean{boundaryCharge_}\allowbreak\lean{two_}\allowbreak\lean{eq_}\allowbreak\lean{leeYangC}; $n=3$: $-68/7$). For $n=1$ it
gives $c_{2,3} = 0$; all theorems below assume $n\ge2$. $\rho_2 = 36\cdot11/90 = 22/5$; the library proves
$\rho_m$ strictly decreasing (\lean{lengthTwoRho_succ_lt}) with $0<\rho_m\le22/5$ for $m\ge2$.
\end{dictionary}

\begin{physics}
$-\rho_m$ is the central charge at which the new length-two quasi-primary of weight $2m$ becomes null
(norm zero) in $V(c,0)$; $\rho_2 = 22/5$ is the Lee--Yang value.
\end{physics}

\begin{status}
Definitions and inequalities proved, axiom-clean.
\end{status}

\subsection{The length-two quasi-primaries}

\begin{definition}[$Q_{2m}$]\label{l19:def:Q}
For $m\ge2$, with $a_{m,j} = (-1)^j\binom{2m-1}{j}/\binom{j+3}{3}$ and
$b_m = (-1)^{m+1}\dfrac{3\binom{2m}{m}}{m(m+1)(2m+1)}$,
\begin{equation*}
 Q_{2m} = 2\sum_{j=0}^{m-3}a_{m,j}L_{-(2m-2-j)}L_{-(j+2)}\lnineteenvac + a_{m,m-2}L_{-m}^2\lnineteenvac +
 b_mL_{-2m}\lnineteenvac\in V(c,0)_{2m}.
\end{equation*}
\end{definition}

\noindent The rational coefficients $a_{m,j}$ and $b_m$ live in their own file,
\leanpath{VOA/\allowbreak Examples/\allowbreak VirasoroLengthTwoArithmetic.lean}:
\begin{leancode}
def lengthTwoACoeff (m j : ℕ) : ℚ :=
  (-1 : ℚ) ^ j * (Nat.choose (2 * m - 1) j : ℚ) / (Nat.choose (j + 3) 3 : ℚ)

def lengthTwoBCoeff (m : ℕ) : ℚ :=
  (-1 : ℚ) ^ (m + 1) * 3 / ((m : ℚ) * (m + 1) * (2 * m + 1)) *
    (Nat.choose (2 * m) m : ℚ)
\end{leancode}

\noindent The vector $Q_{2m}$ itself and its quasi-primarity are in the file
\leanpath{VOA/\allowbreak Examples/\allowbreak VirasoroLengthTwoQuasiPrimary.lean}, which imports the
coefficients above:
\begin{leancode}
noncomputable def lengthTwoQOffDiag (m : ℕ) : VirUniv c :=
  (2 : ℂ) • ∑ j ∈ Finset.range (m - 2),
    aCoeffC m j • negWordVir c [2 * m - 4 - j, j]

noncomputable def lengthTwoQ (m : ℕ) : VirUniv c :=
  lengthTwoQOffDiag c m
    + aCoeffC m (m - 2) • negWordVir c [m - 2, m - 2]
    + bCoeffC m • negWordVir c [2 * m - 2]

theorem lengthTwoQ_quasiPrimary (m : ℕ) (hm : 2 ≤ m) :
    Lop c 1 (lengthTwoQ c m) = 0 := by
\end{leancode}

\begin{dictionary}
\lean{aCoeffC}, \lean{bCoeffC} are the coefficients cast to $\C$ (private definitions in the file).
\lean{negWordVir c [2*m-4-j, j]} $= L_{-(2m-2-j)}L_{-(j+2)}\lnineteenvac$ by the word convention;
\lean{[2*m-2]} gives $L_{-2m}\lnineteenvac$. \lean{Lop c 1} is $L_1$ on $V(c,0)$. So the theorem says
$L_1Q_{2m} = 0$: $Q_{2m}$ is \emph{quasi-primary}. Since $L_2Q_{2m}\ne0$ in general it is not primary.
$Q_{2m}$ is independent of $c$; for $m=2$, $Q_4 = L_{-2}^2\lnineteenvac-\tfrac35L_{-4}\lnineteenvac$
(\lean{lengthTwoQ_two}), and $5Q_4$ is the Lee--Yang singular vector (\lean{five_smul_lengthTwoQ_two}).
\end{dictionary}

\paragraph{Proof on paper.} Since $[L_1,L_{-a}] = (a+1)L_{1-a}$ and $L_1\lnineteenvac = 0$,
$L_1L_{-a}L_{-b}\lnineteenvac = (a+1)L_{-(a-1)}L_{-b}\lnineteenvac + (b+1)L_{-a}L_{-(b-1)}\lnineteenvac$, where a factor
$L_{-1}$ adjacent to the vacuum is moved through using $L_{-1}L_{-b}\lnineteenvac = (b-1)L_{-(b+1)}\lnineteenvac$, and
$L_1L_{-2m}\lnineteenvac = (2m+1)L_{-(2m-1)}\lnineteenvac$. Collecting the coefficient of each word of weight $2m-1$,
the length-two words cancel in neighbouring pairs by the identity
$a_{m,j}(2m-1-j) + a_{m,j+1}(j+4) = 0$ for $j+3\le m$ (\lean{lengthTwoACoeff_recurrence}), and the single
length-one word $L_{-(2m-1)}\lnineteenvac$ cancels by $(2m+1)b_m + (m+1)a_{m,m-2} = 0$
(\lean{lengthTwoBCoeff_residual}). \qed

\begin{physics}
Quasi-primaries are the $SL(2)$ highest-weight states: the states annihilated by the special conformal
generator. In a 4d $\mathcal N=2$ SCFT they correspond to the Schur operators that are not derivatives;
$Q_{2m}$ is the ``double-trace'' $:TT:$-type operator of spin $2m$ built from two stress tensors with
$2m-4$ derivatives distributed so as to be quasi-primary.
\end{physics}

\begin{status}
Proved, axiom-clean, unconditional (no use of \lean{VOA.pbw}, no character theorem).
\end{status}

\begin{theorem}[Closed norm formula]\label{l19:thm:norm}
For every $c\in\C$ and $m\ge2$,
\begin{equation*}
 (Q_{2m}|Q_{2m}) = \tfrac12\binom{4m-2}{2m-4}\,c\,(c+\rho_m)
\end{equation*}
for the canonical form of $V(c,0)$.
\end{theorem}

\noindent File \leanpath{VOA/\allowbreak Examples/\allowbreak VirasoroLengthTwoNorm.lean}.
\begin{leancode}
theorem canB_lengthTwoQ_self (c : ℂ) (m : ℕ) (hm : 2 ≤ m) :
    (canB c).form (lengthTwoQ c m) (lengthTwoQ c m) =
      (1 / 2 : ℂ) * ((Nat.choose (4 * m - 2) (2 * m - 4) : ℚ) : ℂ) * c *
        (c + (lengthTwoRho m : ℂ)) := by
\end{leancode}

\begin{dictionary}
\lean{canB c} is the canonical normalised invariant form on $V(c,0)$ (Chapter~\ref{ch:L06},
Definition~\ref{l06:def:canB}). The binomial is cast $\N\to\Q\to\C$. Exactly Theorem~\ref{l19:thm:norm}.
Check at $m=2$: $\tfrac12\cdot1\cdot c(c+22/5) = c(5c+22)/10$, which is
\lean{canB_}\allowbreak\lean{lengthTwoQ_}\allowbreak\lean{self_}\allowbreak\lean{two_}\allowbreak\lean{norm}, and which vanishes exactly at $c=0$ and $c=-22/5$ --- the
level-four Kac zeros.
\end{dictionary}

\paragraph{Proof on paper.} Since $Q_{2m}$ is quasi-primary and of weight $2m$, its norm equals
$(\text{top pairing})\times$normalisation: write $(Q|Q) = (Q|b_mL_{-2m}\lnineteenvac + \ldots)$, move the
$L$'s to the left as $L_n$'s, and use $L_1Q=0$ to kill all terms except those pairing with the
``top'' word. The resulting finite sums are evaluated in closed form by a binomial identity
(\lean{lengthTwoNormSummand_}\allowbreak\lean{sum_}\allowbreak\lean{full}, \lean{lengthTwoChooseThree}, \lean{lengthTwoMixedArithmetic}),
yielding a quadratic polynomial in $c$ with roots $0$ and $-\rho_m$. \qed

\begin{physics}
The two-point function of the double-trace operator: proportional to $c$ (one stress tensor
contraction) and to $c+\rho_m$ (the $1/c$ correction from the connected four-point function of $T$),
vanishing at the central charge where $Q_{2m}$ decouples.
\end{physics}

\begin{status}
Proved, axiom-clean, unconditional, all $c$, all $m\ge2$.
\end{status}

\subsection{Towers, layers and the theorem}

Every state of length $\le2$ and weight $N\ge2$ in $V(c,0)$ is a combination of $L_{-1}$-descendants of
$\omega$ and of the $Q_{2m}$: by the global $\mathfrak{sl}_2 = \langle L_{-1},L_0,L_1\rangle$ structure,
the weight-$N$ part of the length-$\le2$ subspace is spanned by
$L_{-1}^{N-2}\omega$ (``stress tower'', index $m=1$) and $L_{-1}^{N-2m}Q_{2m}$, $2\le m\le N/2$
(\lean{virFil_two_eq_span_towers}). These towers are eigenvectors of the self-adjoint operator
$L_{-1}L_1$ with distinct eigenvalues $r(2h+r-1)$ ($h$ the weight of the base, $r$ the depth;
\lean{globalTowerOperator}, \lean{lengthLETwoTowerEigenvalue_}\allowbreak\lean{injective}), hence mutually orthogonal and
linearly independent, and
\begin{equation*}
 \big(L_{-1}^rQ\,\big|\,L_{-1}^rQ\big) = \prod_{j=0}^{r-1}(j+1)(2h+j)\cdot(Q|Q),
\end{equation*}
a positive multiple of $(Q|Q)$.

\begin{definition}[Retained towers in the simple quotient]\label{l19:def:towers}
For the boundary model $n$ and weight $N$, the retained indices are $1\le m\le N/2$ except $(n,m) =
(2,2)$; the corresponding towers in $L(\lnineteencb,0)$ are the images of $L_{-1}^{N-2}\omega$ ($m=1$) and
$L_{-1}^{N-2m}Q_{2m}$ ($m\ge2$).
\end{definition}

\noindent File \leanpath{VOA/\allowbreak Examples/\allowbreak VirasoroLengthTwoBoundary.lean}.
\begin{leancode}
def BoundaryLengthLETwoTowerIndex (n N : ℕ) :=
  {m : LengthLETwoTowerIndex N // n ≠ 2 ∨ m.1 ≠ 2}

noncomputable def boundaryLengthLETwoTowerAt (n N : ℕ)
    (m : BoundaryLengthLETwoTowerIndex n N) :
    VirUniv (boundaryCharge n) ⧸ maxIdealVir (boundaryCharge n) :=
  if hm : m.1.1 = 1 then boundaryLengthOneTowerAt n N
  else ((virasoroSimple (boundaryCharge n)).L (-1) ^ (N - 2 * m.1.1))
    (lengthTwoQSimple (boundaryCharge n) m.1.1)
\end{leancode}

\begin{dictionary}
\lean{LengthLETwoTowerIndex N} is \lean{{m : ℕ // 1 ≤ m ∧ 2 * m ≤ N}}
(\leanpath{VirasoroLengthTwoSpanning.lean}); the subtype removes the single Lee--Yang exception.
\lean{boundaryLengthOneTowerAt n N} is $L_{-1}^{N-2}\omega$ in the quotient and
\lean{lengthTwoQSimple c m} is the image of $Q_{2m}$ (\leanpath{VirasoroLengthTwoTowers.lean}).
\end{dictionary}

\begin{physics}
The stress-tensor descendants and the descendants of the double-trace quasi-primaries: the complete
list of states with at most two stress-tensor modes.
\end{physics}

\begin{status}
Definitions. The library proves these retained towers are nonzero, orthogonal, a basis of
$\lnineteenG_2L(\lnineteencb,0)_N$ (\lean{goodFil_}\allowbreak\lean{two_}\allowbreak\lean{eq_}\allowbreak\lean{span_}\allowbreak\lean{boundaryTowers}), and that the deleted Lee--Yang
tower is zero (\lean{lengthTwoQSimple_}\allowbreak\lean{two_}\allowbreak\lean{eq_}\allowbreak\lean{zero}).
\end{status}

\begin{theorem}[Exact signs of the retained towers]\label{l19:thm:towersign}
For $n\ge2$ and every retained index: the stress tower has norm $-r$ with $r>0$, and every length-two
tower has norm $+r$ with $r>0$.
\end{theorem}

\noindent File \leanpath{VOA/\allowbreak Examples/\allowbreak VirasoroLengthTwoBoundary.lean}.
\begin{leancode}
theorem Bsimp_boundaryLengthLETwoTowerAt_self_sign (n N : ℕ) (hn : 2 ≤ n)
    (m : BoundaryLengthLETwoTowerIndex n N) :
    if hm : m.1.1 = 1 then
      ∃ r : ℝ, 0 < r ∧
        (Bsimp (boundaryCharge n)).form (boundaryLengthLETwoTowerAt n N m)
          (boundaryLengthLETwoTowerAt n N m) = -(r : ℂ)
    else
      ∃ r : ℝ, 0 < r ∧
        (Bsimp (boundaryCharge n)).form (boundaryLengthLETwoTowerAt n N m)
          (boundaryLengthLETwoTowerAt n N m) = (r : ℂ) := by
\end{leancode}

\begin{dictionary}
\lean{if hm : m.1.1 = 1 then … else …} is a dependent case split on whether the index is the stress
tower. Exactly Theorem~\ref{l19:thm:towersign}. The norms are of the \emph{bilinear} form on the
quotient; since every real structure fixes the towers (\lean{fixes_}\allowbreak\lean{boundaryLengthLETwoTowerAt}), they are
also the Hermitian norms.
\end{dictionary}

\paragraph{Proof on paper.} Stress tower: positive factor times $(\omega|\omega) = c/2<0$, since
$\lnineteencb<0$ for $n\ge2$ (\lean{boundaryMinimalCentralCharge_}\allowbreak\lean{neg}). Length-two tower: positive factor
times $\tfrac12\binom{4m-2}{2m-4}c(c+\rho_m)$, a product of two negative numbers: $c<0$, and
$c+\rho_m<0$ because $c\le-22/5$ with equality only at $n=2$, and $\rho_m\le22/5$ with equality only at
$m=2$ (\lean{boundaryLengthTwoBaseNorm_}\allowbreak\lean{pos}, hypothesis \lean{n ≠ 2 ∨ m ≠ 2}). \qed

\begin{physics}
In $M(2,2n+1)$ every stress-tensor descendant has negative norm and every double-trace descendant has
positive norm: exactly the $(-1)^R$ pattern of graded unitarity with $R = 1$ and $R=2$.
\end{physics}

\begin{status}
Proved, axiom-clean, unconditional.
\end{status}

\begin{definition}[Graded unitarity through degree two]\label{l19:def:throughtwo}
Clauses (iv) and (vi) restricted to filtration degrees $p\le2$.
\end{definition}

\noindent File \leanpath{VOA/\allowbreak Examples/\allowbreak VirasoroLengthTwoLayerPositive.lean}.
\begin{leancode}
def LayerNondegenerateThroughTwo (F : GradedFlag (virasoroSimple c)) (N : ℤ)
    (β : LinearMap.BilinForm ℂ ((virasoroSimple c).grading N)) : Prop :=
  ∀ p : ℕ, p ≤ 2 → (β.restrict (F.pieceIn p N)).Nondegenerate

def LayerPositiveThroughTwo (F : GradedFlag (virasoroSimple c)) (N : ℤ)
    (β : LinearMap.BilinForm ℂ ((virasoroSimple c).grading N))
    (h : SesqForm ((virasoroSimple c).grading N)) : Prop :=
  ∀ p : ℕ, p ≤ 2 → IsPosOn (signedForm h p) (F.layer N β p)

def GradedUnitaryThroughTwoAt (F : GradedFlag (virasoroSimple c)) (N : ℤ)
    (β : LinearMap.BilinForm ℂ ((virasoroSimple c).grading N))
    (h : SesqForm ((virasoroSimple c).grading N)) : Prop :=
  LayerNondegenerateThroughTwo F N β ∧ LayerPositiveThroughTwo F N β h
\end{leancode}

\begin{dictionary}
Literally \lean{LayerNondegenerate} and \lean{LayerPositive} with the extra guard \lean{p ≤ 2}. So
\lean{GradedUnitaryAt} implies \lean{GradedUnitaryThroughTwoAt}, not conversely.
\end{dictionary}

\begin{physics}
The positivity constraints on states with $R\le2$: operators made of at most two stress tensors.
\end{physics}

\begin{status}
Definitions.
\end{status}

\begin{theorem}[Boundary Virasoro graded unitarity through filtration degree two]\label{l19:thm:throughtwo}
For every $n\ge2$ and every $N\in\Z$, the simple boundary model $L(\lnineteencb,0)$ with its canonical
form, its length flag and the canonical real structure satisfies clauses (iv) and (vi) of graded
unitarity at weight $N$ in all filtration degrees $p\le2$. The same holds for an \emph{arbitrary} real
structure.
\end{theorem}

\noindent File \leanpath{VOA/\allowbreak Examples/\allowbreak VirasoroLengthTwoLayerPositive.lean}.
\begin{leancode}
theorem boundary_gradedUnitaryThroughTwoAt (n : ℕ) (hn : 2 ≤ n) (N : ℤ)
    (R : RealStructure (virasoroSimple (boundaryCharge n))) :
    GradedUnitaryThroughTwoAt (virGradedFlag (boundaryCharge n)) N
      (LiRadical.gradedForm (Bsimp (boundaryCharge n)) N)
      (gradedSesq (Bsimp (boundaryCharge n)) R N) := by

theorem boundary_gradedUnitaryThroughTwo (n : ℕ) (hn : 2 ≤ n) (N : ℤ) :
    GradedUnitaryThroughTwoAt (virGradedFlag (boundaryCharge n)) N
      (LiRadical.gradedForm (Bsimp (boundaryCharge n)) N)
      (gradedSesq (Bsimp (boundaryCharge n)) (boundaryVirasoroRealStructure n) N) :=
\end{leancode}

\begin{dictionary}
Every ingredient has been defined above: \lean{boundaryCharge n} (Def.~\ref{l19:def:boundarycharge}),
\lean{virGradedFlag} (Def.~\ref{l19:def:lengthflag}), \lean{Bsimp} (canonical form),
\lean{gradedSesq} (Def.~\ref{l19:def:gradedsesq}), \lean{boundaryVirasoroRealStructure}
(Thm.~\ref{l19:thm:boundaryreal}). The only hypothesis is $n\ge2$ (excluding the trivial $c_{2,3} = 0$).
The first form holds for an arbitrary real structure \lean{R} --- this is \emph{stronger}, and it is
possible because every real structure fixes all towers (Theorem~\ref{l19:thm:rhostable} and its proof).
The second form instantiates the canonical one and is therefore hypothesis-free.

\textbf{Same, stronger or weaker than the target?} Strictly \emph{weaker} than
\lean{VirGradedUnitary (boundaryCharge n) R}: layers $p\ge3$ are not addressed. For $n=2$ and weights
$N\le 11$ it nevertheless covers everything, since Lee--Yang has no admissible length-three partitions
below weight $12$ (\lean{a3count_}\allowbreak\lean{vanishes_}\allowbreak\lean{below_}\allowbreak\lean{twelve}, Section~\ref{l19:sec:programb}) --- but the library
does not draw this corollary, because identifying admissible partitions with a basis is the uncited
\lean{FlagBasisHyp}.

\textbf{Non-vacuity.} The conclusion would be empty if the layers were zero. A Lean check against the
library (not part of this release) proves that for all $n\ge2$, $N\ge2$ the degree-one layer is a nonzero line,
that the degree-two layer is nonzero whenever a length-two tower is retained,
and the concrete instance $M(2,7)$, weight $4$, of the latter. All three
compile with standard axioms only.
\end{dictionary}

\paragraph{Proof on paper.} Fix $n\ge2$. Negative $N$ and $N=1$: the weight space is zero and every clause
holds trivially. $N=0$: the weight space is $\C\lnineteenvac$, all steps equal it, the form is nondegenerate
since $(\lnineteenvac|\lnineteenvac) = 1$, layer $0$ is positive (Theorem~\ref{l19:thm:vacuum}) and layers $1,2$ are
zero. $N\ge2$: $\lnineteenG_0 = 0$; $\lnineteenG_1L_N$ is the line of the stress tower
(\lean{pieceIn_one_eq_span_stress}), so $\lnineteenLay_1$ is that line; $\lnineteenG_2L_N$ is spanned by the
retained towers, an orthogonal family with nonzero norms, so $\beta$ is nondegenerate on $\lnineteenG_1$ and on
$\lnineteenG_2$ (\lean{restrict_}\allowbreak\lean{span_}\allowbreak\lean{nondegenerate_}\allowbreak\lean{of_}\allowbreak\lean{orthogonal_}\allowbreak\lean{family}) and $\lnineteenLay_2$ is the span of the
length-two towers (\lean{layer_two_eq_lengthTwoSpan}). For $x = \sum_iz_it_i$ in the span of an orthogonal
$\rho$-fixed family $(t_i)$ one has $h(x,x) = \sum_i|z_i|^2(t_i|t_i)$; by Theorem~\ref{l19:thm:towersign},
$-h$ is positive on $\lnineteenLay_1$ and $+h$ on $\lnineteenLay_2$ (\lean{isPosOn_}\allowbreak\lean{span_}\allowbreak\lean{of_}\allowbreak\lean{orthogonal_}\allowbreak\lean{family}). \qed

\begin{physics}
In every boundary minimal model, the two lowest rungs of the four-dimensional positivity ladder hold
exactly: single stress-tensor states have negative two-dimensional norm, and the exact-length-two
double-trace states have positive norm, at every spin. This is the first all-weight piece of evidence for
the Kaufmann/ABLR conjecture that is a theorem rather than a numerical check.
\end{physics}

\begin{status}
Proved; \texttt{\#print axioms} of both: \lean{[propext, Classical.choice, Quot.sound]}. Unconditional
(\lean{boundary_}\allowbreak\lean{gradedUnitaryThroughTwo} has only $n\ge2$); \lean{VOA.pbw} not used. The full conjecture
(all $p$) remains open; the repository's only other evidence for it is numerical (about $700$ $(n,N)$ pairs of
the alternating-sum signature law, recorded in the project's internal records; \emph{evidence, not proof}).
\end{status}

\begin{example}[Lee--Yang at weight six, revisited]\label{l19:ex:ly6b}
At $n=2$, $N=6$ the retained towers are $m=1$ ($L_{-1}^4\omega$) and $m=3$ ($Q_6$); $m=2$ is the deleted
$Q_4$ (whose tower $L_{-1}^2Q_4$ vanishes in the quotient). $\rho_3 = 36\cdot29/2520 = 29/70$, so
$(Q_6|Q_6) = \tfrac12\binom{10}{2}\cdot(-\tfrac{22}5)(-\tfrac{22}5+\tfrac{29}{70}) = \tfrac{45}2\cdot
\tfrac{22}{5}\cdot\tfrac{279}{70}>0$, in agreement with the positive pivot $682/35$ of
Example~\ref{l19:ex:ly6}. The two examples are two views of one computation: Example~\ref{l19:ex:ly6} reads
the signs off the Gram-matrix minors (the minors criterion of Section~\ref{l19:sec:minors}), while this example
reads them off the norms of the tower vectors that span the layers; their agreement is the layer decomposition
at work.
\end{example}

\begin{warning}[title={Stale status lines in the Virasoro docstrings}]\label{l19:warn:stale}
In the earlier state of the library against which this chapter was first written,
\leanpath{VirasoroGradedUnitarity.lean} said ``No \lean{Core.RealStructure} is constructed on
$L(c,0)$'' and ``No positivity statement is proved of anything''. Both were already out of date then: real
structures are constructed (Theorem~\ref{l19:thm:virreal}) and the through-two
positivity theorem is proved (Theorem~\ref{l19:thm:throughtwo}). The correct, still-true part is that the
full \lean{VirGradedUnitary} is never asserted. At a pre-release build commit the file header no longer contains these two
sentences; it states that nothing in the repository proves graded unitarity of $L(c,0)$ and points to the
through-two truncation \lean{boundary_gradedUnitaryThroughTwo}.

The book does not repair docstrings: editing a \texttt{.lean} file (even a comment) changes the file's blob hash
recorded in \texttt{manifest.json}, so docstring repairs are left to the library's maintainers. A reader of the
Lean source should therefore trust the theorem statements quoted in this chapter, not the prose status lines in
the file headers. This is a documentation-only defect, listed with the other stale docstrings in
Section~\ref{l19:sec:defects}.
\end{warning}

%======================================================================
\section{The Lee--Yang model}\label{l19:sec:leeyang}
%======================================================================

The Lee--Yang model is the boundary model $n=2$, $M(2,5)$, $c = -22/5$. It is the testing ground
of the whole graded-unitarity programme and the library proves a number of facts about it that are
not about positivity: the derivation of its central charge from a singular vector, the level-four
structure of its maximal ideal, its $C_2$-cofiniteness and strong finiteness, and the Zhu-algebra
relation that pins its two lowest weights. Rigidity (every graded first-order deformation of
$L(-22/5,0)$ is trivial) is Chapter~\ref{ch:L07}, \S\ref{l07:sec:leeyang}.

\begin{physics}
The Yang--Lee edge singularity: the Ising model in an \emph{imaginary} magnetic field $ih$ has
partition-function zeros (Lee--Yang zeros) on the imaginary $h$-axis; at the edge of their
distribution the model is critical, and Cardy (1985) identified the critical theory with the
$M(2,5)$ minimal model, a non-unitary CFT with one nontrivial primary $\varphi$ of weight
$h = -1/5$ (scaling dimension $-2/5$, so $\langle\varphi\varphi\rangle$ grows with distance --- the
hallmark of non-unitarity). In four dimensions, $L(-22/5,0)$ is the chiral algebra of the simplest
Argyres--Douglas theory $(A_1,A_2)$, which \emph{is} unitary; graded unitarity is the reconciliation
of these two facts.
\end{physics}

\subsection{The singular vector and the central charge}

\begin{theorem}[The Lee--Yang central charge is forced]\label{l19:thm:lysing}
Let $c\in\C$. There is a nonzero vector $\alpha L_{-2}^2\lnineteenvac + \beta L_{-4}\lnineteenvac\in V(c,0)_4$ killed
by $L_1$ and $L_2$ if and only if $c = -22/5$. At $c=-22/5$ one such vector is
$s = 5L_{-2}^2\lnineteenvac - 3L_{-4}\lnineteenvac$.
\end{theorem}

\noindent File \leanpath{VOA/\allowbreak Examples/\allowbreak LeeYangSingular.lean}.
\begin{leancode}
noncomputable def leeYangC : ℂ := -22 / 5

noncomputable def leeYangSingular : VirUniv leeYangC :=
  (5 : ℂ) • L2L2vac leeYangC + (-3 : ℂ) • L4vac leeYangC

theorem exists_level4_singular_iff (c : ℂ) :
    (∃ α β : ℂ, α • L2L2vac c + β • L4vac c ≠ 0
        ∧ Lop c 1 (α • L2L2vac c + β • L4vac c) = 0
        ∧ Lop c 2 (α • L2L2vac c + β • L4vac c) = 0)
      ↔ c = -22 / 5 := by
\end{leancode}

\begin{dictionary}
\lean{L2L2vac c} $= L_{-2}L_{-2}\lnineteenvac$ and \lean{L4vac c} $= L_{-4}\lnineteenvac$ in $V(c,0)$
(defined in the same file); \lean{Lop c n} is the Virasoro mode $L_n$. Since $V(c,0)_4$ is spanned by
these two vectors (\lean{gradingVir_four}), the left side says ``$V(c,0)_4$ contains a nonzero singular
vector''. Exactly Theorem~\ref{l19:thm:lysing}. (A vector killed by $L_1,L_2$ is killed by all $L_n$,
$n>0$: \lean{Lop_pos_ann}.)
\end{dictionary}

\paragraph{Proof on paper.} $L_1L_{-2}^2\lnineteenvac = 3L_{-3}\lnineteenvac$, $L_1L_{-4}\lnineteenvac = 5L_{-3}\lnineteenvac$,
$L_2L_{-2}^2\lnineteenvac = (c+8)L_{-2}\lnineteenvac$, $L_2L_{-4}\lnineteenvac = 6L_{-2}\lnineteenvac$. The linear system
$3\alpha+5\beta = 0$, $(c+8)\alpha+6\beta = 0$ has a nonzero solution iff $18-5(c+8) = -(5c+22) = 0$;
the solution is $(5,-3)$. \qed

\begin{physics}
The central charge of the Yang--Lee CFT is not an input: it is the unique value at which the
stress-tensor OPE closes with a null state at level four, the relation $:TT: \propto \partial^2T$ that
truncates the vacuum module. Cardy derived $c = -22/5$ from exactly this null vector.
\end{physics}

\begin{status}
Proved, axiom-clean, unconditional.
\end{status}

\begin{theorem}[Level four of the maximal ideal; non-simplicity of $V(-22/5,0)$]\label{l19:thm:lylevel4}
The weight-four part of the maximal proper ideal of $V(-22/5,0)$ is the line $\C s$; in particular
$V(-22/5,0)$ is not simple.
\end{theorem}

\noindent Files \leanpath{VOA/\allowbreak Examples/\allowbreak LeeYangKac.lean} and
\leanpath{VOA/\allowbreak Examples/\allowbreak LeeYangProper.lean}.
\begin{leancode}
theorem maxIdealVir_leeYang_inf_grading_four :
    maxIdealVir leeYangC ⊓ gradingVir leeYangC 4 = Submodule.span ℂ {leeYangSingular} := by

theorem virasoroVOA_leeYang_not_isSimple : ¬ (virasoroVOA leeYangC).IsSimple :=
  not_isSimple_of_idealGen_proper leeYangSingular_ne_zero vacuum_notMem_idealGen
\end{leancode}

\begin{dictionary}
\lean{gradingVir c 4} is the weight-four subspace of $V(c,0)$ (\leanpath{VirasoroGrading.lean});
\lean{maxIdealVir c} the maximal proper graded ideal; \lean{(virasoroVOA c).IsSimple} the VOA simplicity
predicate (Chapter~\ref{ch:L04}). The second theorem's one-line proof term is informative: the ideal
generated by $s$ is proper because it does not contain the vacuum (all its states have weight
$\ge4$, \lean{four_le_leeYangEnergy}). The docstring of the first warns that it is a level-four
statement only: the full description of $\mathcal M_{-22/5}$ (the ideal generated by $s$, which is the
Rogers--Ramanujan character identity) is \textbf{not} proved.
\end{dictionary}

\paragraph{Proof on paper.} A vector of a proper ideal pairs to zero with everything under the
invariant form (the form is $\lnineteenvac$-valued and the ideal misses $\lnineteenvac$), so
$\mathcal M\cap V_4$ lies in the radical of the level-four Gram form; the library computes that radical to
be $\C s$ (\lean{radBil_}\allowbreak\lean{shapForm_}\allowbreak\lean{leeYang_}\allowbreak\lean{eq_}\allowbreak\lean{span}) and shows $s\in\mathcal M$. \qed

\begin{physics}
The only null state at level four is $s$; in the Lee--Yang CFT the vacuum sector at weight four is
one-dimensional ($\partial^2T$ only).
\end{physics}

\begin{status}
Both proved, axiom-clean, unconditional.
\end{status}

\begin{warning}[title={Two different level-four forms}]\label{l19:warn:shap}
\leanpath{LeeYangKac.lean} computes the Gram matrix of the ``top-mode'' (Shapovalov) form
$\langle a,b\rangle_4 = \varepsilon(a_{(7)}b)$ of \lean{RigidBridge.shapForm} in the basis
$(L_{-2}^2\lnineteenvac,L_{-4}\lnineteenvac)$:
\lean{gram4_eq : gram4 c hvac = !![c ^ 2 / 2 + 40 * c, 63 * c; 63 * c, 105 * c]}, with
\lean{det_gram4 : (gram4 c hvac).det = 21 * (c ^ 2 * (5 * c + 22) / 2)}. This matrix is \emph{not} $21$
times the textbook contravariant Kac matrix $\begin{psmallmatrix}\frac c2(c+8)&3c\\3c&5c\end{psmallmatrix}$
(the $(1,1)$ entries differ by $-2c(5c+22)$), although the determinants agree up to the factor $21$. The
reason is that $L_{-2}^2\lnineteenvac$ is not quasi-primary, and $a_{(2k-1)}b$ equals the invariant pairing (up to
$(-1)^k$) only for quasi-primary $a$. The file documents this correctly. The docstring of
\lean{shapForm} in \leanpath{VOA/\allowbreak Core/\allowbreak RigidityBridge.lean} (Chapter~\ref{ch:L07}), which says
``This is the form whose Gram matrix is the level-$k$ Kac matrix'', is \textbf{inaccurate} in general
(documentation defect).
\end{warning}

\subsection{\texorpdfstring{$C_2$}{C2}-cofiniteness and strong finiteness}

\begin{theorem}[Lee--Yang is strongly finite]\label{l19:thm:lysf}
$L(-22/5,0)$ is $C_2$-cofinite, with $\dim L(-22/5,0)/C_2 \le 2$, and it is strongly finite (simple,
self-dual and $C_2$-cofinite). If it is moreover admissibly rational (or regular), it is strongly rational.
\end{theorem}

\noindent File \leanpath{VOA/\allowbreak Examples/\allowbreak LeeYangC2.lean}.
\begin{leancode}
theorem leeYang_c2Cofinite : C2Cofinite (virasoroSimple leeYangC) :=

theorem finrank_a₂_leeYang_le_two :
    Module.finrank ℂ (A₂ (virasoroSimple leeYangC)) ≤ 2 := by

theorem leeYang_stronglyFinite : StronglyFinite (virasoroSimple leeYangC) where
  simple := virasoroSimple_isSimple leeYangC
  selfDual := virasoroSimple_selfDual leeYangC
  c2Cofinite := leeYang_c2Cofinite

theorem leeYang_stronglyRational_of_admissiblyRational
    (h : AdmissiblyRational (virasoroSimple leeYangC)) :
    StronglyRational (virasoroSimple leeYangC) where
\end{leancode}

\begin{dictionary}
\lean{C2Cofinite W} is \lean{FiniteDimensional ℂ (A₂ W)}, where \lean{A₂ W} is the quotient
$V/C_2(V)$ (\leanpath{VOA/\allowbreak Modules/\allowbreak C2Cofinite.lean}; the name \lean{A₂} is \emph{not} the second higher
Zhu algebra, cf.\ Chapter~\ref{ch:L12}). \lean{StronglyFinite} is the structure of
Chapter~\ref{ch:L12} (Definition~\ref{l12:def:strong}): simple, self-dual, $C_2$-cofinite; the three fields
are supplied explicitly in the quoted proof term, so the reader can see that simplicity and self-duality are the
general Virasoro facts and only $C_2$-cofiniteness is specific. The last theorem has hypothesis
\lean{AdmissiblyRational}, which is \textbf{not} proved for Lee--Yang: rationality of $L(c_{p,q},0)$ (Wang 1993)
is not in the library. Exactly Theorem~\ref{l19:thm:lysf}.

\textbf{Non-vacuity.} \lean{leeYang_}\allowbreak\lean{conformalVector_}\allowbreak\lean{ne_}\allowbreak\lean{zero}: the VOA is not the trivial one.
\end{dictionary}

\paragraph{Proof on paper.} $L(c,0)$ is strongly generated by $\omega$ (every state is a combination of
iterated products of $\omega$, \lean{virasoroSimple_}\allowbreak\lean{iterProdSpan_}\allowbreak\lean{conformalVector_}\allowbreak\lean{eq_}\allowbreak\lean{top}). In
$R_V = V/C_2(V)$, $L_{-4}\lnineteenvac = \omega_{(-3)}\lnineteenvac\in C_2(V)$, so the relation $5L_{-2}^2\lnineteenvac -
3L_{-4}\lnineteenvac = 0$ becomes $5[\omega]^2 = 0$. A commutative algebra generated by one nilpotent element
$x$ with $x^2 = 0$ is spanned by $1,x$; the library's criterion \lean{c2Cofinite_}\allowbreak\lean{of_}\allowbreak\lean{iterProdSpan_}\allowbreak\lean{of_}\allowbreak\lean{isNilpotent}
(from the triplet files) turns this into $\dim R_V\le2$. \qed

\begin{physics}
$R_V = \C[x]/(x^2)$ is the coordinate ring of the \emph{associated variety}, here a fat point:
the Lee--Yang theory has no Higgs branch, as expected for the $(A_1,A_2)$ Argyres--Douglas theory, and
$C_2$-cofiniteness is the finiteness that underlies modularity of its two characters (the
Rogers--Ramanujan functions).
\end{physics}

\begin{status}
\lean{leeYang_c2Cofinite}, \lean{finrank_a₂_leeYang_le_two}, \lean{leeYang_stronglyFinite}: proved,
axiom-clean, unconditional. The \lean{StronglyRational} corollaries are conditional on
\lean{AdmissiblyRational} or \lean{Regular} (open in the library).
\end{status}

\subsection{The Zhu-algebra relation}

\begin{theorem}[Lee--Yang lowest weights]\label{l19:thm:lyzhu}
In Zhu's algebra $A(L(-22/5,0))$ one has $5[\omega]^2+[\omega] = 0$. Consequently, on any nonzero
$A(V)$-module on which $[\omega]$ acts by a scalar $h$, $h\in\{0,-1/5\}$; every lowest-weight module has
lowest weight $0$ or $-1/5$.
\end{theorem}

\noindent File \leanpath{VOA/\allowbreak Examples/\allowbreak LeeYangZhu.lean}.
\begin{leancode}
theorem leeYang_zhu_polynomial :
    (5 : ℂ) • Amul (zhuRingData_holds WLY) omegaA omegaA + omegaA = 0 := by

theorem omegaScalar_eq_zero_or_neg_one_fifth (U : AMod WLY) [Nontrivial U.carrier]
    {h : ℂ} (hU : U.OmegaScalar h) : h = 0 ∨ h = -(1 / 5 : ℂ) := by

theorem lowestWeight_eq_zero_or_neg_one_fifth (P : LWModule WLY) :
    P.h = 0 ∨ P.h = -(1 / 5 : ℂ) := by
\end{leancode}

\begin{dictionary}
\lean{WLY} abbreviates \lean{virasoroSimple leeYangC}; \lean{omegaA} is the class $[\omega]$ in Zhu's
algebra \lean{A WLY}, \lean{Amul} its product (Chapter~\ref{ch:L12}, \S\ref{l12:sec:zhu}).
\lean{AMod} is a bundled $A(V)$-module, \lean{OmegaScalar h} says $[\omega]$ acts as $h\cdot\mathrm{id}$,
\lean{LWModule} a lowest-weight module (Chapter~\ref{ch:L12}). Exactly Theorem~\ref{l19:thm:lyzhu}. Only
the \emph{necessary} condition is proved: that both weights $0$ and $-1/5$ are realised (and that these are all
irreducibles) is not.
\end{dictionary}

\paragraph{Proof on paper.} In Zhu's algebra $[(L_{-1}+L_0)a] = 0$, so $[L_{-1}a] = -\wt(a)[a]$; hence
$[L_{-3}\lnineteenvac] = [L_{-1}\omega] = -2[\omega]$ and $[L_{-4}\lnineteenvac] = \tfrac12[L_{-1}^2\omega] = 3[\omega]$
(\lean{mk_mkQ_L4vac}); and $[L_{-2}^2\lnineteenvac] = [\omega_{(-1)}\omega] = [\omega]*[\omega]+2[\omega]$
(\lean{mk_mkQ_L2L2vac}, using $a*b = \sum_j\binom{\wt a}{j}a_{(j-1)}b$). So
$0 = [s] = 5([\omega]^2+2[\omega]) - 9[\omega] = 5[\omega]^2+[\omega]$, and on a module where $[\omega] = h$,
$h(5h+1) = 0$. \qed

\begin{physics}
The two primaries of the Yang--Lee CFT: the identity ($h=0$) and $\varphi$ ($h = -1/5$). The negative
weight is non-unitarity made visible.
\end{physics}

\begin{status}
Proved, axiom-clean, unconditional. Also \lean{admissibleBottomScalar_}\allowbreak\lean{eq_}\allowbreak\lean{zero_}\allowbreak\lean{or_}\allowbreak\lean{neg_}\allowbreak\lean{one_}\allowbreak\lean{fifth} for
the bottom of any admissible module (with the extra hypotheses that the bottom is nonzero and $L_0$ acts
by a scalar there, which the docstring states are not part of \lean{AdmissibleModule}).
\end{status}

\subsection{Other Lee--Yang files}

\leanpath{LeeYangGram.lean} computes level-four Gram entries at $c=-22/5$; \leanpath{LeeYangRigidity.lean}
instantiates the ``rigidity bridge'' of Chapter~\ref{ch:L07} at Lee--Yang, but only under an explicit
undischarged hypothesis \lean{hself} (\lean{leeYangSingular_}\allowbreak\lean{not_}\allowbreak\lean{liftable}, \lean{leeYang_F2_lt_F1}); the
unconditional rigidity theorem is \lean{VirDeformModes.}\allowbreak\lean{leeYang_}\allowbreak\lean{graded_}\allowbreak\lean{rigidity} of Chapter~\ref{ch:L07}.

\begin{theorem}[Abstract Riccati tower; finite-dimensional core of a proposed length induction]\label{l19:thm:riccati}
Let $(Q_N,S_N,K_N)_{N\in\N}$ be real square matrices and $a_N>0$, with lowering matrices $A_{N,i}$, such that
for all $N$: $Q_N$ is positive definite, $K_N$ is positive semidefinite, $S_N = Q_N^{-1}K_NQ_N^{-1}$,
$S_NQ_NS_N - a_NS_N = \sum_{i<N}A_{N,i}^TS_iA_{N,i}$, and either $S_N$ is invertible or the $A_{N,i}$ have
trivial joint kernel. Then every $S_N$ is positive definite and $S_N - a_NQ_N^{-1}$ is positive semidefinite.
\end{theorem}

\noindent File \leanpath{VOA/\allowbreak Examples/\allowbreak LeeYangRiccatiReduction.lean}.
\begin{leancode}
theorem tower_posDef (Q S K : ∀ N, Matrix (ι N) (ι N) ℝ) (a : ℕ → ℝ)
    (A : ∀ N, (i : Fin N) → Matrix (ι i) (ι N) ℝ)
    (hQ : ∀ N, (Q N).PosDef) (hK : ∀ N, (K N).PosSemidef)
    (hrec : ∀ N, S N = (Q N)⁻¹ * K N * (Q N)⁻¹)
    (hric : ∀ N, S N * Q N * S N - a N • S N = towerSource S A N)
    (ha : ∀ N, 0 < a N)
    (hker : ∀ N, IsUnit (S N) ∨ ∀ v : ι N → ℝ, (∀ i : Fin N, A N i *ᵥ v = 0) → v = 0) :
    ∀ N, (S N).PosDef ∧ (S N - a N • (Q N)⁻¹).PosSemidef := by
\end{leancode}

\begin{dictionary}
\lean{towerSource S A N} $= \sum_{i<N}A_{N,i}^TS_iA_{N,i}$; \lean{*ᵥ} is matrix--vector multiplication.
Exactly Theorem~\ref{l19:thm:riccati}. This is pure linear algebra: \textbf{no Lee--Yang matrix is shown to
satisfy any hypothesis}. The file docstring says so (``the uniform inequality $K_N\ge0$ is an open
hypothesis, not a theorem, and this file must not be cited as a proof of graded unitarity''). Note also that
the hypotheses \lean{hK} and \lean{hrec} already give $S_N\ge0$; the content is the upgrade to definiteness
and the Riccati margin.
\end{dictionary}

\paragraph{Proof on paper.} $S_N = Q_N^{-1}K_NQ_N^{-1}\ge0$. By strong induction the right side of the
Riccati identity is $\ge0$, and $>0$ under the joint-kernel condition; a vector in the radical of $S_N$ is
then killed by the right side, contradiction. The margin: in the $Q$-whitened frame the eigenvalues
$\lambda$ of $S$ satisfy $\lambda(\lambda+a)\ge0$ with $\lambda>0$, forcing $\lambda\ge$ the claimed bound. \qed

\begin{physics}
A proposed proof strategy for all-length graded unitarity of Lee--Yang: positivity at length $r+1$ is
reconstructed from length $r$ through a quadratic (Riccati) identity coming from the null vector. Only the
matrix algebra is formalised.
\end{physics}

\begin{status}
Proved, axiom-clean. Its hypotheses are not instantiated anywhere; it carries no VOA content.
\end{status}

%======================================================================
\section{Program B: finite certificates around the all-length problem}\label{l19:sec:programb}
%======================================================================

Six files record exact finite computations and interfaces produced by the project's attempt on the
full Kaufmann conjecture (``Program B''). None of them proves a statement about a VOA beyond finite
arithmetic; all are axiom-clean (some use \lean{decide}, i.e.\ kernel evaluation, never
\lean{native_decide}). We list them with their headline and their honest scope.

\begin{theorem}[Length-three counting and negative definiteness, weights $12$--$24$]\label{l19:thm:len3}
(a) For $0\le k<33$, the number of Lee--Yang admissible length-three partitions of $k+12$ equals the number of
partitions of $k$ into at most three parts; there are none below weight $12$. (b) The explicit rational
``new-tower'' blocks $N_M$ at Lee--Yang weights $M = 12,\dots,24$ are negative definite.
\end{theorem}

\noindent File \leanpath{VOA/\allowbreak Examples/\allowbreak LengthThreeQPFrame.lean}.
\begin{leancode}
theorem a3count_eq_pLe3_shift :
    ∀ k ∈ List.range 33, a3count (k + 12) = pLe3 k := by decide

theorem negDef_w12 : ∀ x : ℚ, x ≠ 0 → n12_00 * x ^ 2 < 0 :=
\end{leancode}

\begin{dictionary}
\lean{a3count n} counts triples $\ell_1\ge\ell_2+2\ge\ell_3+4$, $\ell_3\ge2$ with sum $n$ (the
Andrews--Gordon gap condition of $M(2,5)$; the Lean encoding enumerates $(\ell_2,\ell_3)$ with
$\ell_1 = n-\ell_2-\ell_3$); \lean{pLe3 k} counts partitions into $\le3$ parts. \lean{n12_00} is a rational
\emph{datum} (an output of an external Python engine), not derived in Lean from the Shapovalov form. The
docstring is explicit: ``the Lean file proves their internal identities and sign structure, not their
provenance''. That these counts are \emph{dimensions} rests on the admissible-basis theorem, external.
\end{dictionary}

\begin{physics}
Evidence for the length-three rung of graded unitarity in Lee--Yang: the exact-length-three layers at
weights $12$--$24$ have the predicted sign $(-1)^3$, assuming the external matrices are correct.
\end{physics}

\begin{status}
Proved by \lean{decide}/\lean{norm_num}, axiom-clean. \textbf{Numerical evidence encoded as data}, not a
theorem about $L(-22/5,0)$.
\end{status}

\begin{theorem}[Weil parity at cutoff three]\label{l19:thm:weil}
For the thirteen balanced-tree channels of the enriched Lee--Yang fusion ring at cutoff three, the channel
polynomials sum to $uv^7+5u^2v^6+6u^3v^5+u^4v^4$, the recorded channel signs equal the Weil signs
$i^{p-q}$ of their bidegrees, and every class-uniform vacuum-anchored landing induces the ``option (a)''
twist.
\end{theorem}

\noindent File \leanpath{VOA/\allowbreak Examples/\allowbreak WeilPolarizedChannelParity.lean}.
\begin{leancode}
theorem weil_polarized_twist_option_a :
    toMv (sumP (channels.map channelPoly))
        = X 0 * X 1 ^ 7 + C 5 * X 0 ^ 2 * X 1 ^ 6 + C 6 * X 0 ^ 3 * X 1 ^ 5
            + X 0 ^ 4 * X 1 ^ 4
      ∧ (∀ s : Sector,
          Xsign s = weil (sectorBidegree s) ∧ kappaSign s = weil (sectorBidegree s))
      ∧ (∀ L : Landing, IsUniformLanding L → (L.k0 : ℕ) % 2 = 0 →
          ∀ s : Sector, twist L s = optionA s)
      ∧ ∀ L : Landing, IsUniformLanding L → ¬ ∀ s : Sector, twist L s = optionC s :=
\end{leancode}

\begin{dictionary}
\lean{toMv} interprets a list-encoded polynomial in Mathlib's \lean{MvPolynomial (Fin 2) ℤ} ($X_0 = u$,
$X_1 = v$). \lean{Xsign}, \lean{kappaSign}, \lean{optionA}, \lean{optionC} and the fusion weights \lean{w} are
\textbf{external data} (Godfard's enriched fusion ring, Dotsenko--Fateev channel norms), stated as
definitions. What is proved is the finite combinatorics connecting them. The docstring's ``Not claimed''
list (no Hodge filtration on conformal blocks, no analytic magnitudes) is accurate.
\end{dictionary}

\begin{physics}
A heuristic Hodge-theoretic explanation of the graded-unitarity signs in terms of conformal blocks of the
Lee--Yang fusion category; exploratory.
\end{physics}

\begin{status}
Proved by \lean{decide}, axiom-clean. Combinatorics on external data; no VOA content.
\end{status}

\begin{theorem}[Minor signs from a congruence normal form]\label{l19:thm:payload}
Let $S$ be a real symmetric $n\times n$ matrix, $d:\{1..n\}\to\N$, $U$ upper unitriangular and $G$ positive
definite with $U^TSU = \operatorname{diag}((-1)^{d_i})\,G$. Then $(-1)^{\sum_{j<k}d_j}\det S_{\le k}>0$ for all
$k$, and every constant-degree layer Schur complement is $(-1)^r$-definite.
\end{theorem}

\noindent File \leanpath{VOA/\allowbreak Examples/\allowbreak PolarizedReesStrongGenerator.lean}.
\begin{leancode}
theorem sign_det_leadingPrincipal_pos (Q : MinorPayload n) {k : ℕ}
    (hk : k ≤ n) :
    0 < (-1 : ℝ) ^ cumulativeDegree Q.d hk
        * (leadingPrincipal Q.S hk).det := by
\end{leancode}

\begin{dictionary}
\lean{MinorPayload n} bundles $S,d,U,G$ and the hypotheses (field \lean{cross : Uᵀ * S * U = signDiag d * G}).
The file's docstring makes an unusually honest ``logical-distance disclosure'': over $\lnineteenR$ the field
\lean{cross} is \emph{equivalent} to the conclusion (by $LDL^T$ decomposition), so the theorem is a
repackaging; all the difficulty lies in producing \lean{cross} for a real model, which the file does not do.
We agree with that assessment.
\end{dictionary}

\begin{physics}
The finite endgame of a proposed ``polarised bi-Rees'' proof: if a deformation argument produced the normal
form, graded unitarity at that weight would follow.
\end{physics}

\begin{status}
Proved, axiom-clean. Interface-level; the hypotheses are not produced for any VOA.
\end{status}

The remaining three files are interfaces:\\ \leanpath{ScreenedReesInterface.lean} (a structure
\lean{ScreenedRees.Hypotheses} of filtration and Felder-descent conditions on a hypothetical screened map,
plus the trivial transfer \lean{isPosOn_of_metricTransport}: an injective map that pulls a positive form back
to a positive multiple of the source form transfers positivity); \leanpath{Logarithmic\allowbreak HeterotypicJet.lean}
(first-order jets $f_0+\epsilon f_1$ of linear maps and the Leibniz/Bockstein bookkeeping, pure algebra);
\leanpath{FreeFieldPi\allowbreak RFiltration.lean} (a combinatorial degree \lean{rDegree} on free-field monomial labels,
``number of derivatives of $\delta-\phi$ plus lattice charge'', after Beem--Meneghelli--Rastelli%
). Their docstrings state their limits accurately.

%======================================================================
\section{Drinfeld--Sokolov reduction for \texorpdfstring{$\mathfrak{sl}_2$}{sl(2)}}\label{l19:sec:ds}
%======================================================================

\subsection{The idea}

Quantum Hamiltonian (Drinfeld--Sokolov, DS) reduction produces $W$-algebras from affine Kac--Moody
algebras by ``gauging'' a nilpotent subalgebra. For $\lnineteensl$ at level $k$ and the principal nilpotent
$f$, one imposes the first-class constraint $e(z) = 1$ on the raising current. BRST-quantising it introduces
one pair of fermionic ghosts $(b,c)$ and the BRST charge
\begin{equation*}
 Q = \oint c(z)\,\big(e(z)+1\big)\,dz,\qquad Q^2 = 0,
\end{equation*}
on the complex $V^k(\lnineteensl)\otimes\mathrm{bc}$; the cohomology is the Virasoro VOA of central charge
\begin{equation*}
 c(k) = 13-6(k+2)-\frac{6}{k+2}.
\end{equation*}
Equivalently (Feigin--Frenkel), the reduction is the kernel of a screening operator
$\oint e^{\alpha_+\phi}$ on a free boson with background charge. At admissible levels $k = -2+p/q$,
$c(k) = c_{p,q}$ and the simple quotient of the reduction is the minimal model (Arakawa). Physically the
reduction is the passage from $SL(2)$ WZW models to Liouville/minimal models by gauging (Polyakov's
$SL(2)$ gravity; Bershadsky--Ooguri).

\begin{theorem}[DS arithmetic]\label{l19:thm:dsarith}
(a) For $p,q\ne0$, $c(-2+p/q) = 1-6(p-q)^2/(pq)$. (b) For $k+2\ne0$, $c(k) = \frac{3k}{k+2} + (-2) - 12k\cdot\frac12$,
where $\frac{3k}{k+2}$ is the library's Sugawara central charge of $V^k(\lnineteensl)$, $-2$ the library's $bc$
central charge, and $\frac12 = \langle\rho^\vee,\rho^\vee\rangle$ computed from the library's Killing form.
\end{theorem}

\noindent File \leanpath{VOA/\allowbreak Examples/\allowbreak DSReductionSl2.lean}.
\begin{leancode}
noncomputable def dsCentralCharge (k : ℂ) : ℂ := 13 - 6 * (k + 2) - 6 / (k + 2)

theorem dsCentralCharge_eq_minimalModel {p q : ℂ} (hp : p ≠ 0) (hq : q ≠ 0) :
    dsCentralCharge (-2 + p / q) = minimalModelCentralCharge p q := by

theorem dsCentralCharge_eq_affine_add_ghost (k : ℂ) (hk : k + 2 ≠ 0) :
    dsCentralCharge k
      = (AffineVOA.affineSl2VOA_closed k hk).centralCharge
        + Ghost.ghostCentralCharge 1 - 12 * k * rhoCheckNormSq := by
\end{leancode}

\begin{dictionary}
\lean{minimalModelCentralCharge p q} $= 1-6(p-q)^2/(pq)$. \lean{affineSl2VOA_closed k hk} is the universal
affine VOA $V^k(\lnineteensl)$ of Chapter~\ref{ch:L08}, with central charge $3k/(k+2)$.
\lean{Ghost.ghostCentralCharge 1} is a \emph{number} $-2$ (the library proves no Virasoro structure on the ghost
system; Chapter~\ref{ch:L17}); \lean{rhoCheckNormSq} $= \kappa(h,h)/4 = 1/2$. So (b) is an arithmetic identity
whose terms are tied to library objects; it is not a theorem that some BRST cohomology has that central charge
(the docstring says so).
\end{dictionary}

\paragraph{Proof on paper.} Field arithmetic with $t = k+2$: $13-6t-6/t = 3k/t-2-6k$. \qed

\begin{physics}
The ghost-number anomaly count: $c_W = c_{\text{WZW}} + c_{\text{ghosts}} - 12k|\rho^\vee|^2$, the last term
coming from the improvement of the stress tensor needed to make the constraint $e=1$ conformally invariant.
\end{physics}

\begin{status}
Proved, axiom-clean. Arithmetic only.
\end{status}

\begin{theorem}[The BRST charge squares to zero]\label{l19:thm:dsq}
For $k+2\ne0$, the operator $Q = \sum_m (e+\lnineteenvac)_{(m)}\otimes c_{(-1-m)}$ on
$V^k(\lnineteensl)\otimes\mathrm{bc}$ satisfies $Q^2 = 0$.
\end{theorem}

\noindent File \leanpath{VOA/\allowbreak Examples/\allowbreak DSReductionSl2.lean}.
\begin{leancode}
noncomputable def dsQ (k : ℂ) (hk : k + 2 ≠ 0) :
    End ℂ (AffineVOA.WakModule k ⊗[ℂ] SympFerm.FockSF) :=
  TensorVA.rawCoeff (affSVA k hk) Ghost.bcGhostSVA (ePlusOne k) cGhost (-1)

theorem dsQ_mul_self (k : ℂ) (hk : k + 2 ≠ 0) : dsQ k hk * dsQ k hk = 0 :=
  rawCoeff_mul_self_eq_zero _ _ (coeffCommute_ePlusOne k hk) coeffAnticommute_cGhost (-1)
\end{leancode}

\begin{dictionary}
\lean{WakModule k} is the carrier of $V^k(\lnineteensl)$, \lean{SympFerm.FockSF} the Fock carrier of the $bc$
system \lean{Ghost.bcGhostSVA} (Chapter~\ref{ch:L17}), \lean{affSVA} is $V^k$ viewed as a super vertex algebra,
\lean{ePlusOne k} $= e_{(-1)}\lnineteenvac+\lnineteenvac$, \lean{cGhost} the state $c$. \lean{TensorVA.rawCoeff ... (-1)} is
the $z^{-1}$-coefficient of $Y(e+\lnineteenvac,z)\otimes Y(c,z)$, i.e.\ the zero mode of the tensor field
(\lean{dsQ_eq_mode_tenY}). The one-line proof term exhibits the two inputs: the coefficients of $Y(e+\lnineteenvac,z)$
commute and those of $Y(c,z)$ anticommute.
\end{dictionary}

\paragraph{Proof on paper.} $Q^2 = \sum_{m,n}a_mb_{-1-m}a_nb_{-1-n}$ with $a_ma_n = a_na_m$ ($[e_m,e_n] = 0$:
no structure constant and $\kappa(e,e) = 0$) and $b_ib_j = -b_jb_i$ ($\{c_i,c_j\} = 0$); the swap $(m,n)\mapsto(n,m)$
negates the summand, so the sum vanishes. \qed

\begin{physics}
Nilpotency of the BRST charge is gauge invariance of the constrained WZW model; the constant ``$+1$'' does not
spoil it because it commutes with everything.
\end{physics}

\begin{status}
Proved, axiom-clean, unconditional. \textbf{What is not done:} the BRST cohomology. \lean{dsReductionSl2} takes
as \emph{input} a \lean{BRST.BRSTComplex} whose differential is \lean{dsQ}; none is constructed. Two obstructions are
recorded: the super Leibniz rule \lean{Q_deriv} is unavailable, and, as \leanpath{DSMorphismSl2.lean} shows,
\lean{dsQ} splits into a weight-preserving half and a half that lowers weight by one
(\lean{dsQ_split}, \lean{dsQlow_lowers_grading}), so the \lean{Q_wt} axiom of \lean{BRSTComplex} (Chapter~\ref{ch:L04},
Definition~\ref{l04:def:brst}) cannot hold for the naive grading --- one needs the improved DS grading, which is not
built.
\end{status}

\begin{theorem}[The DS complex, the embedding, and the screening route]\label{l19:thm:dsmor}
(a) For $k+2\ne0$, $V^k(\lnineteensl)\otimes\mathrm{bc}$ is a super vertex algebra and $x\mapsto x\otimes\lnineteenvac$ is an
injective homomorphism of vertex superalgebras. (b) For $p\ge1$ (level $k = p-2$), the kernel of the screening
operator $(e^{\alpha})_{(0)}$ on the rank-one lattice VOA $V_{\sqrt{2p}\Z}$ with deformed conformal vector is a VOA of
central charge $c(p-2)$, a proper subalgebra, and its inclusion is a VOA homomorphism.
\end{theorem}

\noindent File \leanpath{VOA/\allowbreak Examples/\allowbreak DSMorphismSl2.lean}.
\begin{leancode}
noncomputable def dsComplex (k : ℂ) (hk : k + 2 ≠ 0) :
    SuperVertexAlgebra (AffineVOA.WakModule k ⊗[ℂ] SympFerm.FockSF) :=
  TensorVA.tensorSVA (DSsl2.affSVA k hk) Ghost.bcGhostSVA
    (parityCompatible_toSuper (AffineVOA.Vir.affineVA k hk)) parityCompatible_bcGhostSVA
    (dsTensorGradingHyp k hk)

theorem dsEmbed_injective (k : ℂ) : Function.Injective (dsEmbed k) :=

theorem dsLatticeScreening_centralCharge (p : ℕ) (hp : 1 ≤ p) :
    (dsLatticeScreening p hp).centralCharge = DSsl2.dsCentralCharge (dsLevelOfP p) := by

theorem dsLatticeScreening_ne_top (p : ℕ) (hp : 1 ≤ p) :
    (Screening.screeningKernel (LatGen.Triplet.tripletLatticeVOA p hp)
      (dsScreeningState p) (dsScreeningState_mem_grade_one p hp)).carrier ≠ ⊤ := by

theorem dsScreeningIncl_isVOAHom (p : ℕ) (hp : 1 ≤ p) :
    IsVOAHom (dsLatticeScreening p hp) (LatGen.Triplet.tripletLatticeVOA p hp)
      (dsScreeningIncl p hp) where
\end{leancode}

\begin{dictionary}
\lean{TensorVA.tensorSVA} is the tensor product of Chapter~\ref{ch:L04} (Theorem~\ref{l04:thm:tensor}); its
hypothesis \lean{TensorGradingHyp}, open there, is discharged here (\lean{dsTensorGradingHyp}).
\lean{tripletLatticeVOA p hp} is the lattice VOA $V_{\sqrt{2p}\Z}$ with the background-charge conformal vector
(Chapter~\ref{ch:L18}); \lean{dsScreeningState p} is $e^{\alpha}$, $\alpha^2 = 2p$, of weight one;
\lean{Screening.screeningKernel} is Definition~\ref{l04:def:scrkernel}; \lean{dsLevelOfP p} $= p-2$;
\lean{IsVOAHom} preserves vacuum and all products. Note $c(p-2) = 1-6(p-1)^2/p = c_{1,p}$, the logarithmic
$(1,p)$ central charge. \textbf{Not claimed:} that the screening kernel is the Virasoro VOA, nor anything at
non-integer $k+2$ (Lee--Yang's $k=-8/5$ is out of reach of this route, since $\alpha^2 = 4/5$ is not an even
integer). \lean{RoutesAgreeSl2} (the Feigin--Frenkel comparison) is a \lean{def : Prop}, unproved.
\end{dictionary}

\paragraph{Proof on paper.} (a) The tensor-product construction plus two graded-linear-algebra facts
(independence of the convolution grading, via the weight operator; finite-dimensionality, since both factors
are $\N$-graded). Injectivity: $\lnineteenvac_{bc}\ne0$. (b) The kernel of a weight-one screening zero mode is a
vertex subalgebra containing $\omega$; its central charge is the ambient one, $1-12\lambda_p^2 = c(p-2)$; it is
proper because $e^{-\alpha}$ is not annihilated ($(e^\alpha)_{(0)}e^{-\alpha}\ne0$). \qed

\begin{physics}
(a) is the statement that the affine currents survive in the BRST complex as the ghost-number-zero,
ghost-vacuum sector; (b) is the Feigin--Frenkel ``free-field'' description of the reduced theory at integral
$k+2$, where the screening momenta are lattice vectors.
\end{physics}

\begin{status}
All proved, axiom-clean, unconditional. The identification of either route with $L(c,0)$, and the Feigin--Frenkel
theorem, are open.
\end{status}

%======================================================================
\section{Weight-one positivity of affine \texorpdfstring{$\mathfrak{sl}_2$}{sl(2)} at level \texorpdfstring{$-8/5$}{-8/5}}\label{l19:sec:aff85}
%======================================================================

Before the all-level theorem of Section~\ref{l19:sec:belowcritical}, the library records the first unconditional
positivity fact on the affine side: at the boundary admissible level $k=-8/5$, whose simple quotient is the chiral
algebra of the $(A_1,D_5)$ Argyres--Douglas theory, the three weight-one current states of $V^{-8/5}(\lnineteensl)$ have
the graded-unitarity sign.


\begin{theorem}[Weight-one Hermitian positivity at $k=-8/5$]\label{l19:thm:aff85}
On the three weight-one current states of $V^{-8/5}(\lnineteensl)$ the canonical invariant form is
$(J^a|J^b) = \tfrac85\kappa(a,b)$. With the compact conjugation $\theta(z_e,z_f,z_h) =
(-\bar z_f,-\bar z_e,-\bar z_h)$, the signed form $-(\theta z|w)$ equals
$\tfrac85(\bar z_ew_e+\bar z_fw_f+2\bar z_hw_h)$, which is positive definite.
\end{theorem}

\noindent File \leanpath{VOA/\allowbreak Examples/\allowbreak AffineSl2GradedUnitarityLow.lean}.
\begin{leancode}
theorem canonicalForm85_currentState (a b : Fin 3) :
    canonicalForm85.form (currentState (-8 / 5 : ℂ) a)
      (currentState (-8 / 5 : ℂ) b) = (8 / 5 : ℂ) * kappa a b := by

theorem compactCurrentSignedForm85_eq_canonical (z w : CurrentCoords) :
    compactCurrentSignedForm85 z w =
      -canonicalForm85.form (currentCoordinates85 (compactCurrentConj z))
        (currentCoordinates85 w) := by

theorem compactCurrentSignedForm85_positive :
    IsPosOn compactCurrentSignedForm85 (⊤ : Submodule ℂ CurrentCoords) := by
\end{leancode}

\begin{dictionary}
\lean{canonicalForm85} is the canonical normalised invariant form on $V^{-8/5}(\lnineteensl)$;
\lean{currentState k a} is $J^a_{-1}\lnineteenvac$; \lean{kappa a b} the Killing matrix in the Chevalley basis
$(e,f,h)$. \lean{CurrentCoords} $=\C^3$, \lean{currentCoordinates85 z} $= \sum_az_aJ^a_{-1}\lnineteenvac$,
\lean{compactCurrentConj} is $\theta$. \lean{compactCurrentSignedForm85} is \emph{defined} by the displayed
explicit formula; so the content is the identity with $-(\theta z|w)$ and the sign $\tfrac85>0$. The sign
$+\tfrac85$ (for $k<0$) comes from the FHL convention $u_{(1)}v = -(u|v)\lnineteenvac$. \textbf{Scope:} $\theta$ is
\emph{not} a \lean{RealStructure} on the VOA (the docstring says so: no mode covariance, no descent), and the
statement concerns three vectors only.
\end{dictionary}

\paragraph{Proof on paper.} $J^a_{(1)}J^b_{-1}\lnineteenvac = k\kappa(a,b)\lnineteenvac$, so $(J^a|J^b) = -k\kappa(a,b) =
\tfrac85\kappa(a,b)$. Then $-(\theta z|w) = -\tfrac85\sum\kappa_{ab}(\theta z)_aw_b$, and with $\kappa_{ef} = 1$,
$\kappa_{hh} = 2$ this is the displayed sum of squares on the diagonal. \qed

\begin{physics}
The currents of $L_{-8/5}(\lnineteensl)$ --- the flavour currents of the $(A_1,D_5)$ Argyres--Douglas theory, whose
$SU(2)$ flavour symmetry has $k_{4d} = -2k_{2d} = 16/5>0$ --- have the graded-unitarity sign $(-1)^1$ at
$R$-charge one: a 4d unitarity bound seen in 2d.
\end{physics}

\begin{status}
Proved, axiom-clean, unconditional. Also \lean{simpleCorrectedGoodFiltration85}: an unconditional good
filtration on the simple quotient, induced from the corrected affine filtration.
\end{status}

%======================================================================
\section{Transfer of graded unitarity along isometric embeddings}\label{l19:sec:transfer}
%======================================================================

This section records a piece of linear algebra about the graded-unitarity predicate that concerns no particular VOA:
graded unitarity pulls back along an injective, weight-preserving, isometric and $\rho$-compatible linear map,
provided the map is \emph{orthogonally reducing}, i.e.\ compatible with the flags and with their orthogonal layers. It
is the abstract input of the coset transfer of Chapter~\ref{ch:L28}, which instantiates it through a filtered linear
isomorphism (\leanpath{LeeYangCosetTransferFiltered.lean}). The Lean file and namespace carry the prefix
\lean{InverseDS} for historical reasons; nothing in this section depends on the inverse Drinfeld--Sokolov programme.

\begin{definition}[Graded isometric embedding; orthogonal reduction]\label{l19:def:gie}
Given VOAs $A,T$ with invariant forms and real structures, a \emph{graded isometric embedding} is an injective
linear $\iota:A\to T$ preserving weights, with $B_A(x,y) = B_T(\iota x,\iota y)$ and $\iota\rho_A = \rho_T\iota$.
It is \emph{orthogonally reducing} at weight $n$ for flags $F_A,F_T$ if (1) $\iota(\lnineteenG^A_p) =
\iota(A_n)\cap\lnineteenG^T_p$; (2) $B_T$ is nondegenerate on each $\iota(\lnineteenG^A_p)$; (3)
$\iota(\lnineteenLay^A_p) = \iota(A_n)\cap\lnineteenLay^T_p$.
\end{definition}

\noindent File \leanpath{VOA/\allowbreak Examples/\allowbreak InverseDSGradedUnitarityTransfer.lean}.
\begin{leancode}
structure GradedIsometricEmbedding
    (A : VertexOperatorAlgebra Affine) (T : VertexOperatorAlgebra Ambient)
    (BA : InvariantForm A) (BT : InvariantForm T)
    (RA : RealStructure A) (RT : RealStructure T) where
  /-- The underlying complex-linear embedding. -/
  iota : Affine →ₗ[ℂ] Ambient
  /-- Injectivity of the realization. -/
  iota_injective : Function.Injective iota
  /-- The realization preserves conformal weight. -/
  map_grading : ∀ (n : ℤ) (x : Affine), x ∈ A.grading n → iota x ∈ T.grading n
  /-- The affine invariant form is the exact pullback of the ambient invariant form. -/
  form_isometry : ∀ x y, BA.form x y = BT.form (iota x) (iota y)
  /-- The realization intertwines the source and ambient anti-linear structures. -/
  rho_compatible : ∀ x, iota (RA.rho x) = RT.rho (iota x)

structure OrthogonallyReducingAt
    (E : GradedIsometricEmbedding A T BA BT RA RT)
    (FA : GradedFlag A) (FT : GradedFlag T) (n : ℤ) : Prop where
  /-- The source filtration is exactly induced from the ambient filtration on the image. -/
  piece_strict : ∀ p : ℕ,
    (FA.pieceIn p n).map (E.weightMap n) =
      LinearMap.range (E.weightMap n) ⊓ FT.pieceIn p n
  /-- Each mapped cumulative piece is nondegenerate for the ambient bilinear form. -/
  mapped_piece_nondegenerate : ∀ p : ℕ,
    ((LiRadical.gradedForm BT n).restrict
      ((FA.pieceIn p n).map (E.weightMap n))).Nondegenerate
  /-- Source orthogonal layers map to the image intersections with ambient orthogonal layers. -/
  layer_strict : ∀ p : ℕ,
    (FA.layer n (LiRadical.gradedForm BA n) p).map (E.weightMap n) =
      LinearMap.range (E.weightMap n) ⊓
        FT.layer n (LiRadical.gradedForm BT n) p
\end{leancode}

\begin{dictionary}
\lean{E.weightMap n} is $\iota$ restricted to $A_n\to T_n$; \lean{LinearMap.range} its image. Exactly
Definition~\ref{l19:def:gie}. \lean{GradedIsometricEmbedding} contains \emph{no} vertex-algebra compatibility
(no $\iota(u_{(n)}v) = \iota u_{(n)}\iota v$); the docstring notes that the transfer does not need it.

\textbf{Why (2) and (3) are needed.} The file docstring gives a counterexample showing that strictness (1)
alone does not transfer graded unitarity: on $\C^2$ with form $\operatorname{diag}(1,-1)$, ambient flag
$\lnineteenG_0 = \C e_1$, $\lnineteenG_1 = \C^2$ (graded unitary: $\lnineteenLay_0 = \C e_1$ positive, $\lnineteenLay_1 = \C e_2$
negative-then-twisted-positive), and image line $\C(e_1+ae_2)$, $0<|a|<1$: the induced source flag has
$\lnineteenG^A_0 = 0$, $\lnineteenG^A_1 = $ the line, the source form $1-|a|^2>0$ is positive, but its only nonzero layer
has $p = 1$, where $(-1)^1$ makes it negative. (We checked this by hand.)
\end{dictionary}

\begin{physics}
An isometric realisation (free-field or coset) that respects the $R$-filtration \emph{and} its orthogonal splitting.
The counterexample says: a state can sit in the image of a filtration step without being orthogonal to the
lower-charge ambient states, and then its ``charge'' in the source and in the ambient differ.
\end{physics}

\begin{status}
Definitions. Instantiated in the library by the coset transfer of Chapter~\ref{ch:L28}, for instance through the
filtered linear isomorphisms of
\lean{gradedUnitary_}\allowbreak\lean{of_}\allowbreak\lean{filteredLinearEquiv}.
\end{status}

\begin{theorem}[Transfer of graded unitarity]\label{l19:thm:transfer}
If $\iota$ is a graded isometric embedding, orthogonally reducing at every weight, and $(T,B_T,\rho_T,F_T)$
is graded unitary, then $(A,B_A,\rho_A,F_A)$ is graded unitary.
\end{theorem}

\noindent File \leanpath{VOA/\allowbreak Examples/\allowbreak InverseDSGradedUnitarityTransfer.lean}.
\begin{leancode}
theorem gradedUnitary_of_ambient
    (E : GradedIsometricEmbedding A T BA BT RA RT)
    (FA : GradedFlag A) (FT : GradedFlag T)
    (hreduce : ∀ n : ℤ, E.OrthogonallyReducingAt FA FT n)
    (hambient : GradedUnitary BT RT FT) :
    GradedUnitary BA RA FA := by
\end{leancode}

\begin{dictionary}
Exactly Theorem~\ref{l19:thm:transfer}. All hypotheses explicit.
The weight-by-weight version is \lean{gradedUnitaryAt_}\allowbreak\lean{of_}\allowbreak\lean{ambient}. When $\iota$ maps
every flag step \emph{onto} the ambient step, orthogonal reduction is automatic; this is the form used by the coset
transfer of Chapter~\ref{ch:L28}.
\end{dictionary}

\paragraph{Proof on paper.} Clause (iv): nondegeneracy pulls back along the isometry from (2). Clause (vi): by (3)
$\iota(\lnineteenLay^A_p)\subseteq\lnineteenLay^T_p$, where the twisted form is positive; pull back by
Theorem~\ref{l19:thm:pullback}, using $\rho$-compatibility to identify the sesquilinear forms. \qed

\begin{physics}
Positivity in the bigger theory, sector by sector, gives positivity in the smaller one.
\end{physics}

\begin{status}
Proved, axiom-clean; conditional on the explicit hypotheses \lean{hreduce} and \lean{hambient}.
\end{status}

%======================================================================
\section{Graded unitarity of \texorpdfstring{$V^k(\mathfrak{sl}_2)$}{V\textasciicircum k(sl(2))} below the critical level}\label{l19:sec:belowcritical}
%======================================================================

\paragraph{What this section adds.} (Its Lean quotations were taken at a later revision of the library than the rest of the
chapter; the files are unchanged since.) It documents the twelve files
\leanpath{VOA/\allowbreak Examples/\allowbreak AffineSl2BelowCritical*.lean} (\leanpath{Anchor}, \leanpath{Compact},
\leanpath{Flag}, \leanpath{GradedUnitarity}, \leanpath{GramMatrix}, \leanpath{MatrixAnchor},
\leanpath{MatrixRigidity}, \leanpath{NoWall}, \leanpath{Nondegenerate}, \leanpath{PBWFlag},
\leanpath{RealStructure}, \leanpath{Rigidity}). Their headline is the first theorem of the library that asserts
the predicate \lean{GradedUnitary} of Definition~\ref{l19:def:gu} --- all of it, at every conformal weight
and every filtration degree --- for a concrete vertex operator algebra, with no hypothesis beyond the range of
the level:

\begin{quote}
\emph{For every real $k<-2$, the universal affine vertex operator algebra $V^k(\lnineteensl)$, with its normalised
invariant form, the compact real structure and the current-length flag, is graded unitary.}
\end{quote}

The algebra is the \emph{universal} one, $V^k(\lnineteensl)$, not a simple quotient; the level is \emph{below}
the critical value $-2$; and the flag is the \emph{ordinary} length filtration by number of currents.
Nothing here is about the simple quotients $L_k(\lnineteensl)$ at the admissible levels $k>-2$ that occur in
the SCFT/VOA correspondence (for instance as chiral algebras of Argyres--Douglas theories); those are a different problem
(Remark~\ref{l19:rem:bcscope}).

\paragraph{Reading the Lean displays of this section.} Each display is copied from the cited files. A line
\texttt{...} inside a display marks omitted source lines: other declarations lying between two quoted ones, or the
remaining (proof) fields of a structure instance cut after \lean{where}. Fragments separated only by a blank line are
consecutive in the source, apart from docstrings and comments, which are always omitted; a statement ending in
\lean{:= by} has its proof omitted. Full Lean names are obtained by prefixing the namespace of the file:
\begingroup\raggedright
\begin{itemize}
\item \lean{universal}, \lean{universalForm}: \lean{VOA.AffineSl2GeneralRigidity};
\item \lean{stepMatrix}, \lean{MatrixStepNondeg}, \lean{MatrixLayerSign}, \lean{matrixLayerSign_rigid},
  \lean{matrixLayerSign_eventually} (files \leanpath{AffineSl2BelowCriticalMatrixRigidity.lean},
  \leanpath{AffineSl2BelowCriticalMatrixAnchor.lean}): \lean{VOA.MatrixRigidity};
\item \lean{CurrentSystem} and its theorem \lean{layerNondegenerate_of_lower} (file
  \leanpath{AffineSl2BelowCriticalNoWall.lean}): \lean{VOA.BelowCriticalNoWall} (the theorem inside
  \lean{namespace CurrentSystem});
\item \lean{adjointColour}: \lean{VOA.AffineVOA.HighestModule};
\item everything else in the twelve files: \lean{VOA.AffineSl2BelowCritical}, with \lean{sugawara_sesq} inside the
  sub-namespace \lean{IsCompactRealStructure}.
\end{itemize}
So, for instance, the headline is \lean{VOA.AffineSl2BelowCritical.}\allowbreak\lean{universal_}\allowbreak\lean{gradedUnitary_}\allowbreak\lean{belowCritical_}\allowbreak\lean{unconditional};
the audit listing of Section~\ref{l19:sec:axioms} uses these full names.\par
\endgroup

\paragraph{The route in one paragraph.} Fix a conformal weight $E$. In the basis of ordered PBW monomials of
weight $E$ the Hermitian Gram matrix $G_E(k)$ has entries that are \emph{polynomials} in the level. Three facts
are combined, by induction on $E$:
\begin{enumerate}
\item \emph{Anchor} (Theorem~\ref{l19:thm:bcanchor}): for $k$ sufficiently negative, $G_E(k)$ has all the
  graded-unitary layer signs. Reason: after rescaling a monomial of length $r$ by $|k|^{-r/2}$, the Gram matrix
  tends to the diagonal matrix $\operatorname{diag}((-1)^r d)$ with $d>0$ (Wick's theorem for free currents).
\item \emph{No wall} (Theorem~\ref{l19:thm:bcnowall}): if all lower weights carry the layer signs and the full
  weight-$E$ form is nondegenerate (true for every $k<-2$ by the Kac--Kazhdan determinant), then the form is
  nondegenerate on every flag step at weight $E$. The proof is a positivity argument with the Sugawara
  identity.
\item \emph{Rigidity} (Theorem~\ref{l19:thm:bcrigid}): a continuous family of Hermitian matrices whose flag steps
  are all nondegenerate cannot change its layer signs on a connected parameter set; here the set is the
  interval $(-\infty,-2)$.
\end{enumerate}
Then: no wall gives nondegeneracy at weight $E$ for all $k<-2$; rigidity transports the signs from the anchor
to every $k<-2$; and the induction continues.

%----------------------------------------------------------------------
\subsection{The data: the algebra, the form, the real structure, the flag}
%----------------------------------------------------------------------

Recall from Chapter~\ref{ch:L08} that the vacuum module of affine $\lnineteensl$ at level $k$ is \lean{WakModule k}
(a misleading name: it is \emph{not} a Wakimoto module), with current modes $J^a_m$ (\lean{J k a m}) indexed by
the Chevalley basis $a\in\{0,1,2\} = \{e,f,h\}$ and satisfying
\begin{equation}\label{l19:eq:bccomm}
  [J^a_m,J^b_n] = J^{[a,b]}_{m+n} + m\,k\,\kappa(a,b)\,\delta_{m+n,0},\qquad
  \kappa(e,f) = \kappa(f,e) = 1,\ \kappa(h,h) = 2,
\end{equation}
with $[h,e] = 2e$, $[h,f] = -2f$, $[e,f] = h$ (Theorem~\ref{l08:thm:currentcomm}, Definition~\ref{l08:def:sl2}).
The creation word \lean{wordStateA k w} attached to a list $w = [(a_1,n_1),\dots,(a_r,n_r)]$ of (colour, depth)
pairs is $J^{a_1}_{-n_1-1}\cdots J^{a_r}_{-n_r-1}\lnineteenvac$ (Definition~\ref{l08:def:words}); note the
shift: depth $n$ means mode $-n-1$. Its conformal weight is $\sum_i(n_i+1)$ and its \emph{length} is $r$.

\begin{definition}[The data of the below-critical theorem]\label{l19:def:bcdata}
Let $k\in\C$, $k\ne-2$.
\begin{enumerate}
\item $V^k = V^k(\lnineteensl)$ is the universal affine VOA with the Sugawara conformal vector
  (Definition~\ref{l08:def:Vk}). For a real number $x<-2$ the complex number $x$ is noncritical.
\item $B_k$ is the \emph{normalised canonical invariant form} on $V^k$: the unique invariant symmetric bilinear
  form with $B_k(\lnineteenvac,\lnineteenvac) = 1$ (Theorem~\ref{l04:thm:canonical}). In the FHL convention
  used by the library, currents are \emph{skew}-adjoint: $B_k(J^a_mu,v) = -B_k(u,J^a_{-m}v)$.
\item The \emph{colour reversal} $a\mapsto a^*$ exchanges $e$ and $f$ and fixes $h$. A real structure $\rho$ on
  $V^k$ (Definition~\ref{l19:def:real}) is \emph{compact} if $\rho(J^a_{-1}\lnineteenvac) = -J^{a^*}_{-1}\lnineteenvac$
  for each colour $a$ and its parity scalars $s(n)$ are all nonzero. (Recall from Definition~\ref{l19:def:real} that a
  real structure carries a function $s:\Z\to\C$, its $R$-parity, with $\rho^2 = s(n)$ on the weight-$n$ space. When
  $s(n)\ne0$, $\rho$ is invertible on $V_n$; this is what makes the orthogonal complement of a $\rho$-stable subspace
  the same for the bilinear form $B_k$ and for the Hermitian form $\mathcal H(x,y) = B_k(\rho x,y)$,
  Theorem~\ref{l19:thm:orthsame}. It is the only use of $s$ in this section.)
\item For real $k$ the \emph{compact real structure} is $\rho_k = \theta\circ\beta$, where $\beta$ is the
  anti-linear map fixing $\lnineteenvac$ and commuting with every $J^a_m$ (``complex conjugation of PBW
  coordinates'') and $\theta$ is the Chevalley automorphism ($J^a_m\mapsto -J^{a^*}_m$). Its parity is $s\equiv1$.
  The \emph{compact family} is $x\mapsto\rho_x$ for $x\in(-\infty,-2)$.
\item The \emph{ordinary flag}: $F_p\subseteq V^k$ is the span of all creation words of length $\le p$ (any
  colours, any order, any depths), and the flag piece at weight $n$ is $F_p\cap V^k_n$.
\end{enumerate}
\end{definition}

{\raggedright\noindent Files \leanpath{VOA/\allowbreak Examples/\allowbreak AffineSl2GeneralRigidityBasis.lean} (the abbreviation
\lean{universal}, under \lean{variable (k : ℂ) (hk : k + 2 ≠ 0)}),
\leanpath{VOA/\allowbreak Examples/\allowbreak AffineSl2UniversalGram.lean} (\lean{universalForm}) and
\leanpath{VOA/\allowbreak Examples/\allowbreak AffineSl2BelowCriticalNondegenerate.lean} (\lean{noncritical_of_lt}):\par}
\begin{leancode}
noncomputable abbrev universal := affineSl2VOA_closed k hk

noncomputable def universalForm : InvariantForm (universal k hk) :=
  canonicalInvariantForm (AffineSl2BoundaryMaximalIdeal.affineSl2_vacuum_ne_zero k hk)
    (LQ43.L1_grading_one k hk)

theorem noncritical_of_lt {x : ℝ} (hx : x < -2) : ((x : ℂ) + 2) ≠ 0 := by
\end{leancode}
{\raggedright\noindent File \leanpath{VOA/\allowbreak Examples/\allowbreak AffineSl2VermaDual.lean} (colour reversal) and
\leanpath{VOA/\allowbreak Examples/\allowbreak AffineSl2BelowCriticalCompact.lean} (compactness):\par}
\begin{leancode}
def adjointColour (a : Fin 3) : Fin 3 := if a=0 then 1 else if a=1 then 0 else 2

structure IsCompactRealStructure (R : RealStructure (universal k hk)) : Prop where
  current : ∀ a : Fin 3, R.rho (currentState k a) = -currentState k (adjointColour a)
  parity_ne_zero : ∀ n : ℤ, R.parity n ≠ 0
\end{leancode}
{\raggedright\noindent File \leanpath{VOA/\allowbreak Examples/\allowbreak AffineSl2BelowCriticalRealStructure.lean}
(the construction; \lean{hreal : (starRingEnd ℂ) k = k} says $k$ is real):\par}
\begin{leancode}
def basisConj (x : WakModule k) : WakModule k := ofConj (basisConjₗ hreal x)
...
def compactRho : WakModule k →+ WakModule k where
  toFun x := vacuumChevalley k (basisConj hreal x)
...
def compactRealStructure : RealStructure (universal k hk) where
  rho := compactRho hreal
  antilinear := compactRho_smul hreal
...
  parity := fun _ => 1
...
def compactRealFamily : RealFamily := fun x hx =>
  compactRealStructure (conj_ofReal x) (noncritical_of_lt hx)
\end{leancode}
{\raggedright\noindent File \leanpath{VOA/\allowbreak Examples/\allowbreak AffineSl2BelowCriticalGradedUnitarity.lean}
(the type of families):\par}
\begin{leancode}
abbrev RealFamily : Type :=
  ∀ (x : ℝ) (hx : x < -2), RealStructure (universal (x : ℂ) (noncritical_of_lt hx))
\end{leancode}
{\raggedright\noindent File \leanpath{VOA/\allowbreak Examples/\allowbreak AffineSl2BelowCriticalFlag.lean}:\par}
\begin{leancode}
def ordinaryPiece (p : ℕ) : Submodule ℂ (WakModule k) :=
  Submodule.span ℂ (Set.range (fun w : {w : List (Fin 3 × ℕ) // w.length ≤ p} =>
    wordStateA k w.1))
...
def ordinaryWeightPiece (p : ℕ) (n : ℤ) : Submodule ℂ (WakModule k) :=
  ordinaryPiece k p ⊓ (universal k hk).grading n

def ordinaryFlag : GradedFlag (universal k hk) where
  piece := ordinaryWeightPiece k hk
  le_grading := fun _ _ => inf_le_right
  ...
\end{leancode}

\begin{dictionary}
\textbf{The algebra.} \lean{universal k hk} is \lean{affineSl2VOA_closed k hk}, the VOA $V^k(\lnineteensl)$ of
Definition~\ref{l08:def:Vk} (the same abbreviation is used in Chapter~\ref{ch:L09}). \lean{noncritical_of_lt hx}
is the proof that a real $x<-2$, viewed in $\C$, satisfies $x+2\ne0$; every statement below is about
\lean{universal (x : ℂ) (noncritical_of_lt hx)}, i.e.\ $V^x$ at the real level $x$.

\textbf{The form.} \lean{canonicalInvariantForm hne hqp} is the canonical invariant form of
Theorem~\ref{l04:thm:canonical}: given that the vacuum is nonzero (\lean{hne}, here the axiom-free
\lean{affineSl2_vacuum_ne_zero}) and that $L_1V_1 = 0$ (\lean{hqp}, here \lean{L1_grading_one}), it is the unique
invariant form normalised by $(\lnineteenvac|\lnineteenvac) = 1$. So $B_k$ is determined, not chosen. Its
convention gives $B_k(J^a_{-1}\lnineteenvac,J^b_{-1}\lnineteenvac) = -k\kappa(a,b)$ (compare
Theorem~\ref{l19:thm:aff85}, where $k = -8/5$ gives $\tfrac85\kappa$).

\textbf{Compactness.} \lean{currentState k a} is $J^a_{-1}\lnineteenvac$; \lean{adjointColour} is $a\mapsto a^*$
(colours are numbered $0 = e$, $1 = f$, $2 = h$). \lean{IsCompactRealStructure R} is clause 3 of the definition:
two fields, the action on the three currents and invertibility of the parity. From it alone the library proves
(\leanpath{AffineSl2BelowCriticalCompact.lean}) $\rho J^a_m = -J^{a^*}_m\rho$ (\lean{rho_J}), that $\rho$ preserves weights
and the ordinary flag, and that the Hermitian adjoint of $J^a_m$ for $\mathcal H(x,y) = B_k(\rho x,y)$ is $J^{a^*}_{-m}$
(\lean{sesq_J_adj}). The last is the $\mathfrak{su}(2)$ reality condition $(J^e_m)^\dagger = J^f_{-m}$,
$(J^h_m)^\dagger = J^h_{-m}$: the two minus signs (from $\rho$ and from skew-adjointness) cancel. Since moreover
$\mathcal H(\lnineteenvac,\lnineteenvac) = 1$, repeated use of this adjointness computes every Gram entry as a vacuum expectation
value, \eqref{l19:eq:bcwick} below.

\textbf{The construction.} \lean{basisConjₗ hreal} is the linear map from $V^k$ to the same space with
\emph{conjugated} scalar multiplication, obtained from the universal property of the vacuum module (the
structure constants and the level are real, so the currents still satisfy \eqref{l19:eq:bccomm} there); read back
in $V^k$ it is the anti-linear map $\beta$ with $\beta(\lnineteenvac) = \lnineteenvac$ and $\beta J^a_m = J^a_m\beta$.
\lean{vacuumChevalley k} is the Chevalley automorphism $\theta$ of Chapter~\ref{ch:L09}
(this section uses only $\theta\lnineteenvac = \lnineteenvac$ and $\theta J^a_m = -J^{a^*}_m\theta$,
\lean{vacuumChevalley_current}). Hence $\rho_k = \theta\beta$ sends the creation word of length $r$ to $(-1)^r$
times the word with all colours reversed (\lean{compactRho_wordStateA}). The elided fields of
\lean{compactRealStructure} are proofs: compatibility with every vertex operator (\lean{compactRho_YA}), $\rho\omega = \omega$
(\lean{compactRho_sugawara}), and $\rho^2 = 1$ (\lean{compactRho_compactRho}). \lean{conj_ofReal x} is the proof that
a real number is fixed by conjugation. A \lean{RealFamily} is a choice of real structure on $V^x$ for every real
$x<-2$; \lean{compactRealFamily} is the compact choice; \lean{compactRealFamily_isCompact} proves it compact.

\textbf{The flag.} \lean{Submodule.span ℂ (Set.range …)} is the span of the words with \lean{w.length ≤ p};
\lean{⊓} is intersection with the weight space \lean{(universal k hk).grading n}. The remaining fields of
\lean{ordinaryFlag} (elided) are the axioms of a \lean{GradedFlag} (Definition~\ref{l11:def:flag}): monotone in
$p$, $F_0\cap V_n = 0$ for $n>0$, $F_0\cap V_0 = \C\lnineteenvac$, exhaustive; each is proved. By filtered PBW
straightening (\leanpath{AffineSl2BelowCriticalPBWFlag.lean}, \lean{ordinaryPiece_eq_pbwSpan}) $F_p$ is also the
span of the \emph{ordered} PBW monomials of length $\le p$, so the ordered monomials of weight $E$ form a basis of
$V^k_E$ adapted to the flag.

\textbf{Is this the intended flag?} It is a fixed, explicit flag, not a free parameter, so the vacuity of
Warning~\ref{l19:warn:freeflag} does not apply. It is the natural ``$R$-filtration'' of the universal algebra
when each current carries $R$-charge one and nothing else is imposed. It is \emph{not} the ``corrected''
filtration used for simple quotients at admissible level (in which the Sugawara vector has degree one); here
the Sugawara vector $\omega$ has degree two. The file docstring says so.
\end{dictionary}

\begin{physics}
$J^a(z)$ are the $\mathfrak{su}(2)$ flavour currents. In the SCFT/VOA dictionary they come from the
four-dimensional moment-map operators, which have $SU(2)_R$ spin $R = 1$; a normal-ordered product of $r$ currents
then has $R\le r$, which is what the length flag records. The compact real structure is the statement that
the currents are Hermitian fields of the compact real form $\mathfrak{su}(2)$: $(J^\pm_m)^\dagger = J^\mp_{-m}$,
$(J^3_m)^\dagger = J^3_{-m}$, the standard conjugation of a WZW model. The level enters through the two-point
function $\langle J^a J^b\rangle\propto k\kappa^{ab}$, and $k_{\mathrm{4d}} = -2k$ is the four-dimensional flavour
central charge, positive exactly when $k<0$.
\end{physics}

\begin{status}
Definitions; \lean{compactRealStructure_isCompact} and \lean{compactRealFamily_isCompact} (the constructed $\rho$
is compact) and \lean{ordinaryPiece_eq_pbwSpan} are proved, \texttt{\#print axioms}-clean
(Section~\ref{l19:sec:axioms}).
\end{status}

%----------------------------------------------------------------------
\subsection{The headline}
%----------------------------------------------------------------------

\begin{theorem}[Below-critical graded unitarity of $V^k(\lnineteensl)$]\label{l19:thm:bcheadline}
Let $k$ be a real number with $k<-2$. Then $\bigl(V^k(\lnineteensl),B_k,\rho_k,F\bigr)$, with the data of
Definition~\ref{l19:def:bcdata}, is graded unitary in the sense of Definition~\ref{l19:def:gu}: for every
$n\in\Z$ and every $p\in\N$,
\begin{enumerate}
\item (clause (iv)) $B_k$ is nondegenerate on $F_p\cap V^k_n$;
\item (clause (vi)) on the layer $\lnineteenLay_p = (F_p\cap V^k_n)\cap(F_{p-1}\cap V^k_n)^{\perp}$ the form
  $(-1)^p\,\mathcal H_n$, $\mathcal H_n(x,y) = B_k(\rho_kx,y)$, is positive definite. (The orthogonal complement may be
  taken for $B_k$ or for $\mathcal H_n$: for this $\rho_k$, which preserves the flag pieces and has $s\equiv1$, the two
  agree by Theorem~\ref{l19:thm:orthsame}.)
\end{enumerate}
Clause (v) (the orthogonal decomposition into layers) then follows from Theorem~\ref{l19:thm:layerdec}.
\end{theorem}

{\raggedright\noindent File \leanpath{VOA/\allowbreak Examples/\allowbreak AffineSl2BelowCriticalAnchor.lean}.\par}
\begin{leancode}
theorem universal_gradedUnitary_belowCritical_unconditional (k : ℝ) (hk : k < -2) :
    GradedUnitary (universalForm (k : ℂ) (noncritical_of_lt hk)) (compactRealFamily k hk)
      (ordinaryFlag (k : ℂ) (noncritical_of_lt hk)) :=
  universal_gradedUnitary_belowCritical_of_anchor layerSignAnchor_compact k hk
\end{leancode}

\begin{dictionary}
\textbf{Hypotheses.} Exactly two arguments: a real number \lean{k : ℝ} and a proof \lean{hk : k < -2}. There are no
other hypotheses, no typeclass assumptions beyond those already fixed, no \lean{Prop}-valued structure fields to
be supplied later. The one-line proof term is informative and is shown: it combines the
conditional assembly \lean{universal_gradedUnitary_belowCritical_of_anchor} (which still takes the anchor as an
argument, Theorem~\ref{l19:thm:bcrigid}) with the proved anchor \lean{layerSignAnchor_compact}
(Theorem~\ref{l19:thm:bcanchor}).

\textbf{Decoding the conclusion.} \lean{(k : ℂ)} is the real level viewed in $\C$. \lean{GradedUnitary B R F} is
Definition~\ref{l19:def:gu}: for every \lean{n : ℤ}, \lean{GradedUnitaryAt F n (LiRadical.gradedForm B n)
(gradedSesq B R n)}, i.e.\ \lean{LayerNondegenerate} (clause (iv)) and \lean{LayerPositive} (clause (vi), with the
twist $(-1)^p$ of \lean{signedForm}). \lean{LiRadical.gradedForm B n} is $B$ restricted to $V_n\times V_n$,
\lean{gradedSesq B R n} is $\mathcal H_n(a,b) = B(\rho a,b)$ (Definition~\ref{l19:def:gradedsesq}, where it is called $h_n$; in this section we write $\mathcal H$ for the Hermitian form, because $h$ is the Cartan generator of $\lnineteensl$). The three data are
those of Definition~\ref{l19:def:bcdata}, each pinned: the form is the canonical one, $\rho$ is the constructed
compact one, the flag is the ordinary one.

\textbf{Same as the textbook statement?} Yes, for the data named. Two points of scope. First, the statement is
for one specific compact real structure; it is not stated for an arbitrary compact $\rho$ (although the abstract
assembly of Theorem~\ref{l19:thm:bcrigid} is). Second, it is for real $k<-2$ only; nothing is claimed for
complex $k$ or for $k>-2$. For every real $k>-2$ the conclusion is in fact false for this flag
(Example~\ref{l19:ex:bcweight2}): for $k>0$ clause (vi) fails at weight one ($G_1(k) = \operatorname{diag}(k,k,2k)$ is
positive, not negative, definite); for $-2<k<0$ and for $0<k<1$ it fails at weight two (the layer-2 Schur complement $S$
has $\det S<0$); at $k = 0$ the form vanishes on $V_1$ ($G_1(0) = 0$), so clause (iv) fails at weight one; and at $k = 1$,
$\det S = 0$, so the form is degenerate on the weight-two layer and clause (vi) fails there (a positive definite form
cannot have a null direction). The case $k = -2$ is excluded by the construction (critical level).

\textbf{Could it be vacuous?} The weight spaces are nonzero for $n\ge0$ (they contain the PBW monomials) and
zero for $n<0$; at negative weight both clauses hold trivially, which is the correct reading. The layers are
genuinely nonzero: at weight one the layer $\lnineteenLay_1$ is all of $V^k_1\cong\C^3$, and positivity of
$-\mathcal H_1$ there is the inequality $-k>0$ (Example~\ref{l19:ex:bcweight2}). The predicate \lean{IsPosOn} has teeth
(Definition~\ref{l19:def:isposon}): the same layers fail the condition for $k>0$.
\end{dictionary}

\paragraph{Proof on paper.} Write $\mathrm{(iv)}_E(k)$ and $\mathrm{(vi)}_E(k)$ for the two clauses at weight
$E$ and level $k$. We prove by strong induction on $E$ the statement $P(E)$: \emph{$\mathrm{(vi)}_E(k)$ holds for
every real $k<-2$}; along the way we obtain $\mathrm{(iv)}_E(k)$. For $E<0$ there is nothing to prove.
Assume $P(E')$ for all $E'<E$.
\begin{enumerate}
\item By Theorem~\ref{l19:thm:bcnondeg} the full form $B_k$ on $V^k_E$ is nondegenerate for every $k<-2$.
\item By the no-wall lemma (Theorem~\ref{l19:thm:bcnowall}), applied at each $k<-2$ with the induction hypothesis
  at the weights $E-m$, $m\ge1$, every flag step at weight $E$ is nondegenerate: $\mathrm{(iv)}_E(k)$ for every
  $k<-2$.
\item By the anchor (Theorem~\ref{l19:thm:bcanchor}) there is $K$ such that $\mathrm{(vi)}_E(x_0)$ holds for any
  $x_0<\min(K,-2)$.
\item By rigidity (Theorem~\ref{l19:thm:bcrigid}), which needs exactly $\mathrm{(iv)}_E(x)$ for all $x<-2$,
  $\mathrm{(vi)}_E(x_0)$ implies $\mathrm{(vi)}_E(k)$ for every $k<-2$.
\end{enumerate}
This is $P(E)$, and the induction closes. The quantifier ``for every $k<-2$'' inside the induction is essential:
step~2 at a single level $k$ uses lower weights at that same level, but step~4 needs step~2 at \emph{all} levels.
\qed

\begin{physics}
Graded unitarity is the shadow, in the chiral algebra, of four-dimensional unitarity: an operator built from $r$
moment maps has four-dimensional norm of sign $+$, and the two-dimensional form sees it with sign $(-1)^r$. For the
universal algebra the theorem says that the ``free'' current sector --- multi-current states with no relations
imposed --- passes this test whenever $k_{\mathrm{4d}} = -2k>4$. At very large $k_{\mathrm{4d}}$ this is the
generalized-free-field limit, in which the currents' correlators are Wick contractions and the signs are manifest;
the content of the theorem is that nothing goes wrong as $k$ increases all the way up to the critical level, where the Sugawara
vector becomes null (Example~\ref{l19:ex:bcweight2}). \emph{Heuristic caveat:} $V^k(\lnineteensl)$ at $k<-2$ has
Sugawara central charge $c = 3k/(k+2)>3$, so it is not the chiral algebra of a unitary four-dimensional theory
\emph{with that stress tensor} ($c_{\mathrm{4d}} = -c/12$ would be negative); the theorem is a statement about the
current sector, not a physical model.
\end{physics}

\begin{status}
\textbf{Proved, unconditional}: the only hypothesis is $k<-2$, $k$ real. \texttt{\#print axioms
universal\_gradedUnitary\_belowCritical\_unconditional} gives \lean{[propext, Classical.choice, Quot.sound]}
(isolated-module check; aggregate build not re-run for this section --- see the build-state note in
Section~\ref{l19:sec:axioms}). It does not use \lean{VOA.pbw}.
\end{status}

\begin{remark}[Scope]\label{l19:rem:bcscope}
The theorem concerns the universal algebra below the critical level. The simple quotients
$L_k(\lnineteensl)$ at admissible levels $k = -2+p/q>-2$ (for instance $k = -4/3$, $-8/5$), whose graded unitarity is
predicted by the four-dimensional correspondence, are not covered:
there the universal form is degenerate, the right flag is the corrected one, and the level is above $-2$. The
abstract no-wall lemma (Theorem~\ref{l19:thm:bcnowall}) is stated for an arbitrary VOA and would apply to such a
quotient, but the anchor-and-rigidity part of the argument has no analogue at an isolated level.
\end{remark}

%----------------------------------------------------------------------
\subsection{The three route theorems}
%----------------------------------------------------------------------

\paragraph{Order of the argument.} We first put the Gram matrices into coordinates (the next paragraph and
Definition~\ref{l19:def:bcinputs}), then record two properties of the Gram matrix that both route theorems use
(Lemma~\ref{l19:lem:bcgram}: Hermitian, polynomial in the level). Then come the anchor
(Theorem~\ref{l19:thm:bcanchor}, which uses the lemma), rigidity (Theorem~\ref{l19:thm:bcrigid}, which uses the lemma and
step nondegeneracy), the nondegeneracy of the full form (Theorem~\ref{l19:thm:bcnondeg}) and the no-wall lemma
(Theorem~\ref{l19:thm:bcnowall}), which together supply that step nondegeneracy. The headline's induction
(Theorem~\ref{l19:thm:bcheadline}) assembles them.

\paragraph{Gram matrices and layer signs in coordinates.} Fix $E$. Let $I_E$ be the finite set of PBW exponents
of energy $E$ (\lean{EIdx E}, independent of $k$), with length $\ell(i)$ (\lean{eLen E i}, the number of letters).
The Gram matrix is $G_E(k)_{ij} = \mathcal H(b_i,b_j) = B_k(\rho_kb_i,b_j)$ for the PBW monomials $b_i$ (\lean{gram}). For a
Hermitian matrix $G$ on $I_E$ and $r\in\N$, the \emph{step} $r$ is the coordinate subspace $\{\ell<r\}$, and the
matrix conditions are
\begin{align*}
  \textsc{StepNondeg}(G)&:\ \det\bigl(G|_{\{\ell<r\}}\bigr)\ne0\ \text{for every }r,\\
  \textsc{LayerSign}_r(G)&:\ (-1)^r\,y^*Gy>0\ \text{for every }y\ne0\text{ supported on }\{\ell\le r\}\\
  &\phantom{:\ }\text{with }(Gy)_i = 0\text{ for }\ell(i)<r.
\end{align*}

{\raggedright\noindent Files \leanpath{VOA/\allowbreak Examples/\allowbreak AffineSl2BelowCriticalMatrixRigidity.lean} and
\leanpath{VOA/\allowbreak Examples/\allowbreak AffineSl2BelowCriticalGramMatrix.lean}:\par}
\begin{leancode}
def stepMatrix (G : Matrix ι ι ℂ) (r : ℕ) : Matrix {i // ℓ i < r} {i // ℓ i < r} ℂ :=
  G.submatrix Subtype.val Subtype.val

def MatrixStepNondeg (G : Matrix ι ι ℂ) : Prop :=
  ∀ r : ℕ, (stepMatrix ℓ G r).det ≠ 0

def MatrixLayerSign (G : Matrix ι ι ℂ) (r : ℕ) : Prop :=
  ∀ y : ι → ℂ, (∀ i, r < ℓ i → y i = 0) → (∀ i, ℓ i < r → (G *ᵥ y) i = 0) → y ≠ 0 →
    0 < (-1 : ℂ) ^ r * (star y ⬝ᵥ (G *ᵥ y))

def EIdx (E : ℤ) : Type := {mu : CurrentIndex // currentEnergy mu = E}
...
def gram (R : RealStructure (universal k hk)) (E : ℤ) : Matrix (EIdx E) (EIdx E) ℂ :=
  fun i j => sesqForm (universalForm k hk) R (currentPBWBasis k hk i.val)
    (currentPBWBasis k hk j.val)
...
def eLen (E : ℤ) (i : EIdx E) : ℕ := pbwLength i.val

variable (hR : IsCompactRealStructure R)
include hR
...
theorem matrixStepNondeg_of_layerNondeg {E : ℤ}
    (hnd : LayerNondegenerate (ordinaryFlag k hk) E (LiRadical.gradedForm (universalForm k hk) E)) :
    MatrixStepNondeg (eLen E) (gram hk R E) := by
...
theorem layerPositive_iff_matrixLayerSign {E : ℤ} :
    LayerPositive (ordinaryFlag k hk) E (LiRadical.gradedForm (universalForm k hk) E)
        (gradedSesq (universalForm k hk) R E) ↔
      ∀ r : ℕ, MatrixLayerSign (eLen E) (gram hk R E) r := by
\end{leancode}
\noindent Here \lean{CurrentIndex} is a finitely supported multiplicity function on (colour, depth) pairs,
\lean{currentPBWBasis k hk mu} is the ordered creation word \lean{wordStateA k (orderedPBWColorWord mu)}
(\lean{currentPBWBasis_apply}; the PBW basis of Theorem~\ref{l08:thm:pbw}), \lean{currentEnergy} its weight and
\lean{pbwLength} its number of letters (\leanpath{VOA/\allowbreak Examples/\allowbreak AffineSl2JointPBW.lean},
\leanpath{AffineSl2BelowCriticalPBWFlag.lean}). In the matrix file \lean{ι} is any finite index type and \lean{ℓ}
any length function; \lean{star y ⬝ᵥ (G *ᵥ y)} is $y^*Gy$, and \lean{0 < …} in $\C$ (the library's
\lean{ComplexOrder}) means ``is a positive real number''. In the matrix file the section variables are
\lean{{ι : Type*} [Fintype ι] [DecidableEq ι] (ℓ : ι → ℕ)} (not displayed): $\iota$ is any finite index type with
decidable equality, which every finite type of interest has, so the extra instance is harmless.

\textbf{The two translation theorems and their hypotheses.} Both quoted Gram-matrix theorems lie after the
displayed \lean{variable (hR : IsCompactRealStructure R)} and \lean{include hR}, so \emph{both} take the hypothesis
that $R$ is a compact real structure (Definition~\ref{l19:def:bcdata}, item 3) as an explicit first argument; the
signature printed by \texttt{\#check} (at an earlier commit, not re-run for this edition; namespaces shortened) is
\begin{plaincode}
matrixStepNondeg_of_layerNondeg : ∀ {k : ℂ} (hk : k + 2 ≠ 0) {R : RealStructure (universal k hk)},
  IsCompactRealStructure R → ∀ {E : ℤ},
    LayerNondegenerate (ordinaryFlag k hk) E (LiRadical.gradedForm (universalForm k hk) E) →
      MatrixStepNondeg (eLen E) (gram hk R E)
\end{plaincode}
and likewise \lean{layerPositive_iff_matrixLayerSign} takes \lean{hR} before \lean{{E}}. Mathematically:
\begin{itemize}
\item (\lean{matrixStepNondeg_of_layerNondeg}) if $R$ is compact and $B_k$ is nondegenerate on every flag step
  $F_p\cap V^k_E$ (clause (iv) at weight $E$), then every step minor of $G_E(k)$ is nonzero:
  $\textsc{StepNondeg}(G_E(k))$;
\item (\lean{layerPositive_iff_matrixLayerSign}) if $R$ is compact, clause (vi) at weight $E$ holds if and only if
  $\textsc{LayerSign}_r(G_E(k))$ holds for every $r$.
\end{itemize}
\emph{Proof on paper.} By filtered PBW straightening (\lean{ordinaryPiece_eq_pbwSpan}) the flag piece $F_p\cap V^k_E$
is the coordinate subspace spanned by the ordered monomials of length $\le p$, so ``supported on $\{\ell\le p\}$'' is
membership in $F_p$, and ``$(Gy)_i = 0$ for $\ell(i)<r$'' is $\mathcal H$-orthogonality to $F_{r-1}$. A compact $\rho$
preserves the flag pieces and has nonzero parity scalars, so $\mathcal H$-orthogonality and $B_k$-orthogonality to a flag
piece coincide (as in Theorem~\ref{l19:thm:orthsame}); hence the layer of Definition~\ref{l19:def:gu} is exactly the
set of vectors $y$ in the definition of $\textsc{LayerSign}_r$, and $(-1)^r \mathcal H(y,y) = (-1)^r y^*G y$. For the step
minors: the matrix of $\mathcal H$ on the coordinate subspace $\{\ell<r\}$ is $G|_{\{\ell<r\}}$, and $\mathcal H(x,y) = B_k(\rho x,y)$
with $\rho$ a bijection of $F_{r-1}\cap V^k_E$, so nondegeneracy of $B_k$ there is invertibility of that block.
Nothing beyond linear algebra in an adapted basis is involved. \qed

\begin{physics}
These two statements say that the abstract signature conditions can be read off from the matrix of two-point
functions $\langle b_i|b_j\rangle$ of current states in the PBW basis, organised by current number: the familiar way a
physicist checks unitarity of a Verma module by computing Gram (Kac) matrices level by level. No physics enters
beyond this choice of basis.
\end{physics}

\begin{status}
Both proved, conditional only on their displayed hypotheses (\lean{hR} and, for the first, clause (iv) at weight
$E$); \texttt{\#print axioms}: \lean{[propext, Classical.choice, Quot.sound]} for both (Section~\ref{l19:sec:axioms}).
In the headline both hypotheses are discharged: \lean{hR} by \lean{compactRealFamily_isCompact}, clause (iv) by the
no-wall lemma (Theorem~\ref{l19:thm:bcnowall}).
\end{status}

\begin{definition}[Layer signs, step nondegeneracy, and the two analytic inputs]\label{l19:def:bcinputs}
For a family $R$ of real structures on $V^x$, $x<-2$, write $\mathrm{Sign}_R(x,E)$ for clause (vi) at weight $E$
and $\mathrm{Nd}(x,E)$ for clause (iv) at weight $E$, both for the ordinary flag and the form $B_x$.
\begin{enumerate}
\item The \emph{anchor} for $R$: for every $E$ there is $K\in\lnineteenR$ such that $\mathrm{Sign}_R(x,E)$ holds for
  every $x<-2$ with $x<K$.
\item \emph{Rigidity} for $R$: for every $E$, if $\mathrm{Nd}(x,E)$ holds for every $x<-2$, then
  $\mathrm{Sign}_R(x_0,E)$ for one $x_0<-2$ implies $\mathrm{Sign}_R(x_1,E)$ for every $x_1<-2$.
\end{enumerate}
\end{definition}

{\raggedright\noindent File \leanpath{VOA/\allowbreak Examples/\allowbreak AffineSl2BelowCriticalGradedUnitarity.lean}.\par}
\begin{leancode}
def LayerSigns (R : RealFamily) (x : ℝ) (hx : x < -2) (E : ℤ) : Prop :=
  LayerPositive (ordinaryFlag (x : ℂ) (noncritical_of_lt hx)) E
    (LiRadical.gradedForm (universalForm (x : ℂ) (noncritical_of_lt hx)) E)
    (gradedSesq (universalForm (x : ℂ) (noncritical_of_lt hx)) (R x hx) E)

def LayerNondeg (x : ℝ) (hx : x < -2) (E : ℤ) : Prop :=
  LayerNondegenerate (ordinaryFlag (x : ℂ) (noncritical_of_lt hx)) E
    (LiRadical.gradedForm (universalForm (x : ℂ) (noncritical_of_lt hx)) E)

def LayerSignAnchor (R : RealFamily) : Prop :=
  ∀ E : ℤ, ∃ K : ℝ, ∀ (x : ℝ) (hx : x < -2), x < K → LayerSigns R x hx E

def LayerSignRigidity (R : RealFamily) : Prop :=
  ∀ E : ℤ, (∀ (x : ℝ) (hx : x < -2), LayerNondeg x hx E) →
    ∀ (x₀ : ℝ) (hx₀ : x₀ < -2) (x₁ : ℝ) (hx₁ : x₁ < -2),
      LayerSigns R x₀ hx₀ E → LayerSigns R x₁ hx₁ E
\end{leancode}

\begin{dictionary}
\lean{LayerSigns R x hx E} is $\mathrm{Sign}_R(x,E)$ and \lean{LayerNondeg x hx E} is $\mathrm{Nd}(x,E)$, literally the
two halves of \lean{GradedUnitaryAt} at weight $E$ (Definition~\ref{l19:def:gu}) for the data of
Definition~\ref{l19:def:bcdata}, with the real structure taken from the family. \lean{LayerSignAnchor} and
\lean{LayerSignRigidity} are items 1 and 2, with the quantifiers in the displayed order: in the anchor $K$ may
depend on $E$ (the proof produces an $E$-dependent threshold; see below); in rigidity the nondegeneracy hypothesis is at \emph{all} levels $x<-2$. Both
are \lean{Prop}s about a family; as hypotheses they are \emph{inputs} of the abstract assembly, and the docstrings
call them ``named open input''. For the compact family both are now \emph{theorems}
(Theorems~\ref{l19:thm:bcanchor} and~\ref{l19:thm:bcrigid}); see Warning~\ref{l19:warn:bcstale}.

\textbf{Why is $K$ allowed to depend on $E$?} The weight-$E$ Gram matrices have polynomial entries of degree up to
$E$ in $k$. The anchor's \emph{proof} is a large-$|k|$ asymptotic argument (Theorem~\ref{l19:thm:bcanchor}), whose
threshold is extracted from the weight-$E$ Gram matrix and so depends on $E$; no uniformity in $E$ is claimed, and the
input is stated as ``$\forall E\,\exists K$''. This is a feature of the proof, not of the truth:
a posteriori the headline theorem (Theorem~\ref{l19:thm:bcheadline}) shows that $K=-2$ works for every $E$.
\end{dictionary}

\begin{physics}
$\mathrm{Sign}(x,E)$ says ``all states of conformal weight $E$ and current number exactly $p$ have the four-dimensional
sign $(+)$''; $\mathrm{Nd}(x,E)$ says ``no null states appear in any truncated current-number sector at weight $E$''.
The anchor is the free (large flavour central charge) limit; rigidity says signs cannot change without a null state
appearing on the way.
\end{physics}

\begin{status}
Definitions. The conditional assembly \lean{universal_gradedUnitary_belowCritical} (same file) proves
\lean{GradedUnitary} for any family \lean{R} with \lean{hR : ∀ x hx, IsCompactRealStructure (R x hx)},
\lean{hanchor : LayerSignAnchor R}, \lean{hrigid : LayerSignRigidity R}; it is \texttt{\#print axioms}-clean and,
as stated, conditional on those three hypotheses. For \lean{compactRealFamily} all three are discharged.
\end{status}

\begin{warning}[title={Stale ``open input'' language in \texttt{AffineSl2BelowCriticalGradedUnitarity.lean}}]\label{l19:warn:bcstale}
The module docstring of \leanpath{AffineSl2BelowCriticalGradedUnitarity.lean} says the assembly is conditional on ``three
named inputs, none of which is established in this repository'', and the docstrings of \lean{LayerSignAnchor} and
\lean{LayerSignRigidity} call them ``Named open input''. The next paragraph of the same module docstring says all three
are discharged downstream for the compact family, which is what the machine confirms
(\lean{compactRealFamily_isCompact}, \lean{layerSignRigidity_compact}, \lean{layerSignAnchor_compact}). The definitions
are correct; only the wording is stale. (Documentation only.)
\end{warning}

\begin{lemma}[The compact Gram matrix is Hermitian and polynomial in the level]\label{l19:lem:bcgram}
Let $k$ be real, $k\ne-2$, and $\rho_k$ the compact real structure. For every weight $E$: (1) $G_E(k)$ is Hermitian (indeed real
symmetric); (2) if moreover $k\ne0$, each entry is $G_E(k)_{ij} = P_{ij}(k)$ for a polynomial $P_{ij}\in\C[X]$ that does not depend
on $k$ (\lean{gramPoly E i j}).
\end{lemma}
\noindent The Lean statements are \lean{gram_isHermitian} and \lean{gram_eq_eval}, quoted with
Theorem~\ref{l19:thm:bcrigid} (file \leanpath{AffineSl2BelowCriticalRigidity.lean}).

\paragraph{Proof on paper.} (1) From \eqref{l19:eq:bcwick} below, $\overline{\mathcal H(b_i,b_j)} = \mathcal H(b_j,b_i)$ follows from two
symmetries of $B_k$, $B_k(\theta u,\theta v) = B_k(u,v)$ and $\overline{B_k(\beta u,\beta v)} = B_k(u,v)$ ($\theta$, $\beta$ as in
Definition~\ref{l19:def:bcdata}), both proved by the same uniqueness principle: a bilinear form for which every $J^a_m$ is
skew-adjoint to $J^a_{-m}$ is determined by its values $B(\lnineteenvac,\cdot)$ (\lean{bilinForm_eq_of_current_adj}), and both
sides of each symmetry have this property and agree on the vacuum. (2) Computing \eqref{l19:eq:bcwick} by moving annihilators to
the right with \eqref{l19:eq:bccomm} produces only integer combinations of powers of $k$; the structure constants are
integers and the level enters only through the central term. \qed

\begin{physics}
Hermiticity is the statement that the two-point functions of the $\mathfrak{su}(2)$ current states are those of a
(possibly indefinite) inner-product space; polynomial dependence on the level is the familiar fact that every correlator
of currents is a polynomial in the central charge $k$ of the current algebra.
\end{physics}

\begin{status}
Proved, unconditional (for real noncritical $k$, and $k\ne0$ for (2)); \texttt{\#print axioms} on \lean{gram_isHermitian}
and \lean{gram_eq_eval}: \lean{[propext, Classical.choice, Quot.sound]} (Section~\ref{l19:sec:axioms}).
\end{status}

\begin{theorem}[Anchor: layer signs at large negative level]\label{l19:thm:bcanchor}
\textbf{(a) Matrix version.} Let $I$ be a finite set with a length $\ell:I\to\N$, $P$ an $I\times I$ matrix of
complex polynomials, and $d:I\to\lnineteenR_{>0}$, such that
\begin{enumerate}
\item $\deg P_{ij}\le\min(\ell(i),\ell(j))$;
\item if $\ell(i) = \ell(j)$, the coefficient of $X^{\ell(i)}$ in $P_{ij}$ is $\delta_{ij}d_i$;
\item $P(x)$ is Hermitian for all real $x<x_*$ (some $x_*\in\lnineteenR$; \lean{x₀} in the Lean statement).
\end{enumerate}
Then there is $K$ such that for every real $x<K$ and every $r$, $\textsc{LayerSign}_r(P(x))$ holds.

\textbf{(b) Affine version.} For the compact family, the anchor of Definition~\ref{l19:def:bcinputs} holds: for every
weight $E$ there is $K$ such that the weight-$E$ layer signs hold at every real level $x<\min(K,-2)$.
\end{theorem}

{\raggedright\noindent Files \leanpath{VOA/\allowbreak Examples/\allowbreak AffineSl2BelowCriticalMatrixAnchor.lean} (a) and
\leanpath{VOA/\allowbreak Examples/\allowbreak AffineSl2BelowCriticalAnchor.lean} (b, and the two polynomial facts).\par}
\begin{leancode}
theorem matrixLayerSign_eventually (P : Matrix ι ι ℂ[X]) (d : ι → ℝ) (hd : ∀ i, 0 < d i)
    (hdeg : ∀ i j, (P i j).natDegree ≤ min (ℓ i) (ℓ j))
    (htop : ∀ i j, ℓ i = ℓ j → (P i j).coeff (ℓ i) = if i = j then (d i : ℂ) else 0)
    (x₀ : ℝ) (hherm : ∀ x : ℝ, x < x₀ → (P.map (Polynomial.eval (x : ℂ))).IsHermitian) :
    ∃ K : ℝ, ∀ x : ℝ, x < K → ∀ r, MatrixLayerSign ℓ (P.map (Polynomial.eval (x : ℂ))) r := by

theorem gramPoly_natDegree_le {E : ℤ} (i j : EIdx E) :
    (gramPoly E i j).natDegree ≤ min (eLen E i) (eLen E j) := by
...
theorem gramPoly_coeff_top {E : ℤ} (i j : EIdx E) (hlen : eLen E i = eLen E j) :
    (gramPoly E i j).coeff (eLen E i) = if i = j then ((topWeight i : ℝ) : ℂ) else 0 := by
...
theorem layerSignAnchor_compact : LayerSignAnchor compactRealFamily := by
\end{leancode}

\begin{dictionary}
\textbf{(a).} \lean{ℂ[X]} is $\C[X]$; \lean{P.map (Polynomial.eval (x : ℂ))} is the numerical matrix $P(x)$;
\lean{natDegree} is the degree (with $\deg 0 = 0$, harmless for an upper bound); \lean{coeff (ℓ i)} the coefficient of
$X^{\ell(i)}$. Hypotheses \lean{hd}, \lean{hdeg}, \lean{htop}, \lean{hherm} are items $d>0$, 1, 2, 3; the conclusion
is as stated, uniform in $r$. Nothing VOA-specific.

\textbf{(b).} \lean{gramPoly E i j} (\leanpath{AffineSl2BelowCriticalRigidity.lean}) is the \emph{level polynomial}
of the Gram entry: up to the sign $(-1)^{\ell(i)}$ it is \lean{wordGramPolynomial} of Chapter~\ref{ch:L09} applied
to the colour-reversed word of $b_i$ and the word of $b_j$, and \lean{gram_eq_eval} proves
$G_E(k)_{ij} = \mathtt{gramPoly}_{ij}(k)$ for real $k\ne0,-2$. \lean{topWeight i} is an explicit positive real
(\lean{realDiag}, a product of depths, Killing coefficients and multiplicities). So
\lean{gramPoly_natDegree_le} and \lean{gramPoly_coeff_top} are hypotheses 1 and 2 of (a) for $P = $
\lean{gramPoly E}, and hypothesis 3 is \lean{gram_isHermitian} below; \lean{layerSignAnchor_compact} is (b),
i.e.\ \lean{LayerSignAnchor compactRealFamily} of Definition~\ref{l19:def:bcinputs}, with no hypotheses.
\end{dictionary}

\paragraph{Proof on paper.} \emph{(a).} Put $x = -\mu^2$ with $\mu>0$ and rescale:
$M(\mu)_{ij} = P_{ij}(-\mu^2)/\mu^{\ell(i)+\ell(j)}$. By item~1, $P_{ij}(-\mu^2) = O(\mu^{2\min(\ell_i,\ell_j)})$, so
$M_{ij}\to0$ when $\ell_i\ne\ell_j$ (the exponent $2\min-\ell_i-\ell_j$ is negative). When $\ell_i = \ell_j = \ell$ the
leading term is the coefficient of $X^\ell$ times $(-\mu^2)^\ell$, so by item~2 $M_{ij}\to(-1)^\ell\delta_{ij}d_i$.
Hence $M(\mu)\to D = \operatorname{diag}((-1)^{\ell_i}d_i)$ as $\mu\to\infty$. For the diagonal matrix $D$ the step
minors are nonzero and the Schur complement of step $\{\ell<r\}$ in $\{\ell\le r\}$ is $\operatorname{diag}((-1)^rd_i)_{\ell_i = r}$, so
$(-1)^r$ times it is positive definite. Now (i) $\textsc{LayerSign}_r(G)$ is equivalent to positive definiteness of
$(-1)^r$ times that Schur complement (\lean{matrixLayerSign_iff_schur}); (ii) Schur complements depend continuously on
the matrix where the step minor is nonzero; (iii) positive definiteness is an open condition among Hermitian
matrices (\lean{posDef_eventually}: positivity on the compact unit sphere persists under small perturbation). So for
$\mu$ large every layer sign of $M(\mu)$ holds; for $r$ beyond the maximal length the condition is vacuous. Finally
$P(-\mu^2) = SM(\mu)S$ with $S = \operatorname{diag}(\mu^{\ell_i})$ positive diagonal, and congruence by a positive
diagonal matrix preserves supports, the orthogonality conditions and the signs (\lean{matrixLayerSign_congr}). Take
$K = -\mu_0^2$.

\emph{(b), the two polynomial facts.} Write $b_i = J^{a_1}_{-m_1}\cdots J^{a_\ell}_{-m_\ell}\lnineteenvac$. Moving the creators of $b_i$ across the form one at a
time with $\mathcal H(J^a_mx,y) = \mathcal H(x,J^{a^*}_{-m}y)$ (\lean{sesq_J_adj}), and using $\mathcal H(\lnineteenvac,z) = B_k(\lnineteenvac,z)$, the
vacuum coefficient of $z$,
\begin{equation}\label{l19:eq:bcwick}
  \mathcal H(b_i,b_j) = \bigl\langle J^{a_\ell^*}_{m_\ell}\cdots J^{a_1^*}_{m_1}\,b_j\bigr\rangle_{\lnineteenvac},
\end{equation}
the vacuum coefficient obtained by letting the annihilators act on $b_j$. Move each annihilator to the right with
\eqref{l19:eq:bccomm}. A factor of $k$ is produced only by a central term, which consumes one annihilator and one
creator; the bracket terms consume one of each and produce a mode of lower absolute value but no $k$. Hence the
degree in $k$ is at most the number of annihilators and at most the number of creators: item~1. At equal lengths
$\ell$, the coefficient of $k^\ell$ comes from complete matchings of each annihilator $J^{a^*}_m$ with a creator
$J^{a}_{-m}$ of the \emph{same} depth and the reversed colour, each contributing $m\,\kappa(a^*,a)>0$; such a
matching exists only if the two multisets of letters coincide, i.e.\ $i = j$ for ordered monomials, and then the
number of matchings is a product of factorials of multiplicities. So the top coefficient is $\delta_{ij}d_i$ with
$d_i>0$: item~2. (The sign: with $k = -\mu^2$ the top term $d_ik^{\ell}$ is $(-1)^{\ell}d_i\mu^{2\ell}$, and after the rescaling
of (a) this is the diagonal entry $(-1)^{\ell}d_i$ of the limit $D$; so the alternating signs required by graded unitarity
come entirely from $k<0$, while $d_i>0$ comes from $\kappa(a^*,a)>0$.) Item~3 is Lemma~\ref{l19:lem:bcgram}. Apply (a). \qed

\begin{physics}
Rescaling a state of $r$ currents by $|k|^{-r/2}$ is the usual normalisation of a current with two-point function
$\propto k$; as $|k|\to\infty$ the currents become free (generalized free fields) and every correlator is computed by
Wick contractions. In that limit a state of $r$ currents has norm $\propto k^r$, of sign $(-1)^r$ for $k<0$: exactly the
graded-unitarity sign. The anchor is the precise form of ``at large flavour central charge the current sector is a Fock
space''.
\end{physics}

\begin{status}
(a) and (b) proved; \texttt{\#print axioms} on \lean{matrixLayerSign_eventually}, \lean{gramPoly_natDegree_le},
\lean{gramPoly_coeff_top}, \lean{layerSignAnchor_compact}: \lean{[propext, Classical.choice, Quot.sound]}.
Unconditional.
\end{status}

\begin{theorem}[Rigidity of layer signs]\label{l19:thm:bcrigid}
\textbf{(a) Matrix version.} Let $X$ be a preconnected topological space and $H:X\to\mathrm{Mat}_I(\C)$ a continuous
family of Hermitian matrices with $\textsc{StepNondeg}(H(t))$ for every $t$. If $\textsc{LayerSign}_r(H(t_0))$ holds for one
$t_0$, it holds for every $t_1$.

\textbf{(b) Affine version.} Rigidity (Definition~\ref{l19:def:bcinputs}) holds for the compact family. Consequently, for
every real $k<-2$, graded unitarity of $V^k(\lnineteensl)$ with the compact structure and the ordinary flag follows from
the anchor alone.
\end{theorem}

{\raggedright\noindent Files \leanpath{VOA/\allowbreak Examples/\allowbreak AffineSl2BelowCriticalMatrixRigidity.lean} (a) and
\leanpath{VOA/\allowbreak Examples/\allowbreak AffineSl2BelowCriticalRigidity.lean} (b).\par}
\begin{leancode}
theorem matrixLayerSign_rigid {X : Type*} [TopologicalSpace X] [PreconnectedSpace X]
    (H : X → Matrix ι ι ℂ) (hc : Continuous H) (hH : ∀ t, (H t).IsHermitian)
    (hnd : ∀ t, MatrixStepNondeg ℓ (H t)) (r : ℕ) {t₀ t₁ : X}
    (h₀ : MatrixLayerSign ℓ (H t₀) r) : MatrixLayerSign ℓ (H t₁) r := by

theorem gram_isHermitian (E : ℤ) : (gram hk (compactRealStructure hreal hk) E).IsHermitian := by
...
theorem gram_eq_eval (hk0 : k ≠ 0) (E : ℤ) (i j : EIdx E) :
    gram hk (compactRealStructure hreal hk) E i j = (gramPoly E i j).eval k := by
...
theorem layerSignRigidity_compact : LayerSignRigidity compactRealFamily := by
...
theorem universal_gradedUnitary_belowCritical_of_anchor
    (hanchor : LayerSignAnchor compactRealFamily) (k : ℝ) (hk : k < -2) :
    GradedUnitary (universalForm (k : ℂ) (noncritical_of_lt hk)) (compactRealFamily k hk)
      (ordinaryFlag (k : ℂ) (noncritical_of_lt hk)) :=
  universal_gradedUnitary_belowCritical_compact hanchor layerSignRigidity_compact k hk
\end{leancode}

\begin{dictionary}
\textbf{(a).} \lean{PreconnectedSpace X}: $X$ cannot be split into two disjoint nonempty open sets (for us
$X = (-\infty,-2)$, \lean{Set.Iio (-2)}). \lean{Continuous H} is entrywise continuity. The hypotheses are exactly those
of (a); the conclusion is for the given $r$.

\textbf{(b).} \lean{gram_isHermitian} and \lean{gram_eq_eval} hold for every real noncritical $k$ (\lean{hreal}), the
latter for $k\ne0$: the compact Gram matrix is Hermitian (indeed real symmetric) and polynomial in the level, hence
continuous on $(-\infty,-2)$. \lean{layerSignRigidity_compact} is item~2 of Definition~\ref{l19:def:bcinputs} for
\lean{compactRealFamily}, no hypotheses. \lean{universal_gradedUnitary_belowCritical_of_anchor} is the headline with
the anchor still an argument; its one-line proof plugs the proved rigidity into the assembly
\lean{universal_gradedUnitary_belowCritical_compact} (\leanpath{AffineSl2BelowCriticalRealStructure.lean}), which in turn
specialises the abstract \lean{universal_gradedUnitary_belowCritical} to the compact family.
\end{dictionary}

\paragraph{Proof on paper.} \emph{(a).} At each $t$ the step matrix $\{\ell<r+1\} = \{\ell\le r\}$ is invertible, and
$\det H|_{\{\ell\le r\}} = \det H|_{\{\ell<r\}}\cdot\det S_r(t)$ where $S_r(t)$ is the Schur complement of the first in
the second; so $S_r(t)$ is a Hermitian matrix with trivial kernel, depending continuously on $t$ (the inverse of
$H|_{\{\ell<r\}}$ is the adjugate divided by a nonvanishing determinant). The set $U = \{t : (-1)^rS_r(t)\text{ positive
definite}\}$ is open (positivity on the compact unit sphere persists) and closed (a limit of positive definite matrices
is positive semidefinite, and positive semidefinite with nonzero determinant is positive definite;
\lean{posDef_rigid}). $X$ is preconnected and $t_0\in U$, so $U = X$. Translate back with
\lean{matrixLayerSign_iff_schur}.

\emph{(b).} The Gram matrix is Hermitian and polynomial in the level (Lemma~\ref{l19:lem:bcgram}), hence
continuous on $(-\infty,-2)$. The step-nondegeneracy hypothesis of (a) is the translation of
$\mathrm{Nd}(x,E)$ (\lean{matrixStepNondeg_of_layerNondeg}). Apply (a) on $X = (-\infty,-2)$, which is connected.
\qed

\begin{physics}
This is the familiar statement that a norm can change sign only by passing through zero: along a continuous family of
theories, a state can turn from positive to negative norm only at a point where some state becomes null. Below the
critical level there is no such point (no ``wall''), so the signs found in the free limit persist.
\end{physics}

\begin{status}
(a), (b) proved; \texttt{\#print axioms} on \lean{matrixLayerSign_rigid}, \lean{layerSignRigidity_compact},
\lean{universal_gradedUnitary_belowCritical_of_anchor}: \lean{[propext, Classical.choice, Quot.sound]}. The last is,
as stated, conditional on \lean{hanchor}; the hypothesis is discharged by Theorem~\ref{l19:thm:bcanchor} in the
headline.
\end{status}

\begin{theorem}[Nondegeneracy of the full form below critical level]\label{l19:thm:bcnondeg}
For every real $k<-2$ and every $E$, the form $B_k$ is nondegenerate on $V^k_E$: if $v\in V^k_E$ and $B_k(y,v) = 0$
for all $y\in V^k_E$, then $v = 0$.
\end{theorem}

{\raggedright\noindent File \leanpath{VOA/\allowbreak Examples/\allowbreak AffineSl2BelowCriticalNondegenerate.lean}.\par}
\begin{leancode}
theorem universalForm_separating_of_lt {x : ℝ} (hx : x < -2) (E : ℤ) :
    ∀ v ∈ (universal (x : ℂ) (noncritical_of_lt hx)).grading E,
      (∀ y ∈ (universal (x : ℂ) (noncritical_of_lt hx)).grading E,
        (universalForm (x : ℂ) (noncritical_of_lt hx)).form y v = 0) → v = 0 := by
\end{leancode}

\begin{dictionary}
The statement is separation in the second slot, which in finite dimension is nondegeneracy
(\lean{nondegenerate_of_separatingRight}). No hypothesis beyond $x<-2$. It is the input ``the full weight-$E$ form is
nondegenerate'' of the no-wall lemma.
\end{dictionary}

\paragraph{Proof on paper.} \emph{Step 1: blocks.} The zero mode $h_0 = J^h_0$ commutes with $L_0$ and is skew-adjoint for
$B_k$ in the sense $B_k(h_0u,v) = -B_k(u,h_0v)$; so weight vectors of different $h_0$-charges $c\ne -c'$ are $B_k$-orthogonal,
and the Gram matrix of $V^k_E$ in the PBW basis is block diagonal after pairing charge $c$ with charge $-c$ (the
``Chevalley twist'' $\theta$ of Definition~\ref{l19:def:bcdata} turns this into a block-diagonal matrix with blocks of
fixed charge, which is the form in which Chapters~\ref{ch:L08} and~\ref{ch:L09} compute it). It suffices to show that
each block has nonzero determinant.

\emph{Step 2: the Kac--Kazhdan formula.} For the Verma module of $\widehat{\lnineteensl}$ with highest weight $\Lambda$ (level
$k$, $h_0$-eigenvalue $\lambda$), the determinant of the Shapovalov form on the weight space
$\Lambda-\eta$ is, up to a nonzero constant, a product of factors
\begin{equation*}
  \bigl(2(\Lambda+\hat\rho,\beta)-n(\beta,\beta)\bigr)^{\,P(\eta-n\beta)\cdot\mathrm{mult}(\beta)},
\end{equation*}
one for each positive affine root $\beta$ and each integer $n\ge1$ ($\hat\rho$ the affine Weyl vector, $P$ the partition
function). The positive roots are of two kinds. \emph{Real} roots $\beta = \pm\alpha+m\delta$ ($m\ge1$, or $m = 0$ with the
sign $+$), with $(\beta,\beta) = 2$: since $(\Lambda,\alpha) = \lambda$, $(\hat\rho,\alpha) = 1$, $(\Lambda,\delta) = k$ and
$(\hat\rho,\delta) = h^\vee = 2$, the factor is $2\bigl(\pm(\lambda+1)+m(k+2)-n\bigr)$, which vanishes iff
$\pm(\lambda+1)+m(k+2) = n$; at $\lambda = 0$, iff $m(k+2)\pm1 = n$. \emph{Imaginary} roots $\beta = m\delta$, $m\ge1$, with $(\beta,\beta) = 0$: the
factor is $2(\Lambda+\hat\rho,m\delta) = 2m(k+h^\vee) = 2m(k+2)$, a power of $(k+2)$, vanishing only at the critical level.
In the library this is Theorem~\ref{l08:thm:kac}: the imaginary-root factors are the constant $C_{n,c}(t)$, nonzero for
$t = k+2\ne0$, and the real-root factors are the roots of the monic polynomial $K_{n,c}(t)$ in $\lambda$.

\emph{Step 3: from Verma to vacuum.} $V^k(\lnineteensl)$ is the quotient of the Verma module with $\lambda = 0$ by the submodule
generated by $f_0\lnineteenvac$, i.e.\ by the horizontal real root $\beta = \alpha$ ($m = 0$), whose factor ($n = 1$) vanishes
identically at $\lambda = 0$. The Gram determinant of the vacuum module is the Kac--Kazhdan product with this horizontal family
removed (Section~\ref{l09:sec:determinant}); what remains are the non-horizontal real factors $m\ge1$ and the imaginary
factors.

\emph{Step 4: no factor vanishes for $k<-2$.} A non-horizontal real factor vanishes only if $m(k+2)\pm1 = n$ for integers
$m\ge1$, $n\ge1$. For $k+2<0$ the left side is $m(k+2)+1<1$ in the case $+$ and $m(k+2)-1<-1$ in the case $-$, never a
positive integer. The imaginary factors are powers of $k+2\ne0$. So no factor vanishes
(\lean{vacuumOrder_eq_zero_of_lt}), every block is invertible, and the form is nondegenerate. \qed

\begin{physics}
Kac--Kazhdan zeros are the levels at which the vacuum module acquires null vectors (e.g.\ $(J^e_{-1})^{k+1}\lnineteenvac$ at
positive integer $k$, or the admissible singular vectors at $k = -2+p/q$). Below the critical level there are none:
$V^k(\lnineteensl)$ is simple for every real $k<-2$.
\end{physics}

\begin{status}
Proved, unconditional; \texttt{\#print axioms}: \lean{[propext, Classical.choice, Quot.sound]}.
\end{status}

\begin{theorem}[The prefix no-wall lemma]\label{l19:thm:bcnowall}
\textbf{(a) Abstract.} Let $V$ be a VOA with invariant form $B$, real structure $\rho$ and graded flag $F$, and let
$(X^i_m)_{i\in I,m\in\Z}$ ($I$ finite) be operators on $V$ with an involution $i\mapsto i^\dagger$ of $I$ and numbers
$\lambda_n,z_i\in\C$, $w_i>0$, such that: $X^i_m$ lowers the weight by $m$; $\mathcal H(X^i_ma,b) = \mathcal H(a,X^{i^\dagger}_{-m}b)$
for $\mathcal H(a,b) = B(\rho a,b)$; $X^i_m$ preserves $F_p$ for $m\ge0$ and maps $F_p$ into $F_{p+1}$ for $m<0$; every
$V_n$, $n>0$, is spanned by vectors $X^i_{-m}y$ with $m>0$, $y\in V_{n-m}$; $\rho$ preserves weights and flag pieces,
with $s(n)\ne0$; and the energy identity
\begin{equation}\label{l19:eq:bcsug}
  \lambda_n\,\mathcal H(v,v) = \sum_iz_i\,\mathcal H(X^i_0v,X^i_0v) + \sum_i\sum_{m=1}^{n}w_i\,\mathcal H(X^i_mv,X^i_mv)\qquad(v\in V_n)
\end{equation}
holds. Fix $E$. If clause (vi) holds at every weight $E-m$, $m\ge1$, and $B$ is nondegenerate on $V_E$, then $B$ is
nondegenerate on every flag step $F_p\cap V_E$ (clause (iv) at weight $E$).

\textbf{(b) Affine instance.} For $V^k(\lnineteensl)$ at any $k\ne-2$, the ordinary flag, any compact $\rho$, the
currents $X^a_m = J^a_m$, $a^\dagger = a^*$, and \eqref{l19:eq:bcsug} is the Hermitian Sugawara identity
\begin{equation*}
  2(k+2)E\,\mathcal H(v,v) = \sum_cs_c\,\mathcal H(J^c_0v,J^c_0v) + \sum_c\sum_{m=1}^{E}2s_c\,\mathcal H(J^c_mv,J^c_mv),\qquad
  (s_e,s_f,s_h) = (1,1,\tfrac12).
\end{equation*}
In the notation of (a) this is $\lambda_E = 2(k+2)E$, $z_c = s_c$ and $w_c = 2s_c>0$ (the Lean fields
\lean{lam}, \lean{zc}, \lean{wt} of \lean{currentSystem}, with $s_c$ = \lean{sugawaraWeight c}).
\end{theorem}

{\raggedright\noindent File \leanpath{VOA/\allowbreak Examples/\allowbreak AffineSl2BelowCriticalNoWall.lean} (a; the field
docstrings are omitted):\par}
\begin{leancode}
structure CurrentSystem (ι : Type*) [Fintype ι] (B : InvariantForm W) (R : RealStructure W)
    (F : GradedFlag W) where
  mode : ι → ℤ → Module.End ℂ 𝒱
  adj : ι → ι
  adj_adj : ∀ i, adj (adj i) = i
  mode_grading : ∀ (i : ι) (m : ℤ) {n : ℤ} {x : 𝒱}, x ∈ W.grading n →
    mode i m x ∈ W.grading (n - m)
  sesq_adj : ∀ (i : ι) (m : ℤ) (a b : 𝒱),
    sesqForm B R (mode i m a) b = sesqForm B R a (mode (adj i) (-m) b)
  nonneg_mem : ∀ (i : ι) {m : ℤ} {p : ℕ} {n : ℤ} {x : 𝒱}, 0 ≤ m → x ∈ F.piece p n →
    mode i m x ∈ F.piece p (n - m)
  neg_mem : ∀ (i : ι) {m : ℤ} {p : ℕ} {n : ℤ} {x : 𝒱}, 0 < m → x ∈ F.piece p n →
    mode i (-m) x ∈ F.piece (p + 1) (n + m)
  creation_span : ∀ {n : ℤ}, 0 < n → W.grading n ≤ Submodule.span ℂ
    {x | ∃ (i : ι) (m : ℕ) (y : 𝒱), 0 < m ∧ y ∈ W.grading (n - m) ∧ x = mode i (-(m : ℤ)) y}
  wt : ι → ℝ
  wt_pos : ∀ i, 0 < wt i
  lam : ℤ → ℂ
  zc : ι → ℂ
  sugawara : ∀ {n : ℤ} {v : 𝒱}, v ∈ W.grading n →
    lam n * sesqForm B R v v =
      (∑ i, zc i * sesqForm B R (mode i 0 v) (mode i 0 v)) +
        ∑ i, ∑ m ∈ Finset.Icc 1 n.toNat,
          (wt i : ℂ) * sesqForm B R (mode i (m : ℤ) v) (mode i (m : ℤ) v)
  rho_piece : ∀ {p : ℕ} {n : ℤ} {x : 𝒱}, x ∈ F.piece p n → R.rho x ∈ F.piece p n
  rho_grading : ∀ {n : ℤ} {x : 𝒱}, x ∈ W.grading n → R.rho x ∈ W.grading n
  parity_ne_zero : ∀ n, R.parity n ≠ 0
...
variable {ι : Type*} [Fintype ι] (S : CurrentSystem ι B R F)
...
theorem layerNondegenerate_of_lower (E : ℤ)
    (hfull : ∀ v ∈ W.grading E, (∀ x ∈ W.grading E, B.form x v = 0) → v = 0)
    (hlower : ∀ m : ℕ, 0 < m →
      LayerPositive F (E - m) (LiRadical.gradedForm B (E - m)) (gradedSesq B R (E - m))) :
    LayerNondegenerate F E (LiRadical.gradedForm B E) := by
\end{leancode}
{\raggedright\noindent File \leanpath{VOA/\allowbreak Examples/\allowbreak AffineSl2BelowCriticalCompact.lean} (b):\par}
\begin{leancode}
def sugawaraWeight : Fin 3 → ℝ := ![1, 1, 1 / 2]
...
theorem sugawara_sesq {E : ℤ} {v : WakModule k} (hv : v ∈ (universal k hk).grading E) :
    (2 * (k + 2) * (E : ℂ)) * sesqForm (universalForm k hk) R v v =
      (∑ c : Fin 3, (sugawaraWeight c : ℂ) *
          sesqForm (universalForm k hk) R (J k c 0 v) (J k c 0 v)) +
        ∑ c : Fin 3, ∑ m ∈ Finset.Icc 1 E.toNat, ((2 * sugawaraWeight c : ℝ) : ℂ) *
          sesqForm (universalForm k hk) R (J k c (m : ℤ) v) (J k c (m : ℤ) v) := by
...
def currentSystem {R : RealStructure (universal k hk)} (hR : IsCompactRealStructure R) :
    CurrentSystem (Fin 3) (universalForm k hk) R (ordinaryFlag k hk) where
  ...
\end{leancode}

\begin{dictionary}
\textbf{The structure.} \lean{CurrentSystem ι B R F} packages the hypotheses of (a) field by field: \lean{mode i m} is
$X^i_m$, \lean{adj} the involution $\dagger$, \lean{sesq_adj} the Hermitian adjointness, \lean{nonneg_mem}/\lean{neg_mem}
the flag behaviour, \lean{creation_span} the spanning, \lean{wt}/\lean{wt_pos} the positive weights $w_i$, \lean{lam} and
\lean{zc} the coefficients $\lambda_n$, $z_i$ (no sign condition on either), \lean{sugawara} identity
\eqref{l19:eq:bcsug} (\lean{n.toNat} is $\max(n,0)$; \lean{Finset.Icc 1 n.toNat} is $\{1,\dots,n\}$), and the last three
fields the conditions on $\rho$. A \lean{structure} with \lean{Prop} fields is only as good as its instances; here it is
instantiated by \lean{currentSystem hR} for $V^k(\lnineteensl)$, every field proved (\lean{lam E := 2 * (k + 2) * (E : ℂ)},
\lean{wt c := 2 * sugawaraWeight c}, \lean{sugawara hv := hR.sugawara_sesq hv}), so it is not vacuous.

\textbf{The theorem.} \lean{hfull} is separation of $B$ on $V_E$ (Theorem~\ref{l19:thm:bcnondeg} supplies it);
\lean{hlower} is clause (vi) at each lower weight $E-m$, $m\ge1$, \emph{at the same} $(B,\rho,F)$; the conclusion is
clause (iv) at weight $E$. It does \emph{not} conclude clause (vi) at weight $E$ (the docstring stresses this). In the
quoted source the theorem lives in \lean{namespace CurrentSystem} with the displayed \lean{variable} line, so it is
\lean{S.layerNondegenerate_of_lower}; the affine specialisation \lean{AffineSl2BelowCritical.layerNondegenerate_of_lower}
(\leanpath{AffineSl2BelowCriticalCompact.lean}) is
\lean{(currentSystem hR).}\allowbreak\lean{layerNondegenerate_of_lower}.

\textbf{The Sugawara identity.} \lean{sugawara_sesq} holds for every noncritical $k$ and every compact $\rho$; the
weights $(1,1,\tfrac12)$ are the entries of $\kappa^{-1}$ in the basis $(e,f,h)$. Note the coefficient $2(k+2)E$ is
\emph{negative} for $k<-2$; the argument does not need its sign.
\end{dictionary}

\paragraph{Proof on paper.} \emph{The Sugawara identity.} $L_0 = \frac{1}{2(k+2)}\sum_{a,b}\kappa^{ab}
\bigl(J^a_0J^b_0 + 2\sum_{m>0}J^a_{-m}J^b_m\bigr)$ on $V^k$, where $\kappa^{ab}$ is the inverse Killing matrix
($\kappa^{ef} = \kappa^{fe} = 1$, $\kappa^{hh} = \tfrac12$). Pair with $v\in V_E$ under $\mathcal H$, use $L_0v = Ev$ and move
$J^a_{-m}$ to the left as its adjoint $J^{a^*}_m$; since $\kappa^{ab}\ne0$ only for $b = a^*$, every term becomes a
norm $\mathcal H(J^c_mv,J^c_mv)$.

\emph{The lemma.} Suppose $v\in F_r\cap V_E$ is in the radical of $B$ restricted to $F_r\cap V_E$. Because
$\rho$ preserves $F_r$ and $s(E)\ne0$, $v$ is also $\mathcal H$-orthogonal to $F_r\cap V_E$ (as in Theorem~\ref{l19:thm:orthsame}).
\begin{enumerate}
\item For $m>0$, $u = X^i_mv$ lies in $F_r\cap V_{E-m}$, and for $z\in F_{r-1}\cap V_{E-m}$,
  $\mathcal H(z,u) = \mathcal H(X^{i^\dagger}_{-m}z,v) = 0$ since $X^{i^\dagger}_{-m}z\in F_r$. So $u$ lies in the layer $\lnineteenLay_r$ at
  weight $E-m$, where by hypothesis $(-1)^rh(u,u)\ge0$ with equality only for $u = 0$.
\item Zero modes preserve $F_r$, so $\mathcal H(X^i_0v,X^i_0v) = \mathcal H(v,X^{i^\dagger}_0X^i_0v) = 0$; and $\mathcal H(v,v) = 0$.
\item Multiply \eqref{l19:eq:bcsug} by $(-1)^r$: the left side and the zero-mode terms vanish, so
  $\sum_{i,m}w_i\,(-1)^r\,\mathcal H(X^i_mv,\,X^i_mv)$ $= 0$, a sum of nonnegative reals with $w_i>0$. Every term vanishes, so
  $X^i_mv = 0$ for all $i$ and $m>0$.
\item Then for every spanning vector $X^i_{-m}y$ of $V_E$, $\mathcal H(X^i_{-m}y,v) = \mathcal H(y,X^{i^\dagger}_mv) = 0$. So $v$ is
  orthogonal to all of $V_E$, and $v = 0$ by full nondegeneracy.
\end{enumerate}
(For $E\le0$ the argument degenerates harmlessly: $V_E = 0$ for $E<0$ and $F_r\cap V_0 = V_0$.) \qed

\begin{physics}
A ``wall'' is a level at which a truncated sector develops a null state, the only place where signs can flip. The lemma
says a null state in a current-number-truncated sector would have to be a highest-weight state for the currents (all
positive modes kill it, step 3) --- because its descendants under the positive modes live in lower sectors whose norms all
have the same sign, and the Sugawara identity (which expresses $L_0$ as a sum of norms) forces them to cancel. A
current-primary null state at weight $E>0$ is a singular vector of the full vacuum module, and there are none below the
critical level. (Heuristic summary; the rigorous statement is Theorem~\ref{l19:thm:bcnondeg}.)
\end{physics}

\begin{status}
(a) and (b) proved; \texttt{\#print axioms} on \lean{CurrentSystem.layerNondegenerate_of_lower},
\lean{AffineSl2BelowCritical.layerNondegenerate_of_lower} and \lean{IsCompactRealStructure.sugawara_sesq}:
\lean{[propext, Classical.choice, Quot.sound]}. The abstract lemma (a) is conditional on its displayed hypotheses
\lean{S}, \lean{hfull}, \lean{hlower}, which in the headline are supplied by \lean{currentSystem},
Theorem~\ref{l19:thm:bcnondeg} and the induction hypothesis.
\end{status}

%----------------------------------------------------------------------
\subsection{A worked example: weights one and two}
%----------------------------------------------------------------------

\begin{example}[Weight one and weight two of $V^k(\lnineteensl)$]\label{l19:ex:bcweight2}
The following is a paper computation, \emph{not} a Lean theorem. Every entry is elementary and can be checked by hand
from \eqref{l19:eq:bccomm} and \eqref{l19:eq:bcwick}; the author also checked them with a short computer-algebra script
(not part of the repository), and the fidelity referee recomputed them independently. We write $e_{-n}$ for
$J^e_{-n}$ etc.

\emph{Weight one.} $V_1$ has basis $e_{-1}\lnineteenvac$, $f_{-1}\lnineteenvac$, $h_{-1}\lnineteenvac$, all of length one, and
step $\{\ell<1\}$ is empty at this weight, so $\lnineteenLay_1 = V_1$. By \eqref{l19:eq:bcwick},
$\mathcal H(e_{-1}\lnineteenvac,e_{-1}\lnineteenvac) = \langle f_1e_{-1}\rangle = \langle -h_0+k\rangle = k$, and similarly
\begin{equation*}
  G_1(k) = \operatorname{diag}(k,\,k,\,2k).
\end{equation*}
Clause (vi) at $p = 1$ asks $-G_1(k)$ positive definite: $k<0$. (For $k>0$ --- the regime whose simple quotients
$L_k(\lnineteensl)$, $k\in\Z_{>0}$, are the unitary WZW models --- the graded-unitarity sign is violated: the two notions of positivity are opposite at odd $R$, as in
Remark~\ref{l19:rem:incompatible}.)

\emph{Weight two.} $\dim V_2 = 9$: three monomials of length one ($e_{-2},f_{-2},h_{-2}$) and six of length two. The
$h_0$-charge splits $V_2$ into blocks of charge $4,2,0,-2,-4$. In the charge-$0$ block take the PBW monomials
\begin{equation*}
  v_1 = h_{-2}\lnineteenvac,\qquad v_2 = e_{-1}f_{-1}\lnineteenvac,\qquad v_3 = h_{-1}h_{-1}\lnineteenvac,
\end{equation*}
of lengths $1,2,2$. Their Gram matrix is
\begin{equation*}
  G = \begin{pmatrix} 4k & 2k & 0\\ 2k & k^2+2k & 4k\\ 0 & 4k & 8k^2\end{pmatrix}.
\end{equation*}
For instance $\mathcal H(v_2,v_2) = \langle e_1f_1e_{-1}f_{-1}\rangle$; first
$f_1e_{-1}f_{-1}\lnineteenvac = (-h_0+k)f_{-1}\lnineteenvac = (k+2)f_{-1}\lnineteenvac$, then $e_1(k+2)f_{-1}\lnineteenvac = k(k+2)\lnineteenvac$.
\begin{itemize}
\item \emph{Layer 1} is $\C v_1$, with $\mathcal H(v_1,v_1) = 4k$: sign $-$ for $k<0$, as required.
\item \emph{Layer 2} is $\{av_1+bv_2+cv_3 : \mathcal H(v_1,\cdot) = 0\} = \{4ka+2kb = 0\}$, spanned by
  $w_2 = v_2-\tfrac12v_1 = \tfrac12(e_{-1}f_{-1}+f_{-1}e_{-1})\lnineteenvac$ and $w_3 = v_3$. On it the form is the Schur complement
  \begin{equation*}
    S = \begin{pmatrix} k^2+k & 4k\\ 4k & 8k^2\end{pmatrix},\qquad \det S = 8k^2(k+2)(k-1).
  \end{equation*}
  For $k<-2$: $k^2+k = k(k+1)>0$ and $\det S>0$, so $S$ is positive definite --- sign $+$, as required. At $k = -2$,
  $\det S = 0$; for $-2<k<1$, $k\ne0$, $\det S<0$ and $S$ is indefinite, so clause (vi) \emph{fails}.
\end{itemize}
The other blocks (computed by the same rules; the details are not shown) give layer-2 Schur complements $2k(k-1)$
(charge $2$: basis $e_{-2}\lnineteenvac$, $e_{-1}h_{-1}\lnineteenvac$, the Schur complement of $e_{-2}\lnineteenvac$; charge $-2$
likewise with $f$) and
$\mathcal H(e_{-1}^2\lnineteenvac,e_{-1}^2\lnineteenvac) = 2k(k-1)$ (charges $\pm4$), positive for $k<0$; the full weight-two determinant is
$2^{11}k^9(k-1)^5(k+2)$, nonzero for $k<-2$ (clause (iv) and Theorem~\ref{l19:thm:bcnondeg}).

\emph{The Sugawara vector sits in layer 2.} The unnormalised Sugawara vector is
$\omega' = \tfrac12h_{-1}^2+e_{-1}f_{-1}+f_{-1}e_{-1} = 2w_2+\tfrac12w_3$ (applied to $\lnineteenvac$), and
$\mathcal H(\omega',\omega') = 4(k^2+k)+2\cdot4k+\tfrac14\cdot8k^2 = 6k(k+2)$. With $\omega = \omega'/(2(k+2))$,
\begin{equation*}
  \mathcal H(\omega,\omega) = \frac{3k}{2(k+2)} = \frac c2,
\end{equation*}
the familiar norm $\langle T|T\rangle = c/2$. It is positive exactly when $k<-2$ or $k>0$, and the wall at $k = -2$ is
the critical level, where $\omega'$ becomes null. Together with weight one ($k<0$) this shows that $k<-2$ is exactly the
range in which weights one and two satisfy the theorem.
\end{example}

\begin{physics}
The example is the whole theorem in miniature: one-current states have the four-dimensional sign through the
factor $k<0$, the two-current states (including the stress tensor, a composite of two currents in the universal
algebra) have the sign $+$, and the only obstruction is the critical level, where the Sugawara stress tensor
degenerates into the Feigin--Frenkel centre.
\end{physics}

%----------------------------------------------------------------------
\subsection{The remaining files}
%----------------------------------------------------------------------

The twelve files, collectively (each is unconditional except where a hypothesis is named):
\begingroup\raggedright
\begin{itemize}
\item \leanpath{AffineSl2BelowCriticalFlag.lean}: \lean{ordinaryPiece}, \lean{ordinaryFlag}; current modes $J^a_m$
  preserve $F_p$ for $m\ge0$ (\lean{J_mem_ordinaryPiece_of_nonneg}) and map $F_p$ into $F_{p+1}$ for $m<0$; positive weight
  spaces are spanned by negative-mode images (\lean{grading_le_span_creation}).
\item \leanpath{AffineSl2BelowCriticalPBWFlag.lean}: filtered PBW straightening --- a creation word of length $L$ is a
  combination of ordered PBW monomials of length $\le L$ (an adjacent swap keeps the length, its commutator term is one
  letter shorter); \lean{ordinaryPiece_eq_pbwSpan}, \lean{mem_ordinaryPiece_iff_repr}.
\item \leanpath{AffineSl2BelowCriticalCompact.lean}: \lean{IsCompactRealStructure} and its consequences
  (\lean{rho_J}, \lean{sesq_J_adj}, $\rho$ preserves weights and the flag), \lean{sugawara_sesq}, \lean{currentSystem},
  and the affine no-wall lemma.
\item \leanpath{AffineSl2BelowCriticalRealStructure.lean}: the conjugate vacuum module \lean{ConjMod}, \lean{basisConj},
  \lean{compactRho}, covariance of every vertex operator (\lean{compactRho_YA}) and of the Sugawara field,
  \lean{compactRealStructure}, \lean{compactRealFamily}, and the assembly
  \lean{universal_gradedUnitary_belowCritical_compact} (conditional on \lean{hanchor}, \lean{hrigid}).
\item \leanpath{AffineSl2BelowCriticalNondegenerate.lean}: \lean{noncritical_of_lt}, \lean{vacuumOrder_eq_zero_of_lt},
  \lean{vacuumChevalleyPolynomial_det_eval_ne_zero}, \lean{universalForm_separating_of_lt}.
\item \leanpath{AffineSl2BelowCriticalNoWall.lean}: \lean{CurrentSystem} and the abstract no-wall lemma, for an arbitrary
  VOA.
\item \leanpath{AffineSl2BelowCriticalGradedUnitarity.lean}: \lean{RealFamily}, \lean{LayerSigns}, \lean{LayerNondeg},
  \lean{LayerSignAnchor}, \lean{LayerSignRigidity}, the induction (\lean{step}, \lean{layerSigns_all}) and the abstract
  assembly \lean{universal_gradedUnitary_belowCritical} (conditional on \lean{hR}, \lean{hanchor}, \lean{hrigid}).
\item \leanpath{AffineSl2BelowCriticalGramMatrix.lean}: \lean{EIdx}, \lean{coords}, \lean{gram}; clause (iv) $\Rightarrow$
  nonzero step minors, clause (vi) $\iff$ matrix layer signs.
\item \leanpath{AffineSl2BelowCriticalMatrixRigidity.lean} and \leanpath{AffineSl2BelowCriticalMatrixAnchor.lean}: pure
  linear algebra and topology (step matrices, Schur complements, \lean{posDef_rigid}, \lean{posDef_eventually},
  \lean{matrixLayerSign_congr}, the rescaling limit \lean{tendsto_rescaled}).
\item \leanpath{AffineSl2BelowCriticalRigidity.lean}: \lean{bilinForm_eq_of_current_adj}, the two invariances of $B_k$,
  \lean{gram_isHermitian}, \lean{gramPoly}, \lean{gram_eq_eval}, \lean{layerSignRigidity_compact}.
\item \leanpath{AffineSl2BelowCriticalAnchor.lean}: \lean{gramPoly_symm}, \lean{gramPoly_natDegree_le},
  \lean{gramPoly_coeff_top}, \lean{layerSignAnchor_compact}, and the headline.
\end{itemize}
\endgroup



%======================================================================
\section{Appendix: geometric VOAs and the chiral de Rham complex}\label{l19:sec:geometric}
%======================================================================

The directory \leanpath{VOA/\allowbreak Geometric/} contains two files,
\leanpath{VOA/\allowbreak Geometric/\allowbreak GeometricVOA.lean} and \leanpath{VOA/\allowbreak Geometric/\allowbreak ChiralDeRham.lean}.
\textbf{Neither declares anything}: each consists of an import, a module docstring, and an empty
\lean{namespace VOA ... end VOA} with the comment ``OUT OF SCOPE''. The ``target declarations'' in their
docstrings (\lean{SphereWithTubes}, \lean{GeometricVOA}, \lean{geometricEquivAlgebraic}, \lean{VOSA},
\lean{chiralDeRham}, \lean{ellipticGenus}) live inside docstring code fences and are not Lean code. Nothing
about chiral de Rham complexes or elliptic genera is proved, stated or defined in the library; the
``exploratory'' label of the library map means exactly this.

\paragraph{The mathematics the files point to.} A \emph{geometric VOA} (Huang) is a meromorphic, sewing-compatible
assignment of linear maps $V^{\otimes n}\to\overline V$ to spheres with $n+1$ tubes (marked points with local
coordinates); the pair of pants with punctures at $0,w,\infty$ gives $Y(u,w)v$, sewing gives the Jacobi identity, and
the Virasoro algebra encodes changes of local coordinates. Huang proved that geometric VOAs are equivalent to
(algebraic) VOAs. The \emph{chiral de Rham complex} of Malikov--Schechtman--Vaintrob (1999) is a sheaf
$\Omega^{\mathrm{ch}}_X$ of vertex superalgebras on a smooth complex variety $X$, locally the $\beta\gamma$--$bc$ system on
coordinate charts, glued by explicit (and subtly corrected) transformation rules; it carries a differential and, for
Calabi--Yau $X$, an $N=2$ superconformal structure. Its cohomology has the de Rham cohomology of $X$ in conformal
weight zero, and its graded character $\operatorname{tr}q^{L_0}y^{J_0}$ (suitably signed) is the \emph{elliptic genus}
of $X$ (Borisov--Libgober 2000), a weak Jacobi form for Calabi--Yau $X$.

\begin{physics}
$\Omega^{\mathrm{ch}}_X$ is the perturbative chiral algebra of the $(0,2)$ (half-twisted) sigma model with target $X$
(Kapustin 2005; Witten 2007); the elliptic genus is the sigma model's partition function on the torus with the
$U(1)$ R-symmetry fugacity $y$, a protected quantity independent of the metric, which is why it is a topological
invariant and a modular (Jacobi) form. For a Calabi--Yau 3-fold it counts BPS states of the half-twisted theory.
\end{physics}

\begin{status}
Nothing to verify: the files contain no declarations. The file \leanpath{VOA/\allowbreak Examples/\allowbreak MoonshineModule.lean}
(assigned here by the catch-all rule) is likewise an empty stub whose docstring now correctly points to the actual
construction of $\Vnat$ in \leanpath{VOA/\allowbreak Moonshine/} (Chapters~\ref{ch:01}--\ref{ch:15}); its title line
(``next Phase-3 target'') and the comment \lean{-- TODO (Phase 3, next target)} are stale.
\end{status}

%======================================================================
\section{Axiom audit}\label{l19:sec:axioms}
%======================================================================

The following was produced for this chapter at an earlier commit (not re-run for this edition) by a scratch file importing the
relevant modules and running \texttt{\#print axioms}, compiled with
\texttt{lake env lean -DmaxSynthPendingDepth=3} (never with a bare \texttt{lake env lean}, which can silently
change elaboration; no full \texttt{lake build} was run for this chapter, the modules being already built).
Output, one declaration per line:

\begin{plaincode}
'VOA.eq_bot_of_bilinear_pos' depends on axioms: [propext, Classical.choice, Quot.sound]
'VOA.SesqForm.I_smul_self' depends on axioms: [propext, Classical.choice, Quot.sound]
'VOA.layerDecomposition' depends on axioms: [propext, Classical.choice, Quot.sound]
'VOA.finrank_pieceLT_add_layer' depends on axioms: [propext, Classical.choice, Quot.sound]
'VOA.orthogonal_gradedSesq_eq' depends on axioms: [propext, Classical.choice, Quot.sound]
'VOA.layer_zero_ne_bot_and_isPosOn' depends on axioms: [propext, Classical.choice, Quot.sound]
'VOA.not_isPosOn_signedForm_stdHerm_one' depends on axioms: [propext, Classical.choice, Quot.sound]
'VOA.FlagMinors.nondegenerate_flag_iff_minors' depends on axioms: [propext, Classical.choice, Quot.sound]
'VOA.FlagMinors.posDef_of_det_leadingPrincipal_pos' depends on axioms: [propext, Classical.choice, Quot.sound]
'VOA.FlagMinors.schurOfLeading_posDef_of_sign_ratios' depends on axioms: [propext, Classical.choice, Quot.sound]
'VOA.FlagMinors.schurOfLeading_negDef_of_sign_ratios' depends on axioms: [propext, Classical.choice, Quot.sound]
'VOA.FlagMinors.det_leadingPrincipal_unitriangular_congr' depends on axioms: [propext, Classical.choice, Quot.sound]
'VOA.FlagMinors.det_leadingPrincipal_eq_mul_det_schurOfLeading' depends on axioms: [propext, Classical.choice, Quot.sound]
'VOA.fromBlocks_posDef_iff_schur₁₁' depends on axioms: [propext, Classical.choice, Quot.sound]
'VOA.realSymmetricTwo_posDef_iff' depends on axioms: [propext, Classical.choice, Quot.sound]
'VOA.SesqForm.isPosOn_signedForm_of_eq_pullback' depends on axioms: [propext, Classical.choice, Quot.sound]
'VOA.VirGradedUnitarity.gradedUnitaryAt_iff_minors_and_positive' depends on axioms: [propext, Classical.choice, Quot.sound]
'VOA.VirGradedUnitarity.virLayerDecomposition' depends on axioms: [propext, Classical.choice, Quot.sound]
'VOA.VirGradedUnitarity.isPosOn_layer_zero_vacuum_simple' depends on axioms: [propext, Classical.choice, Quot.sound]
'VOA.virasoroSimpleRealStructure' depends on axioms: [propext, Classical.choice, Quot.sound]
'VOA.VirGradedUnitarity.boundaryVirasoroRealStructure' depends on axioms: [propext, Classical.choice, Quot.sound]
'VOA.VirGradedUnitarity.boundary_gradedUnitaryThroughTwo' depends on axioms: [propext, Classical.choice, Quot.sound]
'VOA.VirGradedUnitarity.boundary_gradedUnitaryThroughTwoAt' depends on axioms: [propext, Classical.choice, Quot.sound]
'VOA.VirGradedUnitarity.canB_lengthTwoQ_self' depends on axioms: [propext, Classical.choice, Quot.sound]
'VOA.VirGradedUnitarity.lengthTwoQ_quasiPrimary' depends on axioms: [propext, Classical.choice, Quot.sound]
'VOA.VirGradedUnitarity.lengthTwoRho_mem_Ioc' depends on axioms: [propext, Classical.choice, Quot.sound]
'VOA.VirGradedUnitarity.rhoStableOn_vir_pieceIn' depends on axioms: [propext, Classical.choice, Quot.sound]
'VOA.LeeYangC2.leeYang_stronglyFinite' depends on axioms: [propext, Classical.choice, Quot.sound]
'VOA.LeeYangC2.leeYang_c2Cofinite' depends on axioms: [propext, Classical.choice, Quot.sound]
'VOA.LeeYangC2.finrank_a₂_leeYang_le_two' depends on axioms: [propext, Classical.choice, Quot.sound]
'VOA.LeeYangZhu.leeYang_zhu_polynomial' depends on axioms: [propext, Classical.choice, Quot.sound]
'VOA.LeeYangZhu.lowestWeight_eq_zero_or_neg_one_fifth' depends on axioms: [propext, Classical.choice, Quot.sound]
'VOA.LeeYang.exists_level4_singular_iff' depends on axioms: [propext, Classical.choice, Quot.sound]
'VOA.LeeYang.virasoroVOA_leeYang_not_isSimple' depends on axioms: [propext, Classical.choice, Quot.sound]
'VOA.LeeYangKac.det_gram4' depends on axioms: [propext, Classical.choice, Quot.sound]
'VOA.LeeYangKac.maxIdealVir_leeYang_inf_grading_four' depends on axioms: [propext, Classical.choice, Quot.sound]
'VOA.LeeYangRiccati.tower_posDef' depends on axioms: [propext, Classical.choice, Quot.sound]
'VOA.InverseDS.GradedIsometricEmbedding.gradedUnitary_of_ambient' depends on axioms: [propext, Classical.choice, Quot.sound]
'VOA.DSsl2.dsCentralCharge_eq_minimalModel' depends on axioms: [propext, Classical.choice, Quot.sound]
'VOA.DSsl2.dsCentralCharge_eq_affine_add_ghost' depends on axioms: [propext, Classical.choice, Quot.sound]
'VOA.DSsl2.dsQ_mul_self' depends on axioms: [propext, Classical.choice, Quot.sound]
'VOA.DSMor.dsLatticeScreening_centralCharge' depends on axioms: [propext, Classical.choice, Quot.sound]
'VOA.DSMor.dsLatticeScreening_ne_top' depends on axioms: [propext, Classical.choice, Quot.sound]
'VOA.DSMor.dsScreeningIncl_isVOAHom' depends on axioms: [propext, Classical.choice, Quot.sound]
'VOA.DSMor.dsComplex' depends on axioms: [propext, Classical.choice, Quot.sound]
'VOA.DSMor.dsEmbed_injective' depends on axioms: [propext, Classical.choice, Quot.sound]
'VOA.AffineVOA.GradedUnitarityLow.compactCurrentSignedForm85_eq_canonical' depends on axioms: [propext, Classical.choice, Quot.sound]
'VOA.AffineVOA.GradedUnitarityLow.compactCurrentSignedForm85_positive' depends on axioms: [propext, Classical.choice, Quot.sound]
'VOA.AffineVOA.GradedUnitarityLow.canonicalForm85_currentState' depends on axioms: [propext, Classical.choice, Quot.sound]
'VOA.PolarizedRees.MinorPayload.sign_det_leadingPrincipal_pos' depends on axioms: [propext, Classical.choice, Quot.sound]
'VOA.WeilParity.weil_polarized_twist_option_a' depends on axioms: [propext, Classical.choice, Quot.sound]
'VOA.LengthThreeQPFrame.negDef_w12' depends on axioms: [propext, Classical.choice, Quot.sound]
'VOA.LeeYangC2.leeYang_stronglyRational_of_admissiblyRational' depends on axioms: [propext, Classical.choice, Quot.sound]
'VOA.VirGradedUnitarity.Bsimp_boundaryLengthLETwoTowerAt_self_sign' depends on axioms: [propext, Classical.choice, Quot.sound]
'VOA.VirGradedUnitarity.toMatrix_eq_canB_of_monomials' depends on axioms: [propext, Classical.choice, Quot.sound]
'VOA.LengthThreeQPFrame.a3count_eq_pLe3_shift' does not depend on any axioms
'VOA.InverseDS.GradedIsometricEmbedding.gradedUnitaryAt_of_ambient' depends on axioms: [propext, Classical.choice, Quot.sound]
'VOA.LeeYang.lengthTwoQSimple_two_eq_zero' depends on axioms: [propext, Classical.choice, Quot.sound]
\end{plaincode}

\noindent and, for Section~\ref{l19:sec:belowcritical}, at a later commit (also not re-run for this edition; scratch file importing
\leanpath{VOA/\allowbreak Examples/\allowbreak AffineSl2BelowCriticalAnchor.lean}, same command; line breaks of the
output joined):

\begin{plaincode}
'VOA.AffineSl2BelowCritical.universal_gradedUnitary_belowCritical_unconditional' depends on axioms: [propext, Classical.choice, Quot.sound]
'VOA.AffineSl2BelowCritical.layerSignAnchor_compact' depends on axioms: [propext, Classical.choice, Quot.sound]
'VOA.AffineSl2BelowCritical.layerSignRigidity_compact' depends on axioms: [propext, Classical.choice, Quot.sound]
'VOA.AffineSl2BelowCritical.universal_gradedUnitary_belowCritical' depends on axioms: [propext, Classical.choice, Quot.sound]
'VOA.AffineSl2BelowCritical.universal_gradedUnitary_belowCritical_of_anchor' depends on axioms: [propext, Classical.choice, Quot.sound]
'VOA.AffineSl2BelowCritical.universal_gradedUnitary_belowCritical_compact' depends on axioms: [propext, Classical.choice, Quot.sound]
'VOA.AffineSl2BelowCritical.universalForm_separating_of_lt' depends on axioms: [propext, Classical.choice, Quot.sound]
'VOA.AffineSl2BelowCritical.compactRealStructure_isCompact' depends on axioms: [propext, Classical.choice, Quot.sound]
'VOA.AffineSl2BelowCritical.IsCompactRealStructure.sugawara_sesq' depends on axioms: [propext, Classical.choice, Quot.sound]
'VOA.AffineSl2BelowCritical.layerNondegenerate_of_lower' depends on axioms: [propext, Classical.choice, Quot.sound]
'VOA.BelowCriticalNoWall.CurrentSystem.layerNondegenerate_of_lower' depends on axioms: [propext, Classical.choice, Quot.sound]
'VOA.MatrixRigidity.matrixLayerSign_rigid' depends on axioms: [propext, Classical.choice, Quot.sound]
'VOA.MatrixRigidity.matrixLayerSign_eventually' depends on axioms: [propext, Classical.choice, Quot.sound]
'VOA.AffineSl2BelowCritical.gramPoly_coeff_top' depends on axioms: [propext, Classical.choice, Quot.sound]
'VOA.AffineSl2BelowCritical.gramPoly_natDegree_le' depends on axioms: [propext, Classical.choice, Quot.sound]
'VOA.AffineSl2BelowCritical.layerPositive_iff_matrixLayerSign' depends on axioms: [propext, Classical.choice, Quot.sound]
'VOA.AffineSl2BelowCritical.matrixStepNondeg_of_layerNondeg' depends on axioms: [propext, Classical.choice, Quot.sound]
'VOA.AffineSl2BelowCritical.ordinaryPiece_eq_pbwSpan' depends on axioms: [propext, Classical.choice, Quot.sound]
\end{plaincode}

\noindent and, added in the referee revision at a still later commit (same command; the Lean files of this
section were unchanged since the previous run), the compactness of the family, cited in the status of
Definition~\ref{l19:def:bcdata} (the fidelity referee obtained the same output independently), and the two Gram-matrix facts of
Lemma~\ref{l19:lem:bcgram}:

\begin{plaincode}
'VOA.AffineSl2BelowCritical.compactRealFamily_isCompact' depends on axioms: [propext, Classical.choice, Quot.sound]
'VOA.AffineSl2BelowCritical.gram_isHermitian' depends on axioms: [propext, Classical.choice, Quot.sound]
'VOA.AffineSl2BelowCritical.gram_eq_eval' depends on axioms: [propext, Classical.choice, Quot.sound]
\end{plaincode}

\noindent\textbf{Build state of Section~\ref{l19:sec:belowcritical}.} The axiom checks above are isolated-module runs
(\texttt{lake env lean -DmaxSynthPendingDepth=3} against the compiled modules). The project's internal build record (entry of
2026-10-15/16, nothing committed at that time) reports an aggregate \texttt{lake build} of
13,443 jobs with the twelve below-critical modules wired into \leanpath{VOA/\allowbreak Examples.lean}; that run is at neither
of the commits of the two runs above, and the same record's entry of 2026-10-27 states that no book
chapter had then been verified against an aggregate \texttt{lake build}. The section is therefore \emph{not} certified by an
aggregate build; that the twelve modules are imported by \leanpath{VOA/\allowbreak Examples.lean} was confirmed by reading the
file.

\noindent No declaration of this chapter depends on \lean{VOA.pbw}, on \lean{sorryAx}, or on any axiom beyond
Lean's three standard ones (one, \lean{a3count_eq_pLe3_shift}, proved by \lean{decide}, depends on none).
\emph{Axiom-clean is not unconditional}: the conditional declarations (%
the abstract transfer \lean{gradedUnitary_}\allowbreak\lean{of_}\allowbreak\lean{ambient},
\lean{tower_posDef}, \lean{gradedUnitaryAt_}\allowbreak\lean{iff_}\allowbreak\lean{minors_}\allowbreak\lean{and_}\allowbreak\lean{positive},
\lean{leeYang_}\allowbreak\lean{stronglyRational_}\allowbreak\lean{of_}\allowbreak\lean{admissiblyRational}, and the intermediate below-critical
assemblies \lean{universal_}\allowbreak\lean{gradedUnitary_}\allowbreak\lean{belowCritical}, \lean{…_compact}, \lean{…_of_anchor}) carry their
hypotheses explicitly, as listed in the checklist. The below-critical headline
\lean{universal_}\allowbreak\lean{gradedUnitary_}\allowbreak\lean{belowCritical_}\allowbreak\lean{unconditional} is both axiom-clean and unconditional.

%======================================================================
\section{Defects and traps recorded in this chapter}\label{l19:sec:defects}
%======================================================================

For the maintainers, the findings of this chapter, with severity on a four-step scale:
serious (a statement that is wrong or vacuous), moderate (weaker or stronger than advertised, or a redundant
hypothesis), open gap (a conditional result whose hypothesis is not discharged, or a missing piece), and documentation
only (docstrings and notes).
\begin{enumerate}
\item \textbf{Moderate (design weakness).} \lean{GradedUnitary B R F} quantifies over an arbitrary \lean{GradedFlag},
  which need not be compatible with the vertex operators, $\rho$-stable, or ABLR's $R$-filtration; ``for some
  flag'' is nearly vacuous (Warning~\ref{l19:warn:freeflag}, paper argument). \lean{VirGradedUnitary} pins the flag
  and is unaffected.
\item \textbf{Open gap.} \lean{dsReductionSl2} takes a \lean{BRSTComplex} as input; none exists for \lean{dsQ}
  (\lean{Q_wt} fails for the naive grading, \lean{Q_deriv} unavailable). \lean{RoutesAgreeSl2} is an unproved
  \lean{Prop}.
\item \textbf{Moderate (minor).} \lean{FlagMinors.}\allowbreak\lean{schurOfLeading_}\allowbreak\lean{posDef_}\allowbreak\lean{of_}\allowbreak\lean{sign_}\allowbreak\lean{ratios} and \lean{…negDef…} assume
  $0<D_l$ where $D_l\ne0$ suffices; the general versions exist in \leanpath{PolarizedReesStrongGenerator.lean}.
\item \textbf{Documentation only.} Stale docstrings: \leanpath{FlagMinors.lean}, \leanpath{GradedUnitarity.lean} and
  \leanpath{VirasoroGradedUnitarity.lean} say the pivot/positivity half of the minors criterion is absent
  (it is in \leanpath{MinorPivots.lean}); \leanpath{VirasoroGradedUnitarity.lean} says no real structure is constructed
  on $L(c,0)$ and no positivity is proved (Warning~\ref{l19:warn:stale}); \leanpath{MoonshineModule.lean}'s title and
  TODO comment are stale.
\item \textbf{Documentation only.} The docstring of \lean{RigidBridge.shapForm} (\leanpath{VOA/\allowbreak Core/\allowbreak RigidityBridge.lean})
  says its Gram matrix ``is the level-$k$ Kac matrix''; it agrees with the invariant form only on quasi-primary
  vectors (\leanpath{LeeYangKac.lean} documents the discrepancy at level four; Warning~\ref{l19:warn:shap}).
\item \textbf{Book (not Lean); corrected.} An earlier draft of Chapter~\ref{ch:15} expanded ``ABLR'' wrongly; the sources mean
  Arabi Ardehali--Beem--Lemos--Rastelli, and Chapter~\ref{ch:15} now says so.
\item \textbf{Documentation only.} \leanpath{AffineSl2BelowCriticalGradedUnitarity.lean}: the module docstring calls the
  three inputs of \lean{universal_}\allowbreak\lean{gradedUnitary_}\allowbreak\lean{belowCritical} ``none of which is established in this
  repository'' and the docstrings of \lean{LayerSignAnchor}, \lean{LayerSignRigidity} say ``Named open input''; all three
  are proved for \lean{compactRealFamily} (Warning~\ref{l19:warn:bcstale}).
\item \textbf{Documentation only.} \leanpath{VOA/\allowbreak Core/\allowbreak GradedUnitarity.lean}: since the
  below-critical headline (Theorem~\ref{l19:thm:bcheadline}), the module Status line ``No VOA in this repository is proved
  fully graded unitary'' and the \lean{GradedUnitary} docstring sentence ``No theorem in this repository asserts it of
  anything'' are false: \lean{universal_}\allowbreak\lean{gradedUnitary_}\allowbreak\lean{belowCritical_}\allowbreak\lean{unconditional}
  asserts it of $V^k(\lnineteensl)$ for every real $k<-2$. The module-level warning that the predicate is ``nearly
  vacuous'' (because the flag is free, Warning~\ref{l19:warn:freeflag}) is partly answered by that theorem, whose flag is
  fixed. (Found by the fidelity referee of Section~\ref{l19:sec:belowcritical}.)
\item \textbf{Restriction (not a defect).} Virasoro real structures are constructed only for $c = 1-12\lambda^2$ with
  $\lambda$ real ($c\le1$), although they exist for all real $c$.
\end{enumerate}

%======================================================================
\section{Referee's checklist}\label{l19:sec:checklist}
%======================================================================

\begingroup\small
\begin{longtable}{>{\raggedright\arraybackslash}p{0.29\textwidth}>{\raggedright\arraybackslash}p{0.31\textwidth}>{\raggedright\arraybackslash}p{0.22\textwidth}>{\raggedright\arraybackslash}p{0.1\textwidth}}
\toprule
Lean declaration & Intended meaning & Hypotheses / status & Check\\
\midrule\endhead
\lean{eq_bot_of_bilinear_pos} & bilinear ``positivity'' forces $S=0$ (the old defect) & none; unconditional & Thm.~\ref{l19:thm:bilpos}\\
\lean{SesqForm}, \lean{IsPosOn}, \lean{signedForm} & sesquilinear form; positive on $S$; $(-1)^p$ twist & definitions; witnesses both ways & Def.~\ref{l19:def:sesq}, \ref{l19:def:isposon}\\
\lean{RealStructure} & anti-linear VOA automorphism, $\rho^2 = s$ & definition; $s$ free & Def.~\ref{l19:def:real}\\
\lean{gradedSesq} & $h(a,b) = B(\rho a,b)$ on $V_n$ & definition & Def.~\ref{l19:def:gradedsesq}\\
\lean{layer}, \lean{LayerNondegenerate}, \lean{LayerPositive}, \lean{GradedUnitaryAt}, \lean{GradedUnitary} & ABLR clauses (iv), (vi); graded unitarity & definitions; \textbf{flag free} (Warn.~\ref{l19:warn:freeflag}); asserted only in Thm.~\ref{l19:thm:bcheadline} & Def.~\ref{l19:def:gu}\\
\lean{orthogonal_}\allowbreak\lean{gradedSesq_}\allowbreak\lean{eq} & bilinear and sesquilinear complements agree & $\rho$-stable, $s(n)\ne0$ & Thm.~\ref{l19:thm:orthsame}\\
\lean{layerDecomposition}, \lean{finrank_}\allowbreak\lean{pieceLT_}\allowbreak\lean{add_}\allowbreak\lean{layer} & clause (v) follows from (iv) & reflexive, clause (iv) & Thm.~\ref{l19:thm:layerdec}\\
\lean{isPosOn_}\allowbreak\lean{layer_}\allowbreak\lean{zero_}\allowbreak\lean{vacuum} & vacuum slice of (vi) & $(\lnineteenvac|\lnineteenvac) = 1$ & Thm.~\ref{l19:thm:vacuum}\\
\lean{nondegenerate_}\allowbreak\lean{flag_}\allowbreak\lean{iff_}\allowbreak\lean{minors} & (iv) $\iff$ nonzero leading minors & flag refines basis & Thm.~\ref{l19:thm:minors}\\
\lean{posDef_}\allowbreak\lean{of_}\allowbreak\lean{det_}\allowbreak\lean{leadingPrincipal_}\allowbreak\lean{pos} & converse Sylvester & symmetric real & Thm.~\ref{l19:thm:sylvester}\\
\lean{schurOfLeading_}\allowbreak\lean{posDef_}\allowbreak\lean{of_}\allowbreak\lean{sign_}\allowbreak\lean{ratios} (+\lean{negDef}) & layer definiteness from minor ratios & $D_l\ne0$ (weakened from $D_l>0$) & Thm.~\ref{l19:thm:sylvester}\\
\lean{SesqForm.}\allowbreak\lean{isPosOn_}\allowbreak\lean{signedForm_}\allowbreak\lean{of_}\allowbreak\lean{eq_}\allowbreak\lean{pullback} & positivity pulls back & injective, exact isometry & Thm.~\ref{l19:thm:pullback}\\
\lean{VirGradedUnitary} & Kaufmann/ABLR conjecture for $L(c,0)$ & definition; \textbf{open}, never asserted & Def.~\ref{l19:def:virgu}\\
\lean{gradedUnitaryAt_}\allowbreak\lean{iff_}\allowbreak\lean{minors_}\allowbreak\lean{and_}\allowbreak\lean{positive} & Virasoro reduction to minors & adapted basis \lean{hrb} (never supplied) & Thm.~\ref{l19:thm:virreduction}\\
\lean{virasoroSimpleRealStructure}, \lean{boundaryVirasoroRealStructure} & real structure on $L(c,0)$ & $c = 1-12\lambda^2$, $\lambda\in\lnineteenR$; boundary: none & Thm.~\ref{l19:thm:virreal}, \ref{l19:thm:boundaryreal}\\
\lean{rhoStableOn_}\allowbreak\lean{vir_}\allowbreak\lean{pieceIn} & length flag is $\rho$-stable & any $\rho$; unconditional & Thm.~\ref{l19:thm:rhostable}\\
\lean{lengthTwoQ_}\allowbreak\lean{quasiPrimary}, \lean{canB_}\allowbreak\lean{lengthTwoQ_}\allowbreak\lean{self} & $L_1Q_{2m} = 0$; $(Q|Q) = \frac12\binom{4m-2}{2m-4}c(c+\rho_m)$ & $m\ge2$; unconditional & Thm.~\ref{l19:thm:norm}\\
\lean{Bsimp_}\allowbreak\lean{boundaryLengthLETwoTowerAt_}\allowbreak\lean{self_}\allowbreak\lean{sign} & stress towers negative, length-two towers positive & $n\ge2$ & Thm.~\ref{l19:thm:towersign}\\
\lean{boundary_}\allowbreak\lean{gradedUnitaryThroughTwo} & (iv)+(vi) for $p\le2$, all weights, $M(2,2n+1)$ & $n\ge2$ only; \textbf{unconditional} & Thm.~\ref{l19:thm:throughtwo}\\
\lean{exists_}\allowbreak\lean{level4_}\allowbreak\lean{singular_}\allowbreak\lean{iff} & level-4 singular vector iff $c=-22/5$ & none & Thm.~\ref{l19:thm:lysing}\\
\lean{maxIdealVir_}\allowbreak\lean{leeYang_}\allowbreak\lean{inf_}\allowbreak\lean{grading_}\allowbreak\lean{four} & $\mathcal M\cap V_4 = \C s$ & none (level 4 only) & Thm.~\ref{l19:thm:lylevel4}\\
\lean{leeYang_}\allowbreak\lean{stronglyFinite} & $L(-22/5,0)$ simple, self-dual, $C_2$-cofinite & none; unconditional & Thm.~\ref{l19:thm:lysf}\\
\lean{leeYang_}\allowbreak\lean{stronglyRational_}\allowbreak\lean{of_}\allowbreak\lean{admissiblyRational} & strongly rational & \textbf{assumes} rationality & Thm.~\ref{l19:thm:lysf}\\
\lean{lowestWeight_}\allowbreak\lean{eq_}\allowbreak\lean{zero_}\allowbreak\lean{or_}\allowbreak\lean{neg_}\allowbreak\lean{one_}\allowbreak\lean{fifth} & lowest weights $\in\{0,-1/5\}$ & necessary condition only & Thm.~\ref{l19:thm:lyzhu}\\
\lean{LeeYangRiccati.}\allowbreak\lean{tower_}\allowbreak\lean{posDef} & abstract length induction & matrix hypotheses, not instantiated & Thm.~\ref{l19:thm:riccati}\\
\lean{a3count_}\allowbreak\lean{eq_}\allowbreak\lean{pLe3_}\allowbreak\lean{shift}, \lean{negDef_w12}… & length-3 counts and signs & finite; external data & Thm.~\ref{l19:thm:len3}\\
\lean{weil_}\allowbreak\lean{polarized_}\allowbreak\lean{twist_}\allowbreak\lean{option_}\allowbreak\lean{a} & Weil parity of channels & finite; external data & Thm.~\ref{l19:thm:weil}\\
\lean{MinorPayload.}\allowbreak\lean{sign_}\allowbreak\lean{det_}\allowbreak\lean{leadingPrincipal_}\allowbreak\lean{pos} & minor signs from normal form & hypothesis $\equiv$ conclusion & Thm.~\ref{l19:thm:payload}\\
\lean{dsCentralCharge_}\allowbreak\lean{eq_}\allowbreak\lean{minimalModel}, \lean{…eq_affine_add_ghost} & DS central charge & arithmetic & Thm.~\ref{l19:thm:dsarith}\\
\lean{dsQ_mul_self} & $Q^2 = 0$ & $k\ne-2$; cohomology \textbf{not built} & Thm.~\ref{l19:thm:dsq}\\
\lean{dsEmbed_injective}, \lean{dsLatticeScreening_}\allowbreak\lean{centralCharge} & DS complex, embedding, screening kernel & integral $k+2$ for the screening & Thm.~\ref{l19:thm:dsmor}\\
\lean{compactCurrentSignedForm85_}\allowbreak\lean{positive} & weight-one affine positivity at $-8/5$ & three vectors; no VOA real structure & Thm.~\ref{l19:thm:aff85}\\
\lean{GradedIsometricEmbedding.}\allowbreak\lean{gradedUnitary_}\allowbreak\lean{of_}\allowbreak\lean{ambient} & transfer through isometric embedding & \textbf{conditional}: reduction + ambient GU & Thm.~\ref{l19:thm:transfer}\\
\lean{universalForm}, \lean{compactRealFamily}, \lean{IsCompactRealStructure}, \lean{ordinaryFlag} & canonical form, compact $\rho$, current-length flag on $V^k(\lnineteensl)$ & definitions; $\rho$ proved compact & Def.~\ref{l19:def:bcdata}\\
\lean{universal_}\allowbreak\lean{gradedUnitary_}\allowbreak\lean{belowCritical_}\allowbreak\lean{unconditional} & $V^k(\lnineteensl)$ graded unitary for every real $k<-2$ & $k$ real, $k<-2$ only; \textbf{unconditional} & Thm.~\ref{l19:thm:bcheadline}\\
\lean{LayerSignAnchor}, \lean{LayerSignRigidity} & the two analytic inputs of the assembly & definitions; both proved for the compact family & Def.~\ref{l19:def:bcinputs}\\
\lean{matrixLayerSign_}\allowbreak\lean{eventually}, \lean{layerSignAnchor_}\allowbreak\lean{compact} & layer signs at large negative level & degree and top-coefficient bounds (proved for \lean{gramPoly}); unconditional & Thm.~\ref{l19:thm:bcanchor}\\
\lean{matrixLayerSign_}\allowbreak\lean{rigid}, \lean{layerSignRigidity_}\allowbreak\lean{compact} & signs constant along $(-\infty,-2)$ & continuous Hermitian family, nondegenerate steps; unconditional & Thm.~\ref{l19:thm:bcrigid}\\
\lean{universalForm_}\allowbreak\lean{separating_}\allowbreak\lean{of_}\allowbreak\lean{lt} & $B_k$ nondegenerate on $V^k_E$, $k<-2$ & none; unconditional & Thm.~\ref{l19:thm:bcnondeg}\\
\lean{CurrentSystem.}\allowbreak\lean{layerNondegenerate_}\allowbreak\lean{of_}\allowbreak\lean{lower}, \lean{sugawara_sesq} & prefix no-wall lemma; Hermitian Sugawara identity & lower-weight signs + full nondegeneracy (any VOA) & Thm.~\ref{l19:thm:bcnowall}\\
Geometric files & geometric VOAs, chiral de Rham, elliptic genus & \textbf{no declarations} & §\ref{l19:sec:geometric}\\
\bottomrule
\end{longtable}
\endgroup

%======================================================================
\section{Notes and references}\label{l19:sec:notes}
%======================================================================

\paragraph{Graded unitarity.} The datum formalised here is Definition~2.2 of A.~Arabi Ardehali, C.~Beem,
M.~Lemos and L.~Rastelli, \emph{Graded unitarity in the SCFT/VOA correspondence}, arXiv:2507.23781 (JHEP 2026),
as cited in the source files (the files note the paper was not in the repository when the encoding was written).
The conjecture for the boundary minimal models goes back to R.~M.~Kaufmann, preprint hep-th/9405041 (1994),
\S4.2 and eq.~(4.4), as cited in the source files (we do not reproduce its title); the remark that the full structure was not established is from C.~Beem and S.~Garner, arXiv:2509.10364 (2025). The
SCFT/VOA correspondence is C.~Beem, M.~Lemos, P.~Liendo, W.~Peelaers, L.~Rastelli and B.~C.~van Rees,
\emph{Infinite chiral symmetry in four dimensions}, Comm.\ Math.\ Phys.\ 336 (2015) 1359--1433, arXiv:1312.5344. The
admissible-partition (filtered) basis cited as \lean{FlagBasisHyp} is J.~Salazar, arXiv:2405.07126; the fusion-ring
data of \leanpath{WeilPolarizedChannelParity} is from P.~Godfard, arXiv:2406.07459.

\paragraph{Below-critical graded unitarity.} The proof formalised in Section~\ref{l19:sec:belowcritical} follows the
repository note
\texttt{Notes/\allowbreak graded-unitarity-verification/\allowbreak graded\_unitarity\_campaign/\allowbreak integrable\_deformation\_2026\_10\_01/\allowbreak audit/\allowbreak oct15\_auxiliary\_belowcritical\_prefix\_no\_wall\_primary.md}
(prefix no-wall lemma, large-level asymptotics, continuity in the level). The vacuum determinant is the
Kac--Kazhdan formula: V.~G.~Kac and D.~A.~Kazhdan, \emph{Structure of representations with highest weight of infinite-dimensional
Lie algebras}, Adv.\ Math.\ 34 (1979) 97--108. The Sugawara construction and the critical level: H.~Sugawara, Phys.\ Rev.\ 170
(1968) 1659; B.~Feigin and E.~Frenkel, \emph{Affine Kac--Moody algebras at the critical level and Gelfand--Dikii algebras},
Int.\ J.\ Mod.\ Phys.\ A 7, Suppl.\ 1A (1992) 197--215.

\paragraph{Unitary VOAs.} C.~Dong and X.~Lin, \emph{Unitary vertex operator algebras}, J.\ Algebra 397 (2014)
252--277, arXiv:1305.5526. C.~Ai and X.~Lin, \emph{On the unitary structures of vertex operator superalgebras},
J.\ Algebra 487 (2017) 217--243, arXiv:1606.09445. S.~Carpi, Y.~Kawahigashi, R.~Longo and M.~Weiner, \emph{From vertex
operator algebras to conformal nets and back}, Mem.\ Amer.\ Math.\ Soc.\ 254 (2018), no.~1213, arXiv:1503.01260.
D.~Friedan, Z.~Qiu and S.~Shenker, \emph{Conformal invariance, unitarity, and critical exponents in two dimensions},
Phys.\ Rev.\ Lett.\ 52 (1984) 1575; P.~Goddard, A.~Kent and D.~Olive, \emph{Unitary representations of the Virasoro and
super-Virasoro algebras}, Comm.\ Math.\ Phys.\ 103 (1986) 105--119.

\paragraph{Lee--Yang.} J.~L.~Cardy, \emph{Conformal invariance and the Yang--Lee edge singularity in two dimensions},
Phys.\ Rev.\ Lett.\ 54 (1985) 1354. C.~N.~Yang and T.~D.~Lee, Phys.\ Rev.\ 87 (1952) 404 and 410. Rationality of the
minimal models: W.~Wang, \emph{Rationality of Virasoro vertex operator algebras}, Int.\ Math.\ Res.\ Notices 1993, 197--211
(not formalised).

\paragraph{Drinfeld--Sokolov reduction.} B.~Feigin and E.~Frenkel, \emph{Quantization of the Drinfeld--Sokolov
reduction}, Phys.\ Lett.\ B 246 (1990) 75--81; M.~Bershadsky and H.~Ooguri, \emph{Hidden $SL(n)$ symmetry in conformal
field theories}, Comm.\ Math.\ Phys.\ 126 (1989) 49--83; T.~Arakawa, \emph{Rationality of W-algebras: principal nilpotent
cases}, Ann.\ of Math.\ 182 (2015) 565--604, arXiv:1211.7124.

\paragraph{Inverse reduction.} The inverse Drinfeld--Sokolov programme (Adamovi\'c's free-field embedding of the
boundary admissible $L_k(\lnineteensl)$ into $L(c,0)\otimes\lnineteenPi(0)$) is under development in the library and is
not part of this release. The literature behind it: A.~M.~Semikhatov, \emph{Inverting the Hamiltonian reduction in
string theory}, hep-th/9410109 (1994); D.~Adamovi\'c, \emph{Realizations of simple affine vertex algebras and their
modules: the cases $\widehat{sl(2)}$ and $\widehat{osp(1,2)}$}, Comm.\ Math.\ Phys.\ 366 (2019) 1025--1067,
arXiv:1711.11342; C.~Beem, C.~Meneghelli and L.~Rastelli, \emph{Free field realizations from the Higgs branch}, JHEP 09
(2019) 058, arXiv:1903.07624.

\paragraph{Geometric VOAs and chiral de Rham.} Y.-Z.~Huang, \emph{Two-Dimensional Conformal Geometry and Vertex Operator
Algebras}, Progr.\ Math.\ 148, Birkh\"auser 1997. F.~Malikov, V.~Schechtman and A.~Vaintrob, \emph{Chiral de Rham
complex}, Comm.\ Math.\ Phys.\ 204 (1999) 439--473, arXiv:math/9803041. L.~A.~Borisov and A.~Libgober, \emph{Elliptic
genera of toric varieties and applications to mirror symmetry}, Invent.\ Math.\ 140 (2000) 453--485,
arXiv:math/9904126. A.~Kapustin, \emph{Chiral de Rham complex and the half-twisted sigma-model}, hep-th/0504074.
E.~Witten, \emph{Elliptic genera and quantum field theory}, Comm.\ Math.\ Phys.\ 109 (1987) 525--536.

\endgroup
